\documentclass[12pt,a4paper]{book}
\usepackage[T1]{fontenc}
\usepackage[utf8]{inputenc}
\usepackage{lmodern}
\usepackage[margin=2.5cm]{geometry}
\usepackage{setspace,microtype,booktabs,emptypage,fancyhdr,etoolbox}
\usepackage[hidelinks]{hyperref}
\setstretch{1.5}
\raggedbottom
\setlength{\emergencystretch}{3em}
\setlength{\headheight}{15pt}
\pagestyle{fancy}
\fancyhf{}
\fancyhead[LE,RO]{\small\thepage}
\fancyhead[LO,RE]{\small\itshape\nouppercase{\leftmark}}
\renewcommand{\chaptermark}[1]{\markboth{#1}{}}
\renewcommand{\headrulewidth}{0pt}
\setcounter{tocdepth}{0}
\makeatletter
\patchcmd{\l@chapter}{1.5em}{2.3em}{}{\PackageError{reader}{Contents number field patch failed}{}}
\makeatother
\title{Sentimental Anti-Realism\\[0.6em]\large Lifting All Four Limbs and Expecting to Float}
\author{Katherine Vane}
\date{}
\hypersetup{pdftitle={Sentimental Anti-Realism: Lifting All Four Limbs and Expecting to Float},pdfauthor={Katherine Vane}}
\begin{document}
\maketitle
\begingroup
\setstretch{1.2}
\tableofcontents
\endgroup
\part{How We Persuade One Another}

\chapter{Who Pays for Being Good?}

The man was generous with my place in the queue.

We were waiting at a supermarket counter in Upper Austria. I had ordered two sides of salmon and had come to collect them. He had arrived before me. A woman arrived after us, but instead of standing behind me she went to another part of the counter, where she was easier for the assistant to see.

There was a long wait. One assistant was occupied with a customer who had several special requests. Eventually another employee came out and asked who was next.

I pointed to the man. He was next, and he was served.

Then the first assistant became free and turned to the woman who had arrived after me.

``I was here first,'' I said, and asked for my salmon.

The woman began giving her order anyway. I repeated that I had been waiting longer. At this point the man whose place I had just defended intervened.

``No, no, let the lady go first. We take things easy here. It doesn't matter who is served first and who is served later.''

He had already been served.

I kept my place. Both of them shook their heads. The disagreement was now apparently about my impatience.

Nothing much happened. Nobody was injured. The purchase was not a matter of life and death. Yet I remembered the exchange, partly because I resented it and partly because the man's argument was so convenient. He could appear courteous without giving up anything. The woman would receive the courtesy. I would supply it.

Had I agreed, all three of us could have gone home with an agreeable account of the incident. She would have met a helpful stranger. He would have restored good manners. I could have congratulated myself on being above such a petty dispute. The extra wait would still have been mine, but mentioning it would have spoiled the story.

This book concerns the things that spoil the story.

\section*{A favour that has already been promised}

There are plenty of reasons to let someone go ahead at a shop counter. They may be ill, distressed, or about to miss their bus. They may have one item while we have a trolley full. Sometimes there is no particular reason. We feel generous and offer.

The man at the counter could have offered his own place before he was served. He could also have asked whether I minded waiting. What irritated me was his assumption that my agreement was unnecessary. He announced a general principle about taking life easily at the precise moment when somebody else would have to practise it.

The woman was also free to ask. Asking would have acknowledged that I had a choice. Beginning her order while I objected treated that choice as a nuisance.

The difference matters because a queue is useful precisely when strangers have no desire to be generous. We need some way of deciding whose turn it is. Order of arrival is easy to understand and usually easy to check. Without it, the assistant must choose among people who move closer, speak louder, look more important, or know her name. A queue gives the quiet person a chance of being served before closing time.

It also saves us from hearing everybody's argument about why their errand matters more than ours. I do not need to explain the purpose of my salmon. The next customer need not disclose what is happening at home. We can transact our business without making our lives available for public judgment.

There is no great moral achievement in insisting on this. I was annoyed and wanted to leave. But wanting to leave was enough. I did not need a sick relative or a departing train to justify keeping the place I had waited for.

The small pressure to invent such an excuse is familiar. ``I'd love to, but\ldots{}'' We say it when we do not love to. We manufacture an obstacle because a plain refusal seems unacceptable. A dentist's appointment will be respected. Our preference may not be.

With repetition, this becomes a peculiar way of dealing with one another. We invite requests and discourage honest answers. Then we complain that people are insincere.

\section*{The pleasure of saying no}

On another occasion, at a supermarket in Vienna, a woman behind me asked to go first because she had only one item. I had two.

I refused.

``Why?'' she asked.

``Because I am a bad person.''

The answer ended the discussion. I did not have to demonstrate how busy I was or submit my reasons for her approval. I had accepted the worst description she might give me, which left her with little to add.

I enjoyed that answer. I still find it funny. It was also more aggressive than the occasion required. She had asked for a favour. I answered as though she had already accused me of a crime.

These two facts can coexist: I was entitled to refuse, and I enjoyed making the refusal unpleasant. My right to keep my place does not settle the quality of my behaviour.

After paying, I saw her near the entrance, talking to another woman. This pleased me too. Apparently she had time to talk after all. I could now remember the incident as an attempted imposition that I had correctly resisted.

But I had not asked why she was in a hurry. I did not know what the conversation was about. I had found a detail that suited my annoyance and treated it as confirmation.

It would be easy to turn either shopping trip into a little sermon about the decline of manners, the arrogance of strangers, or the cowardice of people who surrender whenever somebody frowns at them. The sermons would say more than the incidents support. A rude customer does not supply reliable information about the future of Europe. Nor does my account become impartial because I acknowledge that I was irritated.

Still, there is no need to abandon the simple judgment that someone tried to take my turn. Uncertainty about a person's motives does not oblige us to pretend that we cannot see what they are doing. We can object to the action without inventing a complete account of the person.

That would have been enough at both counters. It is the part I find hardest to remember afterwards, when the argument is over and I am improving it in my head.

\section*{When the cost is larger}

A few minutes in a shop are easily given away. Other requests are harder to refuse and harder to satisfy.

Imagine a family deciding how to care for an elderly parent who can no longer manage alone. Everyone agrees that she should stay in her own home. It is where she wants to be. She knows the rooms, the neighbours, the view from the kitchen. A move would frighten her.

One daughter lives nearby. The others live farther away. At first the daughter visits after work. Then she begins shopping, washing clothes, arranging appointments, and answering calls in the night. Her brothers and sisters are grateful. They tell her she is wonderful with their mother.

Their praise may be entirely sincere. It does not give her back an evening.

At some point she asks for another arrangement. Now the family must discuss money, visits, paid help, or a move. These are unwelcome subjects. Somebody says that their mother has done so much for them. Somebody else says that putting her in a home would break her heart. The daughter is asked whether she can manage a little longer.

The mother's wishes are real. So is the daughter's exhaustion. Repeating the wishes does not answer the exhaustion. It merely makes the daughter responsible for disappointing her mother if she refuses to continue.

Unlike a stranger at a checkout, she does have obligations. Years of love and care are not equivalent to an unsolicited request for a place in line. She might reasonably do things for her mother that she would never agree to do for somebody else. The question is how much she can do, for how long, and what the rest of the family will contribute.

Those questions are sometimes treated as evidence that affection has failed. In fact, leaving them unanswered can destroy affection. A person who was glad to help may gradually become furious with everyone, including the parent she loves. Every telephone call becomes a demand. Visits are hurried. She begins to dread the familiar voice that she will one day miss.

There may be no arrangement that makes the family happy. Paid help may be expensive or difficult to find. A move may be necessary and remain a source of grief. The daughter may feel guilty even if she has done more than anyone could reasonably ask.

Honest discussion cannot remove all this. It can at least prevent the others from confusing the arrangement they prefer with an arrangement that costs nothing. If they want their mother to remain at home, they must discuss who will make that possible. The nearby daughter is a member of the family, not a service included with the house.

\section*{Good intentions do not finish the work}

What happens within a family can also happen in a public meeting, although the obligations and the means of meeting them are different.

Suppose a town proposes an evening centre for young people. The purpose is sensible. There is little for them to do, and residents complain about groups gathering outside shops. A building is available. Volunteers offer to help.

Before the centre opens, someone asks who will supervise it after ten, what will happen when a visitor threatens another visitor, and who will clean it the next morning. These are necessary questions. They can also make the person asking them sound hostile to young people, especially beside someone describing all the good the centre could accomplish.

An enthusiastic proposal is easier to applaud than a rota.

If the difficult questions are postponed, the work will still arrive. A caretaker may stay late without pay. A volunteer may find herself dealing with behaviour she has neither the authority nor the training to manage. The people living next door may lose sleep. These possibilities do not establish that opening the centre is a mistake. They establish that opening the door is only part of the job.

The proposal becomes credible when it explains how that job will be done. Perhaps the town will pay staff, set closing hours, and exclude visitors who repeatedly threaten others. Such decisions may disappoint some of the people the centre was intended to help. Without them, other visitors may stop coming. A centre that cannot protect its quieter users is failing them, however welcoming its noticeboard sounds.

Opponents must be equally specific. ``It won't work'' is no more informative than ``We must do something.'' A resident may invoke safety because he dislikes having young people nearby. An official may cite expense because another project matters more to her. Practical language can conceal preferences as readily as compassionate language can.

What would help is an account that allows other people to check it: the opening hours, the staffing, the budget, the likely problems, and a way to change the arrangement if experience contradicts the plan. Some uncertainty will remain. The centre may still fail. Taking a proposal seriously means doing this work before announcing success or disaster.

This is what I mean by sentimental anti-realism. The phrase describes a refusal to examine how a pleasing intention is supposed to work. Once a project has been called compassionate, necessary, or just, questions about its effects are treated as defects in the questioner. The good description is made to do the work of a good arrangement.

The same refusal can wear a severe expression. Someone promises to restore discipline, defend tradition, or end waste, and expects the firmness of the announcement to spare him questions about what he intends to do. We are allowed to ask those questions too. An angry voice does not make a plan practical.

\section*{After the argument}

The attraction of these disputes is that they let us find somebody else at fault. The man at the counter was presumptuous. A family might exploit its most conscientious member. A public official might enjoy promising services whose difficulties will become someone else's responsibility. Exposing this can be useful, and occasionally necessary.

It can also become a habit that makes ordinary life miserable.

There is always another inconsistency to notice. A favour is requested awkwardly. A person who lectures us about patience loses his temper. Someone praises generosity while remaining careful with her own money. We can collect these incidents until every encounter seems to confirm our low opinion of people. The collection will contain plenty of true stories. Our attention will have determined which stories entered it.

I would like to be able to object when an objection is needed and then return to what I was doing. At the supermarket that would mean keeping my turn, collecting the fish, and leaving. It would not require me to forgive an injury that barely deserves the name. It would require me to stop enjoying the quarrel once it had served its purpose.

That is more difficult than producing the clever answer. The answer gives immediate satisfaction. Letting the incident end gives us nothing to repeat at dinner.

We do not have to decide whether the woman was habitually selfish or merely having an unpleasant day. We do not have to establish that the man was a hypocrite in every part of his life. There was enough information to settle whose turn it was. There was much less information about everything else.

Larger decisions will leave larger doubts. A daughter may never be certain that she chose the right care for her mother. A town may have to close a project that helped some of its residents. Careful thought does not guarantee a choice we can remember without regret. Sometimes we will have to revise what we did. Sometimes revision will no longer be possible.

We can remain attentive to those failures without requiring every one of them to tell us what life amounts to. There is work to do and there are people we still want to see. The people will not always behave well. Neither will we.

I kept my place at the counter. The few minutes I saved were soon spent on something I cannot now remember. The quarrel lasted longer, because I kept it going after everyone else had left.

\chapter{The Joke That Ends the Argument}

Imagine asking a question at a meeting and hearing someone laugh before you have finished it.

You have asked whether the promised savings include the cost of replacing the equipment. The person presenting the plan looks towards the chairman, smiles, and says that some people would still have us working by candlelight. A few people laugh. The chairman asks for the next question.

The equipment has not become cheaper. Nobody has explained the calculation. Nevertheless, something has happened to your question. It is now associated with the sort of person who obstructs progress. Asking it again will take more nerve than asking it the first time.

You could object to the joke. This would give the presenter an opportunity to tell you to lighten up. You could laugh along, showing that you are a good sport, although it is your objection that the laughter is helping everyone forget. You could produce a joke of your own. Then the meeting could enjoy a contest of wit while continuing to know nothing about the replacement cost.

The presenter may have an excellent answer. Perhaps you overlooked a line in the appendix. Perhaps the question really has been answered three times. But the people in the room cannot learn that from the laugh. They have learned that the presenter feels confident enough to dismiss you, and that several others are willing to help.

\section*{What the audience has been offered}

Ridicule is useful because it can suggest a devastating answer without requiring anyone to give one. The critic behaves as though the error is too obvious to explain. An audience that does not understand the subject may be especially susceptible: everyone else apparently sees the joke, and nobody wants to be the last person to understand it.

This is not an unusual defect confined to foolish audiences. We all use other people's reactions to decide where to direct our attention. There are too many claims to examine from the beginning. If someone whose judgement we trust finds a proposal absurd, that reaction is information. We might reasonably spend our next hour elsewhere.

The difficulty is deciding what the reaction tells us. An engineer's astonishment at a proposed machine could reflect years of relevant experience. A television presenter's astonishment may reflect nothing more than familiarity with the opinions of other television presenters. Both can look equally convincing in a short clip.

Even the engineer owes an explanation when a decision depends on it. Expertise may justify our asking her first. It does not make her raised eyebrow a design review.

The meeting also contains people who are deciding how to behave rather than what to believe. A junior employee may suspect that the question is sensible and still join the laughter. The presenter influences promotions. The questioner does not. Her response tells us something about her situation, but very little about the equipment.

Once several people have laughed, their reactions can reinforce one another. The chairman sees an audience that apparently agrees. Each member of the audience sees other people who apparently agree. A gesture that began as deference to one influential person can be remembered as the considered opinion of the room.

No conspiracy is required. Nobody needs to meet beforehand or issue instructions. The participants know which kinds of disagreement are likely to make the afternoon unpleasant. They act accordingly.

\section*{When a joke contains an argument}

Suppose a minister announces a campaign against unnecessary bureaucracy and creates four committees to administer it. A cartoon showing the committees arguing over the definition of unnecessary may identify a real problem. It directs attention to the minister's conduct and makes the inconsistency easier to understand.

The cartoon is still a simplification. Some administration may be needed to remove worse administration. Perhaps the committees have limited tasks and short deadlines. Their existence alone does not prove the campaign useless. But there is an intelligible criticism to examine. We can ask whether the proposed machinery is disproportionate to the job.

Now imagine a different cartoon that gives the minister enormous ears. It might be funnier. It tells us nothing about the committees.

We do not need to prohibit the second cartoon in order to distinguish it from the first. Nor must every joke supply a public service. People are allowed to be frivolous, cruel, or simply amused. The mistake occurs when we begin treating their success at amusing us as evidence that they have settled the matter under discussion.

An ugly politician can be right. A charming scientist can be wrong. A person who cannot survive five minutes on a comedy programme may understand the drainage system perfectly well. These observations sound embarrassingly elementary until we watch ourselves respond to an actual performance.

The useful question is what remains when the joke is explained without the joke. If the answer is that the proposed saving has omitted a major expense, there is an objection to investigate. If the answer is that the speaker looks ridiculous, we have learned about the performance. We still need to investigate the saving.

Humour can be especially valuable when an institution protects its claims with solemnity. A pompous official can make a mediocre proposal seem beyond ordinary questioning. A joke reminds the audience that the official is a person, perhaps a rather silly one, who should have to explain himself. In that situation laughter can open a discussion. It earns its place by what becomes possible afterwards.

There is no dependable political rule that tells us which laughter serves this purpose. A dissident can expose a ruler's pretensions and then become intolerable to anyone who questions his own. A joke aimed at the powerful can also humiliate an unimportant person who happens to resemble them. We must look at the particular joke, its target, and the criticism it is supposed to make.

\section*{The satisfaction of being above somebody}

Thomas Hobbes called one source of laughter ``sudden glory'': the pleasure of finding something in ourselves to applaud, sometimes through comparison with a defect in somebody else.\footnote{Thomas Hobbes, \emph{Leviathan}, part I, chapter VI, discussion of laughter and ``sudden glory''. \href{https://www.gutenberg.org/files/3207/3207-h/3207-h.htm}{Public-domain text}.} This is an incomplete account of laughter. It is a recognisable account of certain afternoons spent enjoying other people's stupidity.

Someone mispronounces a word. A confident man makes an elementary mistake. A person who has been lecturing us about morality is caught behaving badly. For a moment we are relieved of the need to accomplish anything ourselves. The comparison has improved without effort on our side.

The pleasure does not prove that the criticism is false. Hypocrisy may have been exposed. A mistake may matter. But pleasure makes us less interested in finding out how much the discovery establishes. We prefer the larger conclusion: this person is a fraud, these people are fools, their entire position can be dismissed.

An audience can become attached to a speaker who repeatedly provides that relief. We know the routine. An opponent is quoted, imitated, and reduced to a familiar type. We laugh before the imitation is finished. The pleasure includes recognition: we understand one another, we know what sort of people we are discussing, and we are glad not to be among them.

It is possible to spend years in such company without getting better at answering the opponents' strongest arguments. Their weakest arguments provide more reliable entertainment. Eventually we may know an extensive collection of embarrassing examples and very little about what a capable opponent would actually say.

There is a cost to this even when our broad political preference is sensible. A movement that cannot describe a serious objection accurately cannot prepare for it. A government surrounded by people who laugh at the right moments may continue making an avoidable mistake until the cost is too large to turn into a joke.

The person who refuses to laugh becomes inconvenient before becoming dangerous. He need not denounce the room. He may merely ask what the original speaker meant. This slows things down and spoils a pleasant sense of agreement. It is easy to describe such a person as humourless. Sometimes he is. That does not answer him either.

\section*{A familiar face before a hearing}

In Plato's \emph{Apology}, Socrates complains that he must answer an image of himself that has circulated for years. He refers to the comic Socrates in Aristophanes' \emph{Clouds}, a figure associated with strange speculation and verbal trickery. His difficulty is that many listeners encountered the caricature before hearing his defence. They do not approach him as a stranger whose conduct remains to be examined.\footnote{Plato, \emph{Apology}, especially 18b--19d, translated by Benjamin Jowett. \href{https://www.gutenberg.org/files/1656/1656-h/1656-h.htm}{Public-domain text}.}

This is Plato's presentation of the defence, not a recording of the trial. It does not establish that a comedy caused the execution of Socrates. The more limited problem it describes is familiar enough: once people recognise a caricature, the real person has to speak through it.

A hesitant answer can confirm that he is evasive. A detailed answer can confirm that he is a pedant. An angry answer can confirm that he is unstable. Even an apology can be received as the confession everyone expected. There may be no available manner of speaking that the audience has not already learned to interpret against him.

This does not happen only to famous people. A child becomes the clumsy one in a family. A colleague becomes the person who always complains. A woman becomes the office eccentric. Their later actions may be interpreted through an old story that nobody bothers to reconsider. They may eventually play the part because attempts to escape it have become exhausting.

Some reputations are deserved. A persistent liar should expect difficulty being believed. But even a deserved reputation is a conclusion drawn from earlier conduct. It is not permission to stop noticing new conduct. The question is whether the story can still be corrected.

For the person caught inside it, an endless campaign to establish a better reputation can become a second punishment. Every conversation becomes a hearing. Every unfamiliar person must be persuaded that the familiar description is wrong. There are situations in which explanation helps, and situations in which leaving is the only useful action still available.

We need not turn leaving into a heroic gesture. It can involve defeat, lost friends, and diminished opportunities. There is nothing admirable about making someone choose between humiliation and withdrawal. But the fact that a room refuses to listen does not create an obligation to remain in it indefinitely.

\section*{The shorter the clip, the easier the victory}

The original question at the meeting might take two minutes to state properly. The laugh takes one second. If someone later circulates the exchange, the laugh and a few words of the question can be enough to tell a complete social story: here is a fool being recognised as a fool.

The missing context may include the table the question referred to, a qualification in the previous sentence, or an answer the presenter gave afterwards. None of these omissions needs to be malicious. A person choosing the funniest moment may sincerely think the remaining material dull. The result can still mislead.

A written exchange can suffer the same treatment. A quotation is detached from the paragraph that limited it. An awkward phrase is repeated while the argument disappears. A thousand people can react to the phrase without one of them having read the original. The number of reactions then seems to confirm that the argument was indefensible.

The correction is at a disadvantage. It has to restore distinctions that the joke removed. It sounds laboured. It asks people to spend time discovering that an enjoyable experience was less informative than it appeared. Many would rather enjoy another one.

Those who complain about this sometimes imply that being mocked is a sign of intellectual independence. It is no such thing. An absurd claim can attract justified ridicule. A dishonest speaker may deliberately provoke it and then present the reaction as proof that important people fear the truth. Humiliation can be manufactured into a credential.

The correction remains the same in either case: recover the claim and examine it. This may disappoint both the crowd and its target. The crowd may have dismissed something it did not understand. The target may have mistaken attention for vindication. Neither disappointment needs to become our problem unless it obstructs the inquiry.

\section*{Answering without becoming the entertainment}

At the meeting, the most useful reply is probably a short one: does the calculation include the replacement cost?

This gives the presenter a way to return to the subject without a prolonged dispute about manners. If the answer is yes, the next step is to find the relevant figure. If the answer is no, the omission can be discussed. If another joke follows, the audience has a clearer view of what is being avoided.

There is no guarantee this will work. A chairman may still move on. A hostile audience may enjoy the second joke more than the first. Advice about staying calm often ignores the real inequality between someone who controls the meeting and someone allowed thirty seconds at the microphone. Composure is useful; it is not an alternative source of power.

When the decision matters, a written question may outlast the performance. Asking that an objection be recorded can make it harder to pretend later that nobody raised it. Finding another person who understands the issue may be more useful than winning over the whole room. These are modest possibilities, dependent on the institution and on what retaliation is likely to follow.

Sometimes there is nothing useful to gain by answering. A stranger who has posted a laughing face may have offered no claim, no question, and no sign of willingness to listen. We do not owe him an afternoon of instruction merely because he has made himself irritating.

This is easy to say and difficult to practise. Ridicule leaves an unfinished feeling. We want a final exchange in which everyone recognises that the laugh was cheap and the answer was ours. We compose that exchange on the way home. We improve it while doing something else. The person who laughed may have forgotten the incident long before we deliver the perfect reply in his absence.

The question about the equipment still deserves an answer. The question of whether everyone in the room respects us may remain unanswered. Confusing the two gives the laugher more control over our life than the meeting ever required.

\chapter{What Can Still Be Said}

A project does not become a failure on the day the director admits it.

Long before then, someone has noticed that a necessary component does not work. A supplier has missed a delivery. The staff know that the promised opening date cannot be met. Each fact has travelled through a series of conversations in which people decide how much of it to pass on.

Consider a fictional organisation preparing to replace an old computer system. The new one is supposed to cost less, process work faster, and be ready in September. The director has announced the date. Several senior managers have described the project as a major achievement. The staff who will use it have begun training.

In June, the person supervising the transfer of records discovers that a substantial portion of them cannot be imported reliably. The fault may be repairable, but nobody yet knows how long that will take. She reports the problem to her manager.

He asks whether it is really necessary to put such alarming language in the monthly report. Could she describe it as an issue being resolved? This may seem fair. The technical staff are working on it. A dramatic warning might distract them and provoke interference from people who understand less than they do.

The manager submits a report stating that data transfer remains under active review. His superior summarises this as progress on migration. The director receives a document in which every major heading indicates continued work. Nothing on the page says that September is now improbable.

Each person has made the report a little easier for the next person to receive. The result is difficult for anyone to use.

\section*{The price of accurate language}

It is tempting to describe the staff as cowards. Some may be. But the woman who found the defect faces a problem that her eventual critics may not have faced themselves. She has to decide whether to risk her position over an uncertain technical judgement.

If she warns emphatically and the fault is repaired quickly, she may acquire a reputation for panic. If she warns quietly and the project fails, someone may later ask why she did not make the seriousness clear. If she goes around her manager, he may become hostile. If she does nothing, she will have to work with a system she knew was unreliable.

There may also be a history she understands. Perhaps the last employee to raise a serious objection was excluded from meetings, then criticised for not contributing. No one was formally punished for speaking. The lessons were nevertheless easy to learn.

The director can announce that his door is always open. Unless people who enter it with bad news fare reasonably well afterwards, the announcement is decoration.

An organisation needs more than brave employees. It needs a routine way for unwelcome information to affect decisions. A defect should not have to become a contest between the director's authority and the discoverer's courage. There should be an agreed test, a record of the result, and someone with authority to delay the launch when the test fails.

This does not make judgement unnecessary. Tests can be badly designed. An acceptable failure rate may be disputed. But it gives the disagreement something definite to concern itself with. The question can become whether the records pass the test, rather than whether a particular employee believes in the project.

When an institution treats these as the same question, it invites people to demonstrate loyalty by concealing what it needs to know.

\section*{Agreement that nobody measured}

The director believes the project has broad support. He has heard encouraging remarks at every presentation. No department has formally asked for it to be abandoned. Staff representatives have attended the planning meetings. The internal newsletter has published photographs of the training sessions.

These observations are compatible with genuine enthusiasm. They are also compatible with many people privately expecting trouble and assuming that the decision cannot be changed.

Attendance at a training session tells us that an employee expects to need the new system. It does not tell us that she thinks it is ready. A department may submit an implementation plan because it has been instructed to do so. The plan does not retrospectively turn the instruction into a consultation.

The political economist Timur Kuran calls attention to the difference between what people prefer privately and what they express publicly when expression has costs. His term is \emph{preference falsification}.\footnote{Timur Kuran, \emph{Private Truths, Public Lies: The Social Consequences of Preference Falsification}, Harvard University Press, 1995. \href{https://www.jstor.org/stable/j.ctvt1sgqt}{Book record}.} The idea does not require us to believe that everyone is secretly opposed to every official policy. It tells us why public agreement alone may be poor evidence of private agreement.

The distinction becomes especially important when people use one another's public behaviour to decide whether speaking is worthwhile. An employee sees that no one else has objected and keeps her doubts to herself. The next employee sees her silence and draws the same conclusion. Each makes a locally understandable decision. Together they create the appearance that nobody doubts the plan.

One open objection can change that situation. It may reveal that disagreement is less lonely than people thought. But there is no guarantee of a hidden majority waiting to emerge. The objector may discover that the others sincerely support the project. A sensible account of conformity must leave room for actual conviction.

It must also leave room for limited knowledge. Some people have no settled opinion. They are busy with another task and trust the colleagues responsible. Counting their lack of protest as active support would be misleading. Treating it as evidence that they have been terrorised would be misleading too.

An honest consultation distinguishes these states as far as it can. It allows people to say that they support the aim but doubt the timetable, that they lack enough information, or that they prefer an alternative. A questionnaire offering only enthusiasm and hostility produces a particularly useless kind of certainty.

\section*{What a simple experiment can show}

Solomon Asch's experiments on conformity placed participants in a situation where the immediate question was much easier than the reliability of a computer system. They compared the lengths of lines while other members of the group, working with the experimenter, sometimes gave an obviously wrong answer. Many participants went along with the incorrect majority on some trials; many also resisted. Asch explored how the group's unanimity and the presence of support affected the responses.\footnote{Solomon E. Asch, ``Opinions and Social Pressure'', \emph{Scientific American} 193, no. 5 (1955), pp. 31--35. \href{https://web.mit.edu/curhan/www/docs/Articles/15341_Readings/Influence_Compliance/Asch_1955_Opinions_and_social_pressure.pdf}{Original article, hosted by MIT}.}

The experiment is often retold as proof that human beings will deny anything their group tells them to deny. That is a larger claim than the experiment establishes. Its subjects were in a particular setting, answering a narrow question. We should not use a laboratory result to explain an entire political movement without examining what else is going on.

What makes the result relevant is that even a simple judgement can become harder to express when everyone around us appears to reject it. In ordinary life the evidence is usually less clear than the difference between lines. Other people may know something we do not. Deferring can be sensible. Our difficulty is deciding when deference has become surrender to a pressure unrelated to the evidence.

The confident explanation that everyone else has been hypnotised is seldom helpful. It assumes the conclusion we need to investigate: that our view is correct and their agreement needs a special diagnosis. Different people can reach the same view for different reasons. Some may be knowledgeable, some deferential, some frightened, some opportunistic. A crowd does not have one mind that a critic can diagnose from outside.

An organisation that wants better information should therefore reduce the unnecessary costs of reporting it. Independent written judgements before a group discussion can reveal differences that disappear in a show of hands. A person checking a claim should not always report to the person whose promotion depends on it. A warning should be assessed for accuracy even when the person delivering it is irritating.

These arrangements cannot manufacture honesty. They can make honesty less dependent on an unusual willingness to be punished.

\section*{The dissenter who cannot be wrong}

The employee with the warning deserves a hearing. She does not deserve a guarantee that the organisation will agree with her.

Suppose a competent independent team repeats her checks and finds that the suspected defect came from a temporary test file. The live records transfer correctly. She might accept the explanation. She might identify a flaw in the new test. Either response could contribute to the work.

A different response would be to announce that the independent team had obviously been bought off because it failed to confirm her warning. When asked for evidence, she might point to the organisation's wish to launch on time. That wish is a reason to examine the test carefully. It does not establish that the result was fabricated.

This distinction is easily lost by people who have been treated badly. Once the management has concealed one problem, every reassurance becomes suspect. Suspicion may be justified, but it still needs a way of encountering an answer. Otherwise the first deception gives the critic permission to assert whatever follows.

The same habit can develop among political outsiders. Rejection proves that the establishment fears them. Silence proves that the establishment is suppressing them. An invitation to debate proves that the establishment is panicking. Every possible response strengthens the same belief.

That is a comfortable arrangement for the dissenter. He no longer risks finding out that his position is wrong. His audience may be as demanding of conformity as the institution he attacks. Members learn which doubts can be expressed among the rebels and which will get them called agents of the establishment.

We should defend the possibility of dissent without attaching intellectual merit to the experience of exclusion. A person can be censored unjustly and still be mistaken. A rejected proposal can be excellent, mediocre, or worthless. The conduct of the gatekeeper and the quality of the proposal are separate matters, even when anger makes it difficult to examine them separately.

This is one reason that merely replacing an official authority with an unofficial favourite is not much of an advance. The new authority may be less constrained, keep fewer records, and answer criticism with a more devoted audience. Independence requires more work than changing whom we trust.

\section*{Disagreement, exclusion, and punishment}

People use the word censorship for experiences that differ greatly in their consequences. A hostile response, the loss of an invitation, dismissal from a job, and imprisonment are not interchangeable events. Treating them as interchangeable obscures what needs to be opposed in each case.

An editor has to reject material. A journal has to distinguish competent work from bad work. A meeting has limited time. Nobody acquires an unlimited claim on other people's attention by saying something controversial. A person who deliberately insults an audience should not be astonished if its members leave.

Yet the observation that some exclusion is necessary does not justify every exclusion. A journal can apply its standards unevenly. An employer can describe a relevant factual objection as misconduct. A professional network can ensure that an unpopular view costs a person opportunities far beyond the occasion on which it was expressed.

The useful questions are particular. What was said? Was it relevant to the speaker's responsibilities? What standard was allegedly violated? Is that standard applied to allies as well as opponents? Can anyone appeal to a person who did not make the original decision?

An institution that refuses to answer these questions should not expect a declaration of its good intentions to settle the dispute. Its power is part of what must be examined.

Mill's defence of open discussion in \emph{On Liberty} rests in part on our capacity for error. The suppressed opinion may be true, may contain some truth missing from the prevailing view, or may force the accepted view to be understood rather than merely repeated.\footnote{John Stuart Mill, \emph{On Liberty} (1859), chapter II. \href{https://www.gutenberg.org/files/34901/34901-h/34901-h.htm}{Public-domain text}.} None of this promises that every argument will be worth hearing indefinitely. It explains why the confidence of today's majority is insufficient grounds for making its view impossible to challenge.

The argument becomes harder to honour when the challenge is ugly or appears likely to help people we fear. At that point we are tempted to make an exception for ourselves: our opponents must tolerate our speech because freedom matters, but we must limit theirs because the consequences matter. Freedom has always had consequences. The question is which restrictions can be justified under rules we would accept in less sympathetic hands.

\section*{When the warning finally arrives}

Return to the computer project. September comes. The launch is postponed. The public statement explains that extra time is being taken to ensure a successful transition. Inside the organisation, the staff spend evenings repairing records and preparing another timetable.

At the review, people remember that concerns were raised early. The June report is produced. The phrase ``under active review'' now seems to have contained the warning everyone needed. A senior manager says that the seriousness should have been escalated. The employee who tried to escalate it is asked whether she has written evidence of the conversation.

A worthwhile inquiry would reconstruct the information available at each stage, including what happened when someone tried to report it. It would ask why the organisation's normal procedures made serious trouble look manageable. An inquiry limited to finding the person who failed to use a sufficiently alarming adjective will leave those procedures intact.

The delay may be survivable. The system may eventually work well. That does not recover the money wasted, the evenings lost, or the trust damaged between colleagues. When a person says she no longer believes management's assurances, pointing to the eventual success may not answer what happened to her along the way.

Nor does a late admission restore the opportunities that an earlier admission would have preserved. In June the old system might have been extended calmly. By September contracts have expired, staff have left, and other work has been postponed. The same technical problem now sits inside a larger practical one.

This is the serious meaning of being too late. It need not mean that everything is doomed. It can mean that a preventable loss is now unavoidable, and that improvement must begin from a worse position.

No organisation can act on every warning. Some will be mistaken, confused, or self-serving. What it can do is preserve a route by which a warning earns attention through evidence. When that route is closed, people are left choosing between silence and a public fight. The institution will then condemn the silence as negligence and the fight as disloyalty. It has arranged to be surprised.

\chapter{The Answer Hidden in the Question}

Suppose a council decides to close a library. The announcement calls it a modernisation of local services.

There may be a serious case for the decision. The building may be expensive to maintain. Few people may use it. A nearby library may offer longer hours and a better collection. A mobile service might reach readers who cannot travel. None of these possibilities can be assessed from the word modernisation.

The word performs a different job. It makes the proposed change sound like an improvement before the improvement has been demonstrated. A person who objects can be placed in the position of defending something old merely because it is old.

An opponent might call the same decision an attack on the community. That description supplies its own conclusion. We are now invited to imagine an aggressor and a victim before discovering how the service will change.

If we begin with the practical questions, the discussion becomes less exciting. Where will people borrow books? How far will they have to travel? Who uses the present building? What money will actually be saved? What else happens there that a lending service alone would not replace?

These questions may eventually support a harsh judgement. Perhaps the council is closing a useful institution to achieve a saving too small to justify the damage. Perhaps its opponents are defending a nearly empty building while ignoring a better service. Either conclusion would mean more after the questions have been answered.

\section*{A name can contain a recommendation}

We cannot speak without selecting words, and words seldom arrive without associations. Protection suggests a danger and someone who needs to be protected. Discipline suggests conduct that needs correcting. Liberation suggests someone being held back. Reform suggests that the existing arrangement is defective.

These associations are not automatically dishonest. A measure may protect people. A rule may deserve reform. Trouble begins when the favourable association substitutes for an account of what will happen.

Consider a proposal to protect tenants. The phrase could describe repairs to unsafe buildings, limits on deposits, restrictions on eviction, a cap on rents, or a subsidy paid from taxation. These measures differ in what they require, whom they benefit, and what problems they may create. Agreement that tenants should be protected does not settle which measure should be adopted.

An opponent who calls all of them interference has compressed just as much. Renting a home already depends on an arrangement of rights, contracts, and enforcement. The question is which arrangement to use. A word that makes one side sound active and the other side sound natural can conceal how much has already been decided.

The same difficulty arises with tradition. Someone asks us to preserve a tradition, but we need to know what practice is meant and what preserving it would require. A language taught to children, a religious procession, a rule excluding women from a position, and a familiar street market are not one object merely because each can be described as traditional.

We should be able to value a practice without pretending that age proves its worth. We should also be able to criticise one without behaving as though people who care about it have no reason beyond stupidity. What they would lose may be difficult to replace, even if we conclude that the practice should change.

Plain description cannot remove disagreement. It gives the disagreement a subject. The library closes on a stated date. The rent can rise under stated conditions. A person becomes eligible for an office from which she was previously excluded. Once we can say what is happening, we can argue about whether it is justified.

\section*{What has happened to the sentence?}

The proposed library closure also gives an official several ways to avoid a definite claim. She may say that the council is committed to improving access, that it recognises the value of reading, and that it will place residents at the centre of the transition.

Any of these statements could be sincere. None tells a resident where to return a book next month.

When pressed, the official may explain that she was describing aspirations rather than making promises. That would be reasonable if the distinction had been clear from the start. It becomes evasive when the language first encourages the audience to expect an outcome and later retreats to the safer claim that the outcome was merely desired.

Harry Frankfurt's essay \emph{On Bullshit} distinguishes lying from speech that is indifferent to whether it accurately describes the matter at hand. A liar attends to the truth in order to conceal it. The bullshitter's immediate concern is the impression he produces.\footnote{Harry G. Frankfurt, \emph{On Bullshit}, Princeton University Press, 2005. \href{https://www.jstor.org/stable/j.ctt7t4wr}{Publisher's book on JSTOR}.} The difference helps explain why some public statements are so difficult to correct. Their practical purpose has already been achieved when the audience sees a confident, concerned person at the microphone.

If a numerical claim is exposed as wrong, another can replace it. If a promise becomes inconvenient, it can be redescribed as a direction of travel. If a question remains unanswered, a statement of concern can fill the remaining time. The words need not be random. They need to fit the occasion and protect the speaker's position.

People who write such statements may be competent and personally honest. Their job may require them to produce reassurance before anyone can supply enough information to justify it. They learn what kinds of sentence can pass through several offices without provoking an objection. Specificity introduces risk. General commitment survives.

A document produced this way can contain no sentence that any individual regards as a lie. It can still leave its readers with a substantially false impression. Responsibility has been dispersed among people who each made the wording a little safer for themselves.

There is a temptation to fight this language by writing equally inflated denunciations. The institution becomes a machine for deception, its employees become priests of an official creed, and every bland sentence becomes evidence of deliberate manipulation. Such prose offers the critic the same convenience it condemns. It creates an impression before explaining a case.

The more useful response is less spectacular. Identify the statement, explain what a reasonable reader would understand it to mean, and compare that meaning with what is known. Where the evidence is incomplete, say what would be needed to decide. A sentence that cannot survive this treatment probably needs replacing.

\section*{Familiarity without confirmation}

Now imagine that the council's account appears in a local newspaper. Another newspaper summarises the first report. A commentator cites both. Residents begin saying that there seems to be general agreement that the change will improve access.

There may still be only one source for that prediction: the council's press release. The later appearances have enlarged the audience without adding an independent observation.

This is easy to overlook because each repetition arrives in a different setting. We remember hearing the claim several times and may forget where it began. A familiar sentence can feel less doubtful simply because it no longer takes effort to recognise it. Research on the illusory truth effect examines this increased credibility through repetition; work by Lisa Fazio and colleagues found that relevant knowledge did not always prevent it.\footnote{Lisa K. Fazio, Nadia M. Brashier, B. Keith Payne, and Elizabeth J. Marsh, ``Knowledge Does Not Protect Against Illusory Truth'', \emph{Journal of Experimental Psychology: General} 144, no. 5 (2015), pp. 993--1002, doi:10.1037/xge0000098. \href{https://scholars.duke.edu/publication/1090926}{Duke University publication record}.}

The finding is a reason to check the origin of a repeated claim. It is not evidence that repetition can make everyone believe anything. People can recognise propaganda, retain contradictory knowledge, or become irritated by a slogan. Nor does a familiar statement become false through familiarity. The question is whether repetition is being mistaken for a second reason to believe it.

A useful report tells us what evidence has been added. Another interview with the same official may clarify her position. It does not independently verify her prediction. A survey of residents can supply different evidence, but only if we know who was asked and what questions they answered. A record of library use may be more informative about some matters than either of them.

Tracing a claim back can be tedious. It is particularly tedious when the claim confirms what we already think. We would rather spend the time dismantling an opposing claim, where discovering a bad source feels like progress. But the pleasing claim may be the one we are least equipped to catch when it is wrong.

\section*{The belief that has become part of a life}

Not every resistance to correction comes from public pressure. Sometimes nobody else is present. We are alone with information that would make our own life harder to understand.

Imagine a retired teacher learning that a method she defended for years was less effective than an alternative she dismissed. She might be willing to examine the new evidence. But the question does not concern a method alone. It concerns pupils she remembers, colleagues she argued with, and a career in which she tried to do something worthwhile.

An unwelcome conclusion can seem to turn past effort into waste. She may respond by looking for weaknesses in the research. There may be real weaknesses. A comparison might be unfair, the outcomes too narrow, or the pupils different from hers. The problem is deciding whether these objections are being used to improve the judgement or to postpone any judgement she could not bear.

People often reveal the difference through changing standards. At first a single study is insufficient. Then several studies come from the wrong institutions. Then independent agreement is dismissed as professional fashion. Eventually only evidence produced by someone who already shares the belief will be considered trustworthy.

The pattern can also run in the opposite direction. A younger critic may accept every unfavourable study because discrediting the old method helps establish his own importance. He may treat decades of other people's work as an obstacle that deserves no careful account. Being less invested in an old belief does not make him free of interests.

The practical task remains to understand the evidence. But people may do that task better when correction is not presented as the destruction of an entire life. The teacher can have cared for her pupils, worked diligently, and used an inferior method. These claims are compatible. A finding about effectiveness need not become a complete judgement of her character.

This does not eliminate responsibility. If she ignored strong evidence while she could still have changed her teaching, that matters. It means we should identify the decisions for which she was responsible rather than using the result to condemn every hour she spent at work.

We are often much more precise about this distinction when the person is someone we love.

\section*{What would count against it?}

A question can help when we suspect that a belief is being protected from examination: what result would make us reconsider?

The answer need not be a single decisive experiment. Some beliefs depend on many observations, and some decisions involve values that no measurement can settle. But a factual claim should have some relation to possible evidence against it. If every outcome can be explained as confirmation, the claim is doing little to help us distinguish among outcomes.

Suppose the council says the closure will improve access. What does access mean here? More opening hours across the district? More loans? More residents within a reasonable travelling distance? Better provision for people who cannot leave home? These measures may move in different directions. Choosing one after the outcome is known would give the council an easy way to declare success whatever happens.

The measure should therefore be stated before the change, along with the people it may fail to represent. Increased total borrowing could coexist with the loss of access for a small group of elderly residents. A lower cost per loan could reflect their disappearance from the service. One accurate number cannot decide whether that loss is acceptable.

A critic also needs to describe what would change her view. If improved opening hours, reliable transport, and continued use by vulnerable residents would still leave her certain that the closure was a disaster, perhaps she is defending the building itself. She may have reasons to do so. A building can provide continuity, meeting space, and a sense of belonging. But those reasons should be stated rather than concealed inside a prediction about book loans.

This distinction allows disagreements to become more honest. We may agree on what is likely to happen and still disagree about what should matter most. That is a different problem from one side refusing to look at a result.

\section*{After a word has been withdrawn}

The library eventually closes. Perhaps the new arrangements work better than many expected. Perhaps they fail some residents badly. More likely, there are gains and losses that different people experience in different proportions.

The official word modernisation will not tell us which occurred. The opponents' phrase attack on the community will not tell us either. We need the less convenient particulars.

There is no guarantee that accurate language will prevail. Some people will keep using the old descriptions because they remain useful to their side. Others will no longer be interested. The argument that once filled a room may now seem too small to revisit.

For those who lost something, this indifference can be painful. A familiar place disappears, and the public explanation gives no name to what it contained. They are told that the service continues. They may know that perfectly well. What they miss is the person who remembered their taste in books, the neighbour encountered on Tuesday, or somewhere to go without being expected to buy anything.

Such losses do not automatically defeat a case for change. They do deserve to appear in the account. Calling a decision necessary does not require pretending that nothing worthwhile was sacrificed. Calling a loss regrettable does not require preserving every arrangement forever.

We should be suspicious of descriptions that allow us to avoid either admission. They spare the speaker discomfort by giving the reader a less recognisable world.

\chapter{Who Gets Called an Expert}

When a water pipe bursts, we usually want a plumber. We do not insist on discovering the principles of plumbing independently before allowing anyone near the valve. We ask about experience, reputation, price, and when someone can arrive. Then we rely on knowledge we do not possess.

Much of adult life works this way. We depend on people who understand things we have neither the time nor the ability to master. The claim that we should think entirely for ourselves becomes absurd if it means that every bridge calculation, laboratory result, and repair must be personally reproduced before we use it.

The harder question is how to judge an expert when the result will not be visible by the end of the afternoon. We can see whether the pipe still leaks. We may be unable to tell whether a long-term forecast is well founded, whether a research paper uses an appropriate comparison, or whether a specialist speaking on television represents the field accurately.

At that point we rely on institutions as well as individuals. A university appointment, a professional qualification, a journal, or an advisory body supplies information about someone's competence. These are useful signals. They become dangerous when we forget what, exactly, they signal.

A person may have completed demanding training and still be speaking outside it. An impressive publication record may concern a narrow problem quite different from the one now being discussed. Membership of a committee tells us that someone was selected. We still need to know who selected her and for what purpose.

\section*{The person at the microphone}

Suppose a broadcaster needs a scientist to discuss a report. The producer wants someone who understands the subject, can arrive in time, speaks clearly, and will answer short questions without a long preliminary lecture. A researcher who meets these requirements may be excellent. But the requirements are not identical to those for doing excellent research.

An unusually careful specialist may refuse to simplify the report into a confident prediction. Another may understand the field less deeply but be very good at giving the audience a sentence it can remember. If the second person appears regularly, familiarity can begin to substitute for evidence of authority.

The producer does not need to intend this. A programme has to be made. The reliable guest is invited again because previous appearances went smoothly. Other programmes notice the guest and invite her too. Eventually an audience encounters the same person in so many settings that it assumes she must be the natural representative of the subject.

She may also begin to see herself that way. A qualified opinion becomes a public position. Viewers expect consistency. An admission of error will travel with her name. Questions from colleagues can sound like attempts to undermine work she now regards as a public duty.

None of this makes her evidence wrong. It does give us reason to distinguish the evidence from the role she has acquired. A person can be a useful guide to a field without becoming its authorised voice on every question.

The broadcaster has another difficulty. Inviting one dissenter beside one representative of a large body of research can make the evidence appear evenly divided when it is not. Inviting only familiar supporters can conceal a significant dispute. The solution cannot be a rule requiring either permanent balance or permanent unanimity. It requires editorial knowledge of the particular disagreement.

We would learn more if the programme explained where the field agrees, what it continues to dispute, and which disagreement would change the decision under discussion. A shouting match between two people introduced as experts does little of that work.

\section*{Who owns the science?}

During a World Economic Forum discussion in 2022, the UN communications official Melissa Fleming said, ``we own the science''. She was discussing efforts to make UN climate information more prominent in Google searches after encountering what she described as distorted information.\footnote{World Economic Forum, ``Tackling Disinformation'', discussion transcript published September 2022, Melissa Fleming's remarks on Google and UN climate resources. \href{https://www.weforum.org/stories/industries-in-depth/tackling-disinformation-agenda-dialogues/}{Full transcript}.}

The context matters. This was a claim about communicating climate information, not an announcement that the UN had purchased the laws of nature. It still deserves criticism. An institution can commission research, assemble findings, and explain why an account deserves confidence. Ownership is a poor description of the authority any of those activities should confer.

Search results must be ordered somehow, and reliable material should be easier to find than demonstrable nonsense. But deciding which institution receives prominence is an exercise of power over public attention. The arrangement should be open to examination: what qualifies a source, how disputed claims are handled, and how a mistaken classification can be corrected.

A communications official may regard such questions as obstacles to getting accurate information before the public. They are also part of establishing why the public should trust the arrangement. If the answer amounts to saying that the institution exists to do good, we have received a description of its intentions where an explanation of its procedures was needed.

\section*{Before any result exists}

Research costs money. Someone must pay salaries, maintain equipment, obtain materials, and allow people time to try things that may fail. The source of that money can shape the research before anyone has a result to distort.

Imagine a company developing a new building material. It funds tests of durability, ease of installation, and performance under conditions likely to favour its intended uses. Those can be legitimate tests. They may produce accurate and useful findings.

A local authority buying the material may need additional answers. How does it perform in older buildings? What does it cost to remove? Can damaged sections be repaired by ordinary contractors? How does it compare with a cheaper material already available? A technically sound report on the first set of questions does not supply answers to the second.

The company need not falsify a measurement to benefit from choosing what gets measured. A field can become rich in evidence about profitable uses and poor in evidence about neglected alternatives. Later, a review of the available literature may faithfully reproduce this imbalance. Its authors can only assess the studies that exist.

Public funding has priorities too. A ministry may favour work aligned with an announced programme. A foundation may choose questions that fit its mission. A university may prefer projects likely to attract further funding. Calling money public or charitable does not remove the choices made through it.

Some choices are unavoidable. No funder can support every question. The problem is whether the priorities are visible and whether other routes remain available. A research system becomes less trustworthy when one approach is both favoured by the dominant funders and treated as the only approach that deserves serious discussion.

This is a stronger criticism than the casual accusation that scientists will say whatever they are paid to say. The accusation makes research sound like a simple market in lies. In practice, honest researchers can spend productive careers answering a selected set of questions while important unselected questions remain unanswered.

\section*{What a conflict of interest establishes}

A company's sponsorship is a reason for scrutiny. It is not a calculation that lets us reverse the reported result.

An effective product does not become ineffective because its manufacturer funded a study. A poorly conducted independent study does not become reliable because its authors dislike the manufacturer. We must inspect the methods and the evidence, including the choices made before the measurements began.

There is nevertheless empirical reason to take sponsorship seriously. A 2017 Cochrane review of research on drug and device studies found that industry-sponsored studies more often reported favourable efficacy results and favourable overall conclusions than studies with other funding. The difference was not adequately explained by the conventional methodological indicators the review examined.\footnote{Andreas Lundh, Joel Lexchin, Barbara Mintzes, Jeppe B. Schroll, and Lisa Bero, ``Industry Sponsorship and Research Outcome'', \emph{Cochrane Database of Systematic Reviews} (2017), MR000033, doi:10.1002/14651858.MR000033.pub3. \href{https://www.cochrane.org/evidence/MR000033_industry-sponsorship-and-research-outcome}{Review and plain-language account}.}

That finding describes an association across bodies of research. It does not identify every sponsored paper as biased or tell us the mechanism in each case. It supports a demand for safeguards that reach beyond a declaration that the authors used respectable methods.

Among the relevant safeguards are a clear account of the planned study, access to its results even when they disappoint the sponsor, and freedom to publish conclusions the sponsor dislikes. An evaluator also needs to know whether the chosen comparison was genuinely demanding. Beating an unsuitable alternative can produce an impressive result with limited practical meaning.

Disclosure is useful because it tells readers where a risk may lie. But disclosure alone does not remove the risk. Printing a conflict in small type cannot compensate for a contract that allows inconvenient results to disappear.

Nor should scrutiny stop at money. An investigator can be attached to a theory, a political cause, a professional rivalry, or a celebrated earlier finding. These interests are harder to list on a form. They still matter when deciding how much independent checking a consequential claim requires.

The appropriate response is therefore neither automatic trust nor automatic accusation. It is to make important conclusions less dependent on the unexamined integrity of one person or one sponsor. Independent teams, accessible methods, preserved records, and a visible history of correction help because people with different interests can inspect one another's work.

\section*{What publication promises}

Peer review usually means that a journal has asked other people with relevant knowledge to assess a submission before publishing it. A reviewer may catch an invalid calculation, an inadequate comparison, a missing qualification, or a conclusion that goes beyond the evidence. This can substantially improve a paper.

It is still not the same as repeating the experiment. A reviewer may have little access to raw data, limited time, and no practical means of examining how the measurements were collected. Even a conscientious review cannot guarantee that a result is true.

The words peer reviewed therefore need a modest interpretation: the work has passed a particular process of assessment. They do not mean that every factual statement is certified, that all relevant alternatives were considered, or that the paper's conclusion should become public policy.

A journal can also reject good work. Reviewers have habits and loyalties. An unfamiliar approach may be genuinely weak or merely difficult for the selected reviewers to assess. A result that challenges an established account may receive demanding criticism. Some of that criticism will be exactly what it needs. Some may protect the established account from a challenge it ought to face.

The rejected author cannot settle which by pointing to the rejection. We need the objections and the answers. The existence of more than one journal, opportunities to share work, and the possibility of later replication matter because a single editorial decision should not become a final judgement on knowledge.

Publication begins a wider examination. Researchers try the method elsewhere, compare the result with other findings, or discover that an assumption does not hold. When a conclusion changes through this process, it is misleading to describe the change automatically as proof that science has failed. Sometimes the earlier work was poor. Sometimes the later work has made visible a limitation that was difficult to see before.

The institution's response matters. Does it correct the record, explain the consequences, and make the change easy to find? Or does it protect the publicity generated by the original claim while leaving a technical correction where few readers will encounter it?

\section*{From evidence to an instruction}

Suppose engineers agree that a bridge has deteriorated and that continued heavy traffic increases the risk of failure. Their agreement should carry great weight. A council that ignores it because some voters dislike restrictions is behaving recklessly.

Several decisions may still remain. Should the bridge close immediately? Could lighter vehicles cross for a limited period? Is a temporary repair possible? What would happen to emergency access and local businesses? How quickly could a replacement be built?

Some of these questions require further technical advice. Others involve deciding which burdens can reasonably be imposed on whom. The evidence about the bridge constrains the choices. It does not allow the council to claim that no judgement has been made by the people in charge.

An official who says that science required a particular decision may be compressing an argument she should make explicitly. What outcome is she trying to avoid? What alternatives were examined? Which costs were accepted? Under what conditions will the decision be reviewed?

Calling these political questions does not make them arbitrary. They can be discussed with evidence and reasons. It means that the official should take responsibility for them rather than presenting a recommendation as an order issued by nature.

Experts also have opinions as citizens. An engineer may favour spending much more on safety than the council is willing to spend. She should be able to argue for that preference. Her technical qualification does not make the preference disappear into a calculation.

This separation protects expertise as well as public judgement. If a policy fails because its financial assumptions were poor, it is unfair to blame the specialist whose narrower warning was correct. Conversely, an accurate technical warning does not excuse every policy that claims to follow it.

\section*{The right to see how a conclusion was reached}

People affected by an expert decision cannot always reproduce the underlying work. A patient cannot personally inspect every clinical trial. A resident cannot rebuild a structural model of a bridge. The demand for transparency must therefore include people able to use the material on their behalf.

Releasing a mass of undocumented files may satisfy a formal requirement while making examination practically impossible. Useful disclosure needs definitions, methods, relevant assumptions, and an explanation of changes. Where privacy or security prevents full public access, there should be a credible arrangement for independent scrutiny rather than a request that everyone simply trust the custodian.

The same applies when a public body buys a technical system from a private vendor. The vendor may claim that its method is commercially sensitive. The public body may claim that the method is too technical for it to explain. A person denied an opportunity can then find that nobody is willing to give an intelligible reason.

That arrangement should be challenged before it becomes routine. Purchasing a tool does not discharge the purchaser's responsibility for using it. If an institution cannot explain or review a consequential decision, sophistication of the tool is a poor defence.

The standard is demanding because dependence is unavoidable. We need experts, but needing them does not place us under a general obligation to admire their institutions. We can respect a difficult skill while asking how appointments are made, how mistakes are discovered, and what happens to people who point them out.

For the individual researcher this can be an uncomfortable public role. Years of patient work may be represented by a sentence she would never have written. She may be asked to defend a certainty she does not possess or to answer for an institution she does not control. A serious public discussion should give her room to explain the limits of her result without treating the explanation as a confession of uselessness.

Often the person we need is the one willing to say exactly what she knows, how she knows it, and where her answer ends. We should make that person easier to hear.

\chapter{What Was Actually Measured}

An advertisement says that a new teaching programme improved results by twenty per cent.

Before deciding whether to buy it, a school needs to know what results means. Did pupils answer more questions correctly? Did their examination marks rise? Did the proportion passing an examination increase? Were the pupils the same age, and were they compared with similar pupils who did not use the programme?

Even the percentage needs explaining. A rise from fifty to sixty marks is an increase of twenty per cent relative to fifty. A rise from fifty per cent of pupils passing to seventy per cent passing is an increase of twenty percentage points, or forty per cent relative to the original pass rate. An advertisement that moves carelessly between these descriptions can make a modest gain appear much larger.

Suppose the figure is accurate and clearly defined. We still need to know why it changed. Perhaps the new programme helped. Perhaps the pupils were older, had better attendance, or received more attention from the teachers running the trial. Perhaps the first test took place on an unusually bad day.

These questions are not clever ways of refusing to recognise success. They are necessary if another school wants to know what it can reasonably expect to happen there. A result becomes useful when we understand what produced it and under what conditions it is likely to recur.

\section*{The comparison that matters}

Imagine that two classes use different teaching methods. At the end of the year, one has higher marks. If that class already had higher marks at the beginning, the final difference does not tell us what the new method contributed. If parents chose which programme their children joined, the groups may differ in ways that the test score does not reveal.

Random assignment can help make comparisons fairer, particularly with enough participants and careful implementation. It does not solve every problem. Teachers may know which method is supposed to be the exciting one. Pupils may receive unequal attention. Some may leave the study. The outcome may measure a narrow skill that improves at the expense of something the test ignores.

Good research anticipates such difficulties, explains how it addressed them, and reports limitations that remain. A reader need not regard every limitation as fatal. The point is to understand the size and direction of possible error well enough to judge the conclusion.

There is a great difference between saying that the programme helped these pupils on this test and saying that it has transformed education. The first claim may be justified by a well-designed study. The second requires a much larger argument about different pupils, different teachers, costs, durability, and the purposes of education.

This enlargement often occurs after the research has been completed. A careful paper becomes a press release. A press release becomes an interview. The interview becomes a sales claim. The original authors may recognise each individual sentence while being uncomfortable with the confidence of the whole performance.

An institution should take responsibility for this translation. It is not enough to place qualifications in a paper and circulate an announcement that encourages people to ignore them. The short description should preserve the limitations most likely to affect the reader's decision.

\section*{How an interesting result can appear}

Suppose researchers examine many possible outcomes: different test scores, different age groups, boys and girls separately, short-term gains and longer-term gains. Some comparisons may show little. One may show an impressive difference.

That difference might reveal something important. It might also be a chance finding selected from a large collection of unremarkable results. To judge which, we need to know how many questions were tried and whether this particular comparison was planned in advance.

Statistical tests can help assess how unusual a result would be under specified assumptions. They do not relieve researchers of the need to describe what they tested. If enough comparisons are attempted, apparently surprising results become easier to find by chance. Reporting only the attractive one makes the work look more decisive than the full investigation was.

The familiar threshold of statistical significance is frequently asked to do too much. A p-value does not give the probability that a hypothesis is true. Nor does passing a conventional threshold establish that an effect is large enough to matter. The American Statistical Association's 2016 statement stresses these limitations and the importance of context and complete reporting.\footnote{Ronald L. Wasserstein and Nicole A. Lazar, ``The ASA's Statement on p-Values: Context, Process, and Purpose'', \emph{The American Statistician} 70, no. 2 (2016), pp. 129--133, doi:10.1080/00031305.2016.1154108.}

A tiny difference can become statistically detectable with a sufficiently large sample. A potentially important effect can remain uncertain in a small sample. Neither situation can be understood by treating a threshold as a stamp saying true or false.

This does not mean that statistics are a trick. It means that statistics are part of an argument whose assumptions and purpose must be visible. A calculation can be correct and the interpretation careless.

Planning the main analysis before seeing the results makes it harder to choose the question retrospectively. Exploratory work remains valuable: unexpected patterns are often where interesting questions begin. But an unexpected pattern should be presented as something discovered during exploration, with further checking required. It should not be given the appearance of having passed a test designed in advance to challenge it.

An honest account of the work may therefore sound less dramatic than a selective one. This is an incentive problem for journals, funders, and universities. If they reward only striking findings, researchers receive a practical lesson about which parts of the investigation deserve time and which parts can remain in a drawer.

\section*{Doing it again}

One way to test a striking result is to try to obtain it again. That sounds straightforward until another team begins asking what, precisely, was done.

Instructions may be incomplete. Software may have changed. Participants may differ. A measurement may depend on a detail the original researchers did not realise was important. Even when the second study is carefully conducted, the result can differ for more than one reason.

A failed replication therefore requires investigation. The original result may have been a chance finding. A procedure may not have been reproduced accurately. An effect may depend on conditions that were not understood. A later result can also be mistaken. The task is to narrow these possibilities rather than announce an immediate victory for either team.

The Open Science Collaboration's 2015 project attempted to replicate one hundred studies drawn from three psychology journals. It reported generally smaller effects in the replications and substantial difficulty reproducing the original findings.\footnote{Open Science Collaboration, ``Estimating the Reproducibility of Psychological Science'', \emph{Science} 349, no. 6251 (2015), aac4716, doi:10.1126/science.aac4716. \href{https://pubmed.ncbi.nlm.nih.gov/26315443/}{Article record and abstract}.} The project was an important examination of that body of work. It was not a random test of all psychology, still less a demonstration that every scientific field is unreliable.

Its broader institutional lesson is that publication should leave enough information for later checking and that later checking deserves resources. A system that rewards producing new claims much more than examining old ones can accumulate more confidence than its evidence warrants.

Replication is not the only useful check. Independent measurements, predictions in new conditions, different methods addressing the same question, and explanations that survive attempts to break them can all strengthen a conclusion. The appropriate checks depend on the subject. We cannot repeat the formation of a mountain range as we repeat a classroom task. We can still compare competing accounts with evidence.

The public need not master every method to understand this difference. A claim supported through several genuinely independent routes deserves a different degree of confidence from one supported by a single surprising paper. The number of times the paper has been quoted does not create those routes.

\section*{A useful model leaves things out}

A map that showed every stone, drain, and blade of grass would be inconvenient for finding the railway station. Its omissions help it serve its purpose. A scientific model also leaves things out, though deciding which omissions are acceptable can be much harder.

Consider a model of road traffic. It may estimate travel times from road capacity, usual demand, junctions, and patterns of departure. It cannot contain every driver's decision in advance. Its value depends on whether it captures enough of what matters for the question being asked.

A model adequate for comparing two junction designs may be inadequate for predicting how a new charge will change commuting habits over ten years. The second question requires assumptions about work, housing, transport alternatives, and people's responses to prices. More equations do not make those assumptions disappear.

The assumptions should be stated, and the reader should be able to see which ones strongly affect the answer. If small changes in a poorly known input produce large changes in the outcome, the final recommendation needs to reflect that sensitivity. Presenting only the central estimate would hide information relevant to the decision.

Different models can agree because they capture a robust relationship. They can also agree because they share an assumption that has not been tested. A collection of models is therefore not automatically a collection of independent witnesses. We need to know what they have in common.

Conversely, disagreement among models does not make all planning pointless. It may identify a choice that performs reasonably well across several plausible futures. It may suggest that a decision should be delayed, tested on a smaller scale, or designed so it can be altered. Uncertainty can inform the plan rather than merely appear in a paragraph apologising for it.

\section*{When people respond to the measurement}

Traffic predictions have a further complication. People can change their behaviour after receiving the forecast. If many drivers are told that a road will be empty and choose it, it may cease to be empty. The forecast has become one influence on the event it describes.

The same issue arises when institutions attach rewards to measurements. Suppose a call centre rewards staff for keeping calls short. Faster answers may improve the service. But staff could also shorten calls by rushing confused customers, transferring difficult cases, or discouraging questions. The average time falls while some customers have to call again.

A manager looking only at the average may report success. Adding the number of repeat calls would reveal part of the problem. Staff might then find ways to prevent a return call from being linked to the first one. Each new measure can improve oversight while creating another feature to manage.

The lesson is not to abandon measurement. Without records, customers' experiences may be even easier to ignore. The lesson is to retain contact with the purpose of the work. Listen to some calls. Examine unresolved cases. Ask what staff do when the target conflicts with helping the caller. Compare several kinds of evidence.

One number is attractive because it makes comparison easy. It may also conceal the disagreement that management needs to resolve. A call taking twenty minutes could be wasteful, or it could be the first successful attempt to solve a difficult problem. Duration alone cannot tell us which.

When the recorded outcome is produced partly by the rules attached to the record, later analysts must take those rules into account. The data are still observations. They are observations of people responding to an institution, not of a world left untouched by it.

\section*{Being precise about ignorance}

``We do not know'' can be an honest answer. It can also be an excuse for failing to examine something that could be known with reasonable effort.

An unknown value may not yet have been measured. A relationship may be disputed because available studies conflict. A future outcome may depend on choices that have not been made. These are different kinds of uncertainty and may require different responses.

If the missing information is the width of a doorway, someone can measure it. If the question is how a community will respond to a new transport service, observation of a trial may help. If the question concerns a rare event with grave consequences, waiting for enough examples to make the estimate comfortable may be an unacceptable way to learn.

Action therefore cannot always wait for a confident answer. But the decision to act should include the cost of being wrong in different directions. A cheap precaution that can be withdrawn is different from an expensive intervention that creates a lasting dependency. Reversibility matters, though it cannot always be purchased: some delays themselves close possibilities.

The responsible account states what is known, what is assumed, what remains doubtful, and why the proposed action is justified despite the doubt. It also describes what would lead to a change. A promise to review the decision means little unless someone has time, authority, and information to perform the review.

There is no contradiction between acting and remaining unsure. We do it whenever we make a journey, choose a job, or trust another person. The public version requires more explanation because we may be choosing for people who did not volunteer to share our estimate of the risk.

\section*{The result that survives disappointment}

Suppose the teaching programme does help, but by less than the advertisement suggested. The improvement varies between schools. It takes more teacher preparation than expected. Some pupils benefit; others need a different approach.

There is still something worth knowing. A modest, dependable improvement can be useful. It need not be promoted into a transformation to justify the work that discovered it.

Yet the reduced claim may disappoint the institution that sold the programme, the researcher associated with it, and the critic who hoped to expose it as worthless. All three have reasons to prefer a cleaner story. The pupils need a better decision about teaching.

We should not judge an inquiry only by how impressive its first announcement sounded. We should also judge what remains after the result has been checked, limited, and used. Some of the most valuable knowledge is what survives this loss of excitement.

It will rarely spare us judgement. The school must still decide whether the gain justifies the cost and whether another use of the money would help more. Good evidence makes that decision better informed. It does not relieve the people responsible of having to make it.

\chapter{When the Numbers Must Agree}

An electric car described as having zero emissions still has to be built.

Someone mines and processes materials, manufactures the battery, assembles the vehicle, transports it, and supplies the electricity used to charge it. The car also needs roads, tyres, and eventual disposal or recycling. None of this vanishes when the exhaust pipe does.

Yet a measurement of zero exhaust emissions can be perfectly accurate. A battery-electric car does not burn fuel in an engine and discharge the resulting carbon dioxide through a tailpipe. A rule concerned with that source of emissions has a reason to record zero.

The confusion begins when the description loses the words that identify what was measured. Zero at the tailpipe becomes zero-emission, and zero-emission is heard as an account of the entire vehicle.

This is not a difficult scientific dispute. It is a change in the scope of a sentence. We can recognise the change without concluding that electric cars are useless or that people who buy them have been deceived about everything. We can also object to the misleading description without volunteering to defend petrol engines.

Public argument becomes remarkably bad at allowing this modest position. One side hears criticism of the label as hostility to the technology. The other hears recognition of a real benefit as submission to propaganda. A measurement that ought to help a comparison becomes another test of membership.

\section*{An average is not an individual car}

The European Union's standards for new cars provide a further distinction. They regulate manufacturers' average certified carbon-dioxide emissions from new vehicle fleets. An EU-wide reference value is part of the calculation of manufacturer targets; it is not a fuel-consumption ceiling imposed separately on every vehicle.\footnote{European Commission, ``CO$_2$ Emission Performance Standards for Cars and Vans'', including the 2025--2029 reference values and the 2025 compliance flexibility. The numerical example in the text is simplified and is not a full manufacturer-compliance calculation. \href{https://climate.ec.europa.eu/eu-action/transport/road-transport-reducing-co2-emissions-vehicles/co2-emission-performance-standards-cars-and-vans_en}{Commission explanation}.}

The framework described after the 2025 changes used an EU fleet-wide reference of 93.6 grams of carbon dioxide per kilometre for new passenger cars in 2025--2029. The precise rules include adjustments and compliance provisions. For present purposes, the important point is that the number concerns an average of certified tailpipe values.

A simplified example shows the distinction. Suppose a manufacturer sells ten new cars. Five have certified tailpipe emissions of one hundred grams per kilometre. Five are battery-electric cars assigned zero in that measure. The simple average is fifty grams per kilometre.

No individual car in the example emits fifty grams at its tailpipe. Half emit one hundred and half emit zero. Nor does the average tell us the emissions created in making the cars or generating their electricity. It also does not tell us how far each car will actually be driven.

The average remains useful for the purpose for which it was designed. What would be misleading is to present it as something else: the physical performance of every car, the complete climate cost of the fleet, or the total emissions of the journeys it will make.

Regulations can change, and long-term targets are political decisions rather than physical constants. A reader comparing statements from different years needs to know which version of the rule is being discussed. A number without its date can be as misleading as a number without its unit.

It would be possible to design a different standard, perhaps giving more attention to life-cycle emissions. That would introduce further measurement choices about manufacturing, electricity supply, expected use, and vehicle lifetime. Wider accounting is desirable for some questions and more uncertain in some respects. The existence of those difficulties does not justify pretending that a narrow account is complete.

\section*{A boundary restored}

Once production and electricity are included, the question becomes comparative. Which vehicle produces fewer greenhouse-gas emissions over a specified lifetime, under specified conditions?

The answer depends on the vehicles being compared, battery size, manufacturing, mileage, and electricity supply. The European Environment Agency describes the typical electric vehicle in Europe as having lower life-cycle greenhouse-gas emissions than its combustion-engine equivalent, despite higher impacts in production.\footnote{European Environment Agency, ``Why Are Electric Cars a Good Alternative to Conventional Combustion Vehicles?'', drawing on its life-cycle assessment work. \href{https://www.eea.europa.eu/en/about/contact-us/faqs/why-are-electric-cars-a-good-alternative-to-conventional-combustion-vehicles}{Agency explanation}.} That is a qualified comparative conclusion, with substantially more meaning than a claim of zero emissions.

A critic who points to battery production has identified something the tailpipe figure excludes. To establish that the electric car is worse overall, the critic would also have to count the production and use of the alternative, including the continuing supply and combustion of fuel. Finding a cost on one side does not show which side has the larger total.

The same discipline applies to supporters. A favourable life-cycle comparison does not mean that all electric vehicles are equally sensible uses of materials, that electricity is free of other environmental effects, or that replacing a serviceable car immediately is always the best choice. The particular decision requires a particular comparison.

Our difficulty is often less with arithmetic than with the wish to possess a clean category. A clean car, a dirty fuel, a responsible purchase: each offers an answer that can be carried into conversation without the calculations. A better comparison usually leaves us with a preference and some remaining costs. That can be sufficient for action even if it is less satisfying as a description of ourselves.

\section*{When the total grows and the person does not benefit}

Now consider a different number: the size of an economy.

Imagine a small economy producing goods and services valued at one hundred million units a year for a population of ten thousand. Its output per person is ten thousand units. If the population and total output both increase by ten per cent, total output becomes one hundred and ten million, and the population becomes eleven thousand. Output per person remains ten thousand.

The economy has grown. The per-person average has not. Both statements are true.

Even output per person is not a statement about what each person receives. A large gain for some can coexist with stagnation or losses for others. Nor is measured production a complete account of welfare. Time, insecurity, environmental damage, unpaid care, and the quality of public services may matter in ways that a single output measure does not capture.

These limitations do not make economic statistics worthless. They explain what a speaker has to establish before moving from an increase in a national total to a claim that the typical resident is better off.

The distinction matters in discussions of migration, population, and development. Additional people can contribute to production, fill vacancies, create businesses, and also require housing, schools, transport, and care. The effect on the total is one question. The effect on existing residents, new arrivals, different occupations, and particular towns comprises several others.

A government that selects only the growing total can conceal pressures in the places receiving the largest changes. An opponent who selects only the most strained district can conceal improvements elsewhere. The solution is not to average the rhetoric. It is to examine the relevant groups and explain why the chosen comparison addresses the policy.

A number acquires political force partly through the people it allows us to stop seeing. We should be especially attentive when a speaker moves from everyone in aggregate to everyone as individuals without acknowledging the difference.

\section*{The balancing term}

The word net also performs work that a listener can easily miss.

If a business receives one thousand euros and spends eight hundred, its net receipt is two hundred. Nobody imagines that only two hundred euros entered the account. The subtraction is understood. In environmental claims, the corresponding subtraction can be less visible.

Net-zero carbon dioxide means that human-caused carbon-dioxide emissions are balanced by human-caused removals over a specified period.\footnote{IPCC, \emph{Climate Change 2023: Synthesis Report}, Annex I, Glossary, entries on net-zero emissions and carbon-dioxide removal. \href{https://www.ipcc.ch/report/ar6/syr/downloads/report/IPCC_AR6_SYR_Annex-I.pdf}{Report glossary}.} It does not mean that no carbon dioxide is emitted. Residual emissions may continue while removals balance them.

This can be a meaningful physical aim. It can also be represented through accounting of varying quality. The reader needs to know what is being removed, when, for how long, and how the removal is measured. A promise that a forest will absorb carbon in the future is different from a verified amount already removed and durably stored.

There is also a difference between removing emitted carbon and claiming that an activity avoided emissions that otherwise would have occurred. The latter claim depends on a comparison with a world we did not observe. Some such comparisons can be well supported. Others can be made generous by choosing an implausible account of what would have happened without the project.

Double counting introduces a simpler error: two parties claim the same reduction or removal. A pleasing certificate does not stop the atmosphere from receiving the actual quantity emitted. Accounting systems need rules precisely because the pieces of paper can circulate more easily than the physical result can be verified.

Calling every offset a fraud would spare us the work of examining them. Calling every certified offset a genuine cancellation would spare us the same work from the opposite direction. Neither approach is adequate when a claim of neutrality depends on the quality of the balancing term.

The honest sentence may therefore be inconveniently specific. Emissions were reduced by this amount; these emissions remain; this portion is being balanced through these verified removals; another portion depends on future performance. If the short label encourages a listener to believe more than that, it is doing more than summarising.

\section*{A possibility is not an outcome}

A product labelled recyclable presents a similar difficulty. The material may be capable of entering a recycling process. That does not establish that it will be collected, separated from other waste, processed successfully, and used to make something else.

Each step requires equipment, energy, labour, and a reason for someone to undertake it. A product can be technically recyclable in a facility unavailable to its buyers. A collection scheme can gather material for which there is little demand. Contamination can make a theoretically simple process impractical.

These facts do not make recycling pointless. They help identify where a scheme is failing and what would improve it. A manufacturer could use a material that is easier to separate. A town could collect it more reliably. A purchaser could prefer a product that lasts longer or can be repaired. The useful question is what happens to the material, not how reassuring the symbol looks.

In much political language, can becomes will without an argument in between. A technology can remove carbon, therefore the future plan can rely on a large quantity being removed. A building can accommodate a service, therefore the service is available. A person is eligible for training, therefore the shortage of skills is being addressed.

Sometimes the missing steps are being organised competently. Sometimes nobody has responsibility for them. The announcement alone cannot distinguish the cases.

A promise should therefore describe the steps that connect capacity to use. If a necessary step depends on a future invention, an unapproved budget, or people behaving differently from how they currently behave, the dependence belongs in the account. It is not a disloyal detail.

\section*{Two plus two}

In \emph{Nineteen Eighty-Four}, Orwell uses elementary arithmetic to express a minimum resistance to power: the ability to say what is true even when authority demands another answer.\footnote{George Orwell, \emph{Nineteen Eighty-Four} (1949), part I, chapter 7, and part III, chapter 2.} The scene is deliberately stark. We should not compare every ambiguous public statistic with torture.

The ordinary difficulty is subtler. The sum may have been performed correctly while the things being added or subtracted have changed. An average becomes a statement about each member. A tailpipe figure becomes a life-cycle claim. A future possibility becomes a completed result. A target becomes a forecast.

We can challenge these changes without describing the entire statistical system as fraudulent. Indeed, the challenge often depends on the system's own definitions and records. A footnote, an appendix, or an original data table can restore what the slogan removed.

This work lacks the satisfaction of announcing that everything is a lie. It also produces more useful objections. Instead of saying that a government invented a number, we may be able to show that it used a real number to answer a different question. The correction then becomes precise enough for someone to make.

Of course, some numbers are fabricated. Some results are knowingly concealed. Direct deception should be called what it is when there is evidence for it. But if we use the same accusation for fabrication, carelessness, and a legitimate measurement with limited scope, we make it harder to establish which occurred.

The distinction also protects us from a critic who wins our trust by exposing one misleading claim and then offers an equally misleading alternative. The first achievement does not certify the second. We must continue checking even when we like the person doing the checking.

\section*{What the correct number leaves undecided}

Once the labels are restored, the argument is not over. We may know the likely emissions of two cars and still disagree about transport policy. We may know the change in output per person and still disagree about admissions policy. We may know that a recycling scheme costs more than expected and decide that some other benefit justifies it.

This is not a failure of arithmetic. We had asked it to decide more than it could.

Numbers can establish constraints, reveal costs, and expose impossible promises. They cannot by themselves tell us which loss we should accept when worthwhile purposes conflict. That judgement needs to be made openly, with the people bearing the loss present in the discussion as more than a small remainder in a table.

The best outcome may be an honest account of an imperfect choice. There is no need to improve its appearance by removing costs that remain. People who must live with those costs will notice them soon enough. They may forgive a difficult decision more readily than an official explanation in which their difficulty has ceased to exist.

\part{Obligations We Did Not Choose}

\chapter{Born Owing}

A child has done nothing wrong by arriving in the world.

She nevertheless arrives in a world already affected by other people's conduct. The house in which she sleeps depends on energy and materials. The country in which she grows up has a history of laws, conflicts, possessions, and debts. Her opportunities will be shaped by decisions made before her parents met.

Some of these inheritances are gifts. Others are burdens. She did not choose either kind.

As she grows older, she may be told that she is responsible for what she has inherited. Sometimes this means that she should help maintain useful institutions and repair damage she is able to affect. Sometimes it means that she ought to feel guilty for advantages, injuries, or acts that preceded her existence. The same word, responsible, is being used for different demands.

The distinction is practical. A person can have a duty to deal with a problem she did not cause. If water is leaking through the ceiling of her house, discovering that a previous owner installed the pipe will not stop the leak. It may matter greatly in deciding who should pay, but someone still needs to turn off the water.

That does not make her guilty of installing the pipe. Nor does acknowledging the leak authorise every demand a neighbour chooses to make of her. The problem needs a description: what happened, who contributed to it, who can repair it, and what obligations follow.

When these questions are skipped, inherited responsibility can become a particularly effective accusation. The accused cannot undo the event, cannot resign from having been born, and may never be told what would count as an adequate response.

\section*{What the old doctrine says}

The phrase original sin is often used as if Christianity simply declared every newborn personally guilty of a remote ancestor's offence. That is too crude an account of a doctrine whose meaning has been disputed within Christianity itself.

Genesis tells a story of disobedience and expulsion. The first humans eat the prohibited fruit, discover their nakedness, hide, and shift blame. Their descendants inhabit life outside the garden. The story provides material for later doctrines; it does not present a completed system of inherited personal guilt.

Other biblical passages make individual responsibility explicit. Ezekiel 18 rejects the transfer of a father's guilt to his son. In Romans 5, Paul draws a far-reaching connection between Adam, sin and death, and between Christ, justification and life. The relation between an inherited human condition and each person's responsibility becomes a problem that later theology develops in different ways.\footnote{Genesis 3; Ezekiel 18:19--20; Romans 5:12--19. These passages address different aspects of disobedience, individual responsibility, and the Adam--Christ comparison. \href{https://www.biblegateway.com/passage/?search=Ezekiel\%2018&version=KJV}{Ezekiel 18, King James Version}.}

Augustine's account had a major influence on Western Christianity. But even the Catholic catechism distinguishes the inherited condition from a personal act committed by Adam's descendants. It describes original sin as a state rather than personal fault in each descendant, and places the doctrine within an account of grace and redemption.\footnote{\emph{Catechism of the Catholic Church}, paragraphs 404--406, including its distinction between an inherited state and personal fault and its account of Augustine's role in doctrinal development. \href{https://www.vatican.va/content/catechism/en/part_one/section_two/chapter_one/article_1/paragraph_7_the_fall.html}{Official text}.}

One can reject the theology. What one should not do is manufacture an easier theology to reject. If we are interested in how guilt is used, accuracy about the doctrine is more useful than the pleasure of making believers sound absurd.

The question for this book is narrower. What happens when an inherited condition is turned into a standing accusation against people who did not choose it? What happens when the accusation can be renewed indefinitely, while the authority making it retains the power to decide whether any response is sufficient?

Those questions can be asked about religious and secular institutions alike. They concern a use of guilt, not a claim that every institution has identical beliefs.

\section*{Consequence without personal blame}

Imagine someone inheriting a piece of land contaminated by an old industrial operation. She did not choose the operation and may not have known about the contamination when she inherited it. The soil is still contaminated.

There are several questions to settle. How dangerous is it? Who was responsible for the original activity? What claims can be established against those who benefited? What repair is possible? What should the present owner do if the contamination is spreading beyond the property?

These are not all questions about her character. She may have present obligations because she controls the land, because she benefits from owning it, or because she is in a position to prevent further harm. The justification for each obligation needs to be stated. Calling her personally guilty of the original act would not make the cleanup plan better.

The reverse evasion is equally plain. ``I didn't do it'' may be a correct statement about the past and an inadequate answer to a present request for action. A person can be innocent of causing a danger and negligent in allowing a preventable danger to continue after discovering it.

The same distinctions matter when whole societies inherit consequences. Public institutions outlive their members. A state can retain assets and obligations through changes of government and generation. A historical wrong can continue to affect people who are alive now. None of this requires a theory in which guilt is transmitted biologically.

There will be difficult disputes about continuity, evidence, beneficiaries, victims, and remedies. Time can make them harder by removing records, dispersing families, and combining old advantages with later effort. These difficulties are reasons for careful judgement. They are not reasons to declare every present claim valid, or every past injury irrelevant.

An accusation based only on a person's ancestry can obscure precisely the connections that need to be established. It may attach blame to someone with little power while overlooking an institution that retains the relevant property. It can also encourage the person accused to deny the entire history, because acknowledgment has been made to sound like consent to unlimited liability.

A more exact account gives acknowledgment a better chance. We can admit what occurred and examine what follows without pretending that every possible remedy has already been justified by the admission.

\section*{The child with a footprint}

Environmental accounting supplies a contemporary version of this difficulty. A child will use food, housing, transport, materials, and energy. Some of those activities will produce emissions and other harms. Estimating them can help us understand the consequences of different ways of living.

The estimate becomes a moral accusation when the child's existence itself is treated as the offence. She is described as another burden before she has had the opportunity to do anything with the life she has been given.

This goes beyond the ordinary observation that sustaining people requires resources. Every proposed society must face that fact. The question is whether people are still among the beings for whose welfare our institutions exist, or whether an accounting measure has quietly made their disappearance look like the preferred outcome.

Even the physical account can become careless. Breathing releases carbon dioxide, but the carbon comes from food within the recent biological carbon cycle. It is different from introducing additional fossil carbon by burning coal, oil, or gas. Food production can cause substantial emissions; the fact of exhalation does not establish the same process.\footnote{MIT Climate Portal, ``Does the Carbon Dioxide That Humans Breathe Out Contribute to Climate Change?'' \href{https://climate.mit.edu/ask-mit/does-carbon-dioxide-humans-breathe-out-contribute-climate-change}{Explanation of the biological carbon cycle}.}

Similarly, drinking water is a human need. Its provision can involve scarcity, pollution, energy, and difficult allocation. These are reasons to improve systems and examine waste. They do not make thirst a vice.

Heating a home can depend on very different energy sources and buildings. A wealthy household choosing excessive space and inefficient equipment faces different options from a tenant unable to alter a poorly insulated flat. A common accounting unit can be useful for comparing emissions. It cannot by itself assign equal moral responsibility for every unit.

Responsibility depends partly on knowledge, available alternatives, power, and the cost of changing conduct. Those differences do not absolve everyone. They help identify who can do what and what can reasonably be required.

If a policy requires cooperation, making these distinctions is more than a kindness. People are less likely to trust an authority that condemns a necessity while offering no practicable alternative. They may conclude that the language of responsibility is a way of blaming them for systems other people control.

\section*{Paying to feel better}

There is a market for relief from guilt. A purchaser buys a product, makes a donation, or pays for a certificate and receives an assurance that a troubling activity has been compensated.

Sometimes the payment finances something valuable. Sometimes it finances a story about the purchaser. The two may coincide, but they should not be confused.

The familiar comparison with indulgences can be useful if it is handled accurately. In Catholic doctrine an indulgence concerns temporal punishment for sins whose guilt has already been forgiven; it is not simply permission to commit a sin or a purchase of forgiveness.\footnote{\emph{Catechism of the Catholic Church}, paragraph 1471, definition of indulgences.} Historical selling practices and popular understandings raise further questions. The point of the comparison is the attraction of a transaction that appears to settle a moral burden.

An emissions allowance and a voluntary offset also perform different functions. An allowance under a cap-and-trade system belongs to a regulated limit. An offset makes a claim about a reduction or removal elsewhere. Their effectiveness depends on how the particular system is designed and enforced. The person paying need not become morally better for a well-designed system to reduce pollution.

What deserves scrutiny is the additional promise that payment makes the purchaser's conduct clean. A certificate cannot remove consequences beyond the work it actually verifies. If a payment supports a useful project while the original emissions remain partly unbalanced, the useful project should be credited and the remaining emissions should remain in the account.

The distribution of this reassurance can also be offensive. People with money can buy recognised signs of responsibility while people with fewer alternatives meet the policy chiefly through higher prices. The first group may then lecture the second about commitment.

The answer is not that wealth makes every payment suspect. It is that the physical benefit, the economic burden, and the flattering public description should be examined separately. A genuine reduction is valuable even when the donor is vain. A virtuous appearance is not a reduction.

\section*{A demand that cannot be completed}

A workable obligation normally has some account of what fulfilling it would involve. Repay a stated debt. Repair identified damage. Stop a specified harmful practice. Contribute under an agreed rule. Details can be disputed, but there is a subject for the dispute.

A permanent accusation is different. Every response confirms the authority of the accuser while leaving the accused's condition unchanged. Compliance proves that the accusation was justified. Hesitation proves moral failure. A request for limits proves that the accused has not understood the depth of the obligation.

Such an arrangement need not be deliberately planned. An institution may expand its demands because it keeps discovering real problems. A campaigner may sincerely regard each new requirement as modest. The person subject to them experiences a process with no identifiable end and no clear distinction between improvement and continued inadequacy.

The question to ask is what result the demand seeks and how the demand relates to that result. If a compulsory statement of remorse does nothing to repair an injury, its purpose may be humiliation, reassurance, or the demonstration of obedience. Those purposes should not be concealed behind the word repair.

Public remorse can sometimes recognise a wrong in a way victims value. It can correct a record and end an official denial. But its value is diminished when the speaker is asked to claim feelings she does not have about actions she did not commit. A truthful institutional acknowledgment is stronger than a compulsory personal confession.

People should be able to assist with repair without accepting an invented account of their own guilt. They should also be able to question a proposed remedy without being treated as defenders of the original wrong. These permissions matter most when the wrong is serious, because seriousness makes a bad remedy easier to protect from criticism.

\section*{What remains ours to do}

We do not choose the point in history at which we arrive. We inherit useful knowledge, damaged places, loyalties, debts, and institutions whose origins we may find disturbing. There is no practical way to begin entirely outside that inheritance.

Refusing imaginary guilt does not free us from the world it describes. The polluted land still requires attention. The poorly insulated home still consumes energy. A surviving injustice still affects someone. We may have obligations because of what we control now, what we continue to benefit from, or what we can prevent.

These obligations can be demanding enough without making existence itself a fault. Indeed, they become easier to discuss when the accusation is removed. We can ask which action helps, whether its cost is proportionate, and who should contribute. A person need not hate her own life to take those questions seriously.

There will be losses no remedy can reverse. An apology cannot return a dead person's years. A restored landscape is not the same place remembered by someone who knew it before the damage. Money can help survivors without making the history come out right.

Acknowledging these limits is part of honest repair. It prevents the people paying from purchasing a false sense that the matter is finished. It also prevents those making the claim from promising a restoration they cannot deliver.

The child with whom this chapter began will eventually have decisions of her own to answer for. We should give her a world in which she can distinguish those decisions from the circumstances of her birth. Otherwise we will have taught her either to confess without understanding or to deny without looking. Neither is a useful preparation for responsibility.

\chapter{Before the Race Begins}

Two young people earn the same salary. One receives money from her parents for a deposit on a flat. The other pays rent and tries to save from what remains.

They may both work hard. They may both make sensible decisions. A few years later their financial positions can be substantially different without either having changed jobs. One has had help purchasing an asset and may have accumulated equity. The other has paid for somewhere to live and may still be saving the deposit.

The parents who gave the money can have earned it honestly. They may have postponed holidays, worked extra hours, and hoped for years to give their daughter a more secure beginning. There is no need to describe their gift as a conspiracy against the other young person.

Yet the difference between the recipients is real. The gift has changed what one can afford, how much risk she can take, and what choices she can postpone. It may influence whether she can start a business, care for a relative, or leave an unpleasant employer. These are advantages in a life, not merely figures in an account.

We tend to approve of parents doing this for their children. We also tend to approve of a society in which a child's prospects do not depend heavily on which parents happen to be hers. The two approvals do not fit together as comfortably as our political language suggests.

If parental help makes no difference, much of the effort to provide it has been wasted. If it makes the intended difference, children have unequal advantages. Calling the help love describes why it was given. It does not remove what it does.

\section*{Three different promises}

The phrase equal opportunity can mean several things. At its simplest, it means that a person is allowed to apply. A job is not reserved by law for members of particular families. A child is not prohibited from attending a school because of her ancestry. This is an important protection, and its absence can make every later achievement irrelevant.

A second meaning concerns the fairness of selection. An examination should not be rigged. A hiring panel should not secretly reserve a vacancy for the director's nephew. Candidates should be assessed against requirements relevant to the work. The formal right to apply is worth less when the result has been arranged beforehand.

A third meaning concerns the chances of becoming a strong applicant. Children with similar potential should not have radically different prospects because one has adequate food, stability, teaching, and time while the other does not. Here the question begins long before the application form. This stronger concern is central to discussions of fair equality of opportunity, including Rawls's account.\footnote{John Rawls, \emph{A Theory of Justice}, revised edition, Harvard University Press, 1999, especially the discussion of fair equality of opportunity. \href{https://www.jstor.org/stable/j.ctvkjb25m}{Publisher's edition on JSTOR}.}

These promises require different actions. Publishing vacancies can improve access to information. Independent marking can make an examination fairer. Neither can give a candidate the years of preparation she did not receive.

Conversely, improving preparation does not make every selection procedure fair. A well-educated outsider may still find that the decisive conversation took place at a private dinner. A society can spend generously on schools while preserving a closed circle at the point where authority is distributed.

The language becomes dishonest when success on one promise is used to claim that the others have been fulfilled. ``Anyone can apply'' does not establish that anyone has a realistic chance. ``We provide schooling'' does not establish that the schooling offers comparable preparation. ``The examination is impartial'' does not establish that the differences it records arose impartially.

This is not an argument against examinations or demanding standards. A bridge should be designed by someone who understands engineering. The fact that learning engineering was easier for one candidate does not make it sensible to give the work to someone unable to do it. But selecting the competent person today and improving access to competence for tomorrow are different tasks. We need both.

\section*{What money buys before the inheritance}

The word inheritance brings to mind a will and a transfer after a death. Much of what matters has already been transferred by then.

A quiet place to sleep affects the next school day. A secure home makes it easier to remain at the same school. Parents who can take time off can attend meetings, arrange appointments, and notice a difficulty before it becomes severe. Money can buy equipment, transport, lessons, and time in which a child is not required to earn.

It can also buy recovery from mistakes. A young adult who chooses an unsuitable course may try another. A first business can fail without making the family homeless. An unpaid period of exploration can be described as finding a vocation rather than failing to contribute to the rent.

The resulting confidence may later be mistaken for a quality the person generated alone. She is willing to take risks, employers observe. She thinks ambitiously. She does not become anxious when plans change. Some of these qualities may be personal achievements. Some may be easier to develop when failure has a place to land.

This does not establish a simple ranking in which wealthy children are secure and poor children are damaged. Families differ. Wealth can accompany neglect, coercion, and severe anxiety. Modest means can accompany affection, resourcefulness, and high expectations. A statistical tendency would not decide the character of a particular home.

The practical point is that an account of opportunity cannot stop at current income or formal eligibility. It must ask what support was available during development and what safety remains available when choices go badly.

An estate tax introduced at the end of the parents' lives may affect the next transfer. It cannot remove the earlier education, knowledge, or connections. A person who claims to have solved inherited advantage through one tax has chosen the most visible channel and ignored several others.

\section*{How a small difference becomes several differences}

Consider two families each able to save a little. One begins with fifty thousand euros already available. The other begins with none. Even if they make identical additional savings, the first has more opportunity to earn a return and more capacity to absorb an unexpected expense.

A simple illustration makes the point without pretending to forecast a real investment. At a constant annual return of four per cent, fifty thousand becomes about seventy-four thousand after ten years if nothing is added or withdrawn. Actual returns vary, taxes and inflation matter, and investments can lose value. The example merely shows that an initial asset can produce another advantage while both families are occupied with earning their wages.

Non-financial advantages can develop through a sequence too. A child receives help with reading, begins to enjoy books, reads more, and finds other subjects easier. Teachers notice her progress and encourage her. Later, her performance makes her eligible for opportunities that provide further encouragement.

Another child may encounter difficulty early, receive little useful help, and begin to avoid reading. Work becomes harder, frustration becomes familiar, and teachers come to expect less. There is no inevitable outcome in either sequence. A good teacher, a friend, a change at home, or the child's own persistence can alter it. But waiting until the final examination to discuss fairness misses much of how the candidates arrived there.

The same can happen across generations. Success provides resources and knowledge that help the next child succeed. The adults selecting promising candidates may recognise the manners, interests, and speech they learned in their own families. They can mistake familiarity for ability without consciously intending to exclude anyone.

This is why the apparently small advantage deserves attention. Its later effect can exceed its initial size because it changes access to further advantages. It is also why an equal cash grant at eighteen, though potentially useful, cannot be described as a complete reset. It provides resources at eighteen. It does not provide the childhood that would have made different uses of those resources possible.

\section*{Some advantages improve a life; others exclude a competitor}

If a parent teaches a child to cook, the child gains a useful skill. Other children are not thereby less able to cook. They might even benefit from eating what she makes.

If a parent buys intensive coaching for an examination offering ten places, the child may gain knowledge, but she may also improve her position relative to other applicants. One additional successful applicant means one other applicant misses a place. The benefit has a competitive element that cooking dinner need not have.

The distinction is not absolute. Better cooking can become a qualification for a scarce job. Examination preparation can produce knowledge useful beyond the competition. It nevertheless helps us see why not every parental advantage has the same effect on other children.

Some educational spending produces more capable people. Some largely rearranges who receives a fixed number of prestigious opportunities. A society should care about the difference. If families must buy ever more coaching merely to retain the same relative position, a great deal of effort may produce little improvement in what the candidates can actually do.

The institutions controlling the competition have choices. They can examine whether expensive preparation predicts the ability the position requires. They can use more than one route of entry. They can expand good opportunities where expansion is feasible, reducing the importance of a single narrow entrance.

None of this means that every desired position can be made abundant. Some roles really are limited. But scarcity can be manufactured as well as inherited. A profession may demand an expensive credential partly to protect existing members. An employer may continue recruiting from a small set of institutions because doing otherwise requires more work. The prestige of the old route can conceal the possibility of a competent alternative.

A fairer contest is not necessarily an easier contest. It may be more demanding about the skill that matters and less demanding about the money, manners, or connections used to reach it.

\section*{When advantage buys the rules}

A larger home and influence over a planning decision are both things wealth may help secure. They are not equivalent kinds of advantage.

The home provides comfort and perhaps a better location. Influence over the decision can determine what other people are allowed to build, how much their property is worth, and whether they can enter the market. One benefit changes a private life. The other changes the conditions under which others must live.

This is where inherited wealth becomes a question of political equality as well as distribution. A family may be entitled to own property without being entitled to purchase an exemption from common rules. A successful business may continue within a family without acquiring a right to prevent new competitors from operating.

The line is not always easy to draw. Owners should be able to explain how proposed rules affect them. Experience in an industry can be relevant to legislation. A donation to a public cause may reflect sincere belief. But access should not become a private arrangement in which the richest parties shape the rule while everyone else discovers its consequences later.

Transparency about meetings, conflicts, and financial interests can help. So can independent enforcement and rules that apply to insiders as well as outsiders. These measures do not remove differences in wealth. They restrict some of the ways wealth can be converted into authority over other people.

The family gift with which this chapter began is therefore not a sufficient image for every inheritance. A deposit that gives a young adult more security and a fortune that gives its recipient control over important institutions raise overlapping but different questions. Defending the first does not settle the second.

Nor does condemning the second justify treating every saved euro as an attempted purchase of power. A policy needs distinctions proportionate to what a transfer enables. Without them, the debate becomes a contest between a sentimental picture of loving parents and a sentimental picture of predatory heirs. Both pictures are selected to prevent awkward cases from entering.

\section*{What public institutions can change}

The most urgent task is to make childhood less dependent on a family's ability to purchase necessities. Reliable food, safe housing, healthcare, and effective education matter before a society begins arguing about the precise ranking of applicants to an elite university.

These services require resources and competent administration. Announcing a right does not supply a teacher, repair a flat, or make an appointment available. A poor service labelled universal can leave well-connected families finding ways around it while others receive little beyond the label.

Good common provision can change the role of parental effort. Parents can help with reading without first having to buy access to a functioning school. They can care for a sick child without having to solve the entire system of healthcare privately. The family remains important, but fewer basic conditions depend on its purchasing power.

Institutions can also make their procedures easier to navigate. A capable child should not need a parent who knows which official to telephone. Clear applications, understandable criteria, timely advice, and a route to challenge a mistake reduce the value of inherited knowledge about how to obtain a hearing.

This work can be undramatic. A well-written form will not make a society equal. But a confusing form can exclude someone as effectively as a rule that openly refuses her. The cumulative effect of small administrative barriers matters, especially when people who design them have never faced them without assistance.

There are costs and trade-offs in every approach. Universal services may use public money on families able to pay privately. Targeted services can be cheaper on paper while creating complicated tests, stigma, or gaps. Private options can provide variety and also allow influential families to disengage from common provision. These are questions for design and evidence, not consequences that disappear once the intended beneficiaries are named.

\section*{The advantage we want for our own child}

The inheritance argument becomes less comfortable when we imagine our own family as the beneficiary.

It is easy to favour impartial selection until a child we love narrowly misses a place. We know details the selection process cannot contain: an illness before the examination, an unusual talent, a year of effort, a disappointment that seems more than she can bear. We want someone to look again.

Sometimes looking again is justified. Procedures make mistakes. An appeal can correct one. But there will be another child just below the threshold with a history equally important to someone else. A system should not provide a more sympathetic hearing merely because one parent is more persuasive or better connected.

This is the difficulty that no general declaration removes. The parent sees a particular child and has good reasons to care especially about her. The public institution must make a decision among children who are equally particular to their own families.

It would be cruel to demand that the parent stop caring. It would be unjust to let that care decide a public allocation without limits. The rules must hold even when the loving parent regards them as intolerably impersonal.

There is no reason to expect this conflict to become painless. A family may do everything permitted and still encounter disappointment. A society may improve opportunity and still leave large differences that the children did not earn. We can judge an arrangement better than another without declaring the surviving differences deserved.

The young person still renting at the start of this chapter is not necessarily a victim of a particular wrong committed by the parents who helped her contemporary. She nevertheless has reason to reject a public story in which their different positions reveal only effort, thrift, and character. The gift happened. A decent account of their lives should be able to mention it.

\chapter{What Parents Can Give}

A parent reading to a child is doing several things at once. She is helping the child learn words. She is sharing a story, answering questions, and spending time with someone who wants her attention. The child may remember the evening long after forgetting the story.

The reading may also improve the child's prospects at school. Another child, whose parents are exhausted, absent, unable to read, or simply uninterested, may not receive the same help.

If our only question were how to eliminate unequal preparation, the difference would look like a problem to be removed. But prohibiting the first parent from reading would destroy something valuable without giving the second child an equivalent relationship. The loss would not be adequately described by a change in examination scores.

This is why the family cannot be understood merely as a system for distributing advantages. People within it are attached to one another, and those attachments can be part of what makes their lives good. Political judgement has to take those goods seriously while asking what costs particular forms of parental help impose on others. Harry Brighouse and Adam Swift make this conflict central to their account of parent--child relationships.\footnote{Harry Brighouse and Adam Swift, \emph{Family Values: The Ethics of Parent-Child Relationships}, Princeton University Press, 2014, especially the discussions of relationship goods, equality, and conferring advantage. \href{https://academic.oup.com/princeton-scholarship-online/book/23923}{Publisher's edition}.}

The question is not whether every practice called a family value should be protected. Some practices injure children or prevent them from developing lives of their own. The question is what we would lose by treating every unequal benefit within a family as something requiring public correction.

\section*{Attention that cannot simply be purchased}

A young child needs more than a series of services. Someone has to notice when the services do not meet the need. A coat is too small. The child is frightened at school. The explanation given by an adult does not quite match the child's behaviour. Care includes returning to these matters when no appointment has been made.

Parents are not uniquely capable of such attention. Grandparents, foster parents, other relatives, and people outside a conventional household can provide it. Some children receive much better care from them than from the adults who brought them into the world. What matters is the dependable relationship, not a sentimental picture of one family form.

An institution can organise important forms of continuity too. A good school, a stable care arrangement, or a trusted professional may change a child's life. But the institutional task becomes harder when there is no particular adult who feels responsible for following the child across the separate parts of that life.

The value of that responsibility is partly that it survives inconvenience. A child can fail, become ill, behave badly, or disappoint expectations without losing the relationship altogether. The adult may be angry and still return the next morning. This is different from being rewarded only while one performs well.

Of course, families do not always offer this security. Some withdraw affection as a punishment or use dependence to control. The fact that a relationship can provide a distinctive good does not establish that every instance provides it. It gives us a reason to protect the good where it exists and to recognise the loss where it does not.

An argument about equal opportunity should therefore ask how more children can receive dependable attention. It should not assume that the problem is solved by preventing those who already receive it from benefiting.

\section*{The child is not the parent's project}

Parental sacrifice can become a claim on a child's future. A parent pays for lessons, schooling, and opportunities, then expects a profession, a marriage, a way of living, or a public success in return.

The expectation may be understandable. Years of effort were organised around a hope. But a child is not an investment obliged to produce the return that justified the investor's restraint. The fact that a parent gave up things does not give the parent ownership of the adult who emerges.

This is one of the limits that sentimental accounts of family duty avoid. We are invited to admire the sacrifice without asking what it demands of its recipient. A father who insists that his daughter join the family business may regard himself as offering security. She may experience an obligation to continue a life she would not have chosen.

The business might offer an excellent opportunity. The daughter might later regret refusing it. Those possibilities do not settle who should decide. A meaningful adulthood includes some capacity to make choices that the people who raised us dislike.

Parents also face a genuine difficulty here. They have experience the child lacks. A teenage ambition can be poorly informed. A debt, a harmful relationship, or a decision to abandon education can have consequences the young person underestimates. Respecting developing independence does not require adults to withhold advice or provide unlimited money for every plan.

What changes over time is the justification for control. Protecting a small child from a danger she cannot understand differs from compelling a competent adult to avoid a life her parents find embarrassing. The vocabulary of care can cover both. We need to examine the age, the decision, the dependence, and the means of pressure being used.

The most generous inheritance may include permission to depart from the expectations attached to it. That permission can be painful to give. It means accepting that the life one helped make possible will not necessarily preserve one's own habits or ambitions.

\section*{Teaching a way of life}

Families pass on more than skills. They teach what deserves respect, which events matter, how to respond to disappointment, and whom to trust. A child learns these things partly through explicit instruction and partly through watching what the adults actually do.

A household may teach a religion, a language, a political loyalty, or a demanding conception of duty. Another may teach independence from all of them. There is no upbringing entirely without values. Even the declaration that children should decide everything for themselves supplies a view about what matters and when they should exercise judgement.

The public question is therefore not whether parents may influence their children. They inevitably will. It is how to protect the child's future ability to understand alternatives, obtain useful knowledge, and participate in a wider world.

A family can love its language without preventing a child from learning the language needed outside the home. It can teach a faith while allowing the young person to encounter disagreement. It can preserve memories of injustice without making inherited enemies the organising purpose of the child's life.

These distinctions do not always satisfy the adults. Someone who believes that departure from a faith endangers a soul may not regard exposure to alternatives as neutral. A political family may see disagreement as disloyalty to people who suffered. A parent can sincerely believe that closing a child's options is the only responsible way to protect her.

Sincerity does not make the child's interests disappear. A public order has reason to require enough education and freedom for the eventual adult to enter or leave associations on more than nominal terms. The child should not have to escape into ignorance because all useful preparation was made conditional on remaining obedient.

Public authorities can abuse this responsibility. Officials may mistake unfamiliar customs for neglect or impose a narrow image of respectable life. That risk calls for clear limits, evidence, and review. It does not establish that whatever happens within a family is beyond examination.

The aim is a child capable of becoming an adult among others. It is a more limited aim than making every household resemble the households of those who write the rules.

\section*{The telephone call that changes a decision}

Some parental help consists of knowing whom to ask.

A child has been placed in a course that does not suit her. One parent assumes the decision is final. Another knows that there is a review process, obtains the relevant form, contacts the right person, and presents the evidence effectively. The second child gets another hearing.

The result may be entirely justified. The parent may have corrected a genuine mistake. It would be perverse to insist that she remain silent so that her child suffers equally with those whose parents do not know the procedure.

Yet the institution has a problem if the chance of correction depends strongly on this kind of parental knowledge. A mistake should be reviewable because it may be a mistake, not because a child's family can make itself troublesome.

The repair should begin with the process. Explain decisions clearly. Tell people how to challenge them. Provide assistance to those who cannot formulate a case. Ensure that someone actually reads the appeal. These changes reduce the need for a private advocate without forbidding parents to advocate.

The distinction becomes sharper when the telephone call secures something the evidence would not justify. An influential parent may expect an exception, imply consequences for refusal, or speak to someone who owes a favour. The vocabulary remains concern for a child. The action has become a demand that common rules bend around a family.

A parent can have a powerful motive for making that demand. An official should still refuse it. The public role exists partly to withstand precisely those requests that feel irresistible in private life.

This can place the official in an uncomfortable social position. She may know the family, sympathise with the child, or depend on the parent professionally. Good rules help only if the institution supports the person applying them. An employee instructed to be impartial and left alone when an important family complains has not been given a fair chance to do her job.

\section*{What should not be for sale}

The difference between paying for a music lesson and paying for a public appointment is easy to state. Many harder cases lie between them.

Private tutoring can help a child understand mathematics. It can also provide extensive training in the peculiarities of a selection test, giving purchasers an advantage poorly related to later competence. Private schooling can offer a distinctive education or a refuge from a failing institution. It can also function as an entrance to a restricted network.

These possibilities should affect our judgement of the practice. It is not enough to say that parents naturally want the best. Natural motives can produce harmful arrangements. Nor is it enough to say that the practice creates inequality. The question is what kind, how much, and at what cost the practice could be restricted or its effects reduced.

Suppose an expensive qualification has become the expected entrance to a profession, although competent workers can be trained through a less expensive route. Opening and respecting the alternative route may reduce inherited advantage more directly than a general denunciation of ambitious parents.

Suppose a school admits candidates through an examination highly responsive to purchasable coaching. It can ask whether a different assessment would better identify ability to benefit from the education. It can also make preparation available more widely. Neither option will be perfect, but both address a mechanism rather than a moral category.

Some goods should remain unavailable for private purchase because selling them destroys the reason for having them. A judicial outcome cannot properly become a family gift. A public office cannot be inherited merely because the previous holder loved his child. A place obtained through bribery is not an unfortunate by-product of parental care; it is an injury to the institution and to those excluded.

Saying this clearly should not depend on whether the child is talented. A talented beneficiary can still receive a corrupt advantage. Competence may qualify her to compete for the role. It does not justify arranging the competition in her favour.

\section*{Helping without taking over}

Public support for children is sometimes described as an intrusion into a private responsibility. Yet societies depend on children becoming adults capable of sustaining the shared world. Schools, care, health, and time for parenting are not concerns limited to the household that happens to pay for them.

Support can also enlarge family freedom. Reliable childcare may allow a parent to work without abandoning care. A good local school can spare a family the need to move. A service that helps with a disability can make the difference between a workable home life and an exhausting sequence of emergencies.

But assistance can become supervision in ways that families reasonably resent. A service may require intrusive disclosure unrelated to the help requested. An official may assume that a family with little money is incompetent in every other respect. Conditions attached to support can turn ordinary parenting into a continuing demonstration of respectability.

The questions should be specific. What information is needed to provide the service? What evidence would justify intervention? Who can challenge a mistaken judgement? Is assistance designed to enable the family to function, or to make the family dependent on the approval of a particular office?

Universal services can avoid some humiliating distinctions, but universality does not guarantee quality or respect. Targeted help can reach particular needs, but targeting can miss people and consume effort in proving eligibility. The design has to be judged through the lives of its users, not only the intentions of its administrators.

An effective system should also be able to protect children from serious harm. Family privacy is important because relationships need space. It is not a licence to conceal abuse. The difficulty of distinguishing protection from overreach is real and cannot be settled by treating either parents or officials as naturally benevolent.

\section*{The adult who still needs a beginning}

Some people reach adulthood without the support described here. They cannot recover it by being told that families are valuable. They may need housing, education, advice, treatment, or a dependable person who helps them navigate a difficult period.

The response should not assume that the missed beginning fixes the rest of the life. People can learn, form relationships, acquire skills, and find forms of security later. Institutions can make those possibilities easier through second chances that are not reserved for people able to finance them privately.

There is nevertheless something irrecoverable in the loss. An adult can become an excellent reader without ever having had someone read to her as a child. The skill can be acquired; the remembered evening cannot. A public service that helps now should not be dismissed because it cannot return then.

This distinction matters because arguments about family and equality often promise too much from their preferred remedy. Money cannot reproduce every good of care. Care cannot compensate for every material deprivation. A fair examination cannot erase an unequal childhood. The limits of one remedy do not establish the uselessness of the others.

We are trying to provide more people with opportunities to live well, while recognising that lives do not begin in interchangeable conditions. Some differences should be reduced urgently. Some should be tolerated because removing them would damage relationships we have reason to value. Others should be stopped because they give one family power over the common rules.

The work consists in making those distinctions accurately enough to act. It will remain imperfect even in a much better society.

\section*{A gift that belongs to its recipient}

The parent reading aloud may hope that her child becomes accomplished. She may also simply enjoy the evening. The second reason matters. If every moment of childhood is evaluated by the advantage it produces, a defence of family life begins to resemble the competition it was supposed to place in perspective.

A child should have some experiences that need not become qualifications. A walk can remain a walk. A conversation can be interesting without improving an application. Time together need not be justified by its later economic return.

This does not make the social effects disappear. The conversation may develop confidence. The walk may teach something. It reminds us that those effects are not the whole reason for allowing people to share a life.

Eventually the child may value a different life. She may retain little of the parent's ambition and much of something the parent scarcely noticed giving. She may remember patience, a joke, a meal, or the assurance that a failure would not end the relationship.

Parents cannot reliably choose what survives them. The wish to do so can turn a gift into a continuing demand. Giving a person resources to live should include some acceptance that the life will become hers.

This is a quieter limit on parental power than a law or a tax. It asks the parent to endure being less necessary. That can be painful even when things have gone well. A child capable of leaving is often the evidence that years of care achieved something worthwhile.

\chapter{What We Have Earned}

The best-qualified applicant may deserve the job without deserving everything that made her the best-qualified applicant.

She studied, practised, and persisted. She may also have had excellent teachers, good health, parents who could help, and the good fortune to discover a talent that the labour market rewards. Recognising these conditions does not make her effort imaginary. Recognising her effort does not make the conditions disappear.

We often ask the word merit to decide both questions. It begins as a reason for assigning a particular task and ends as a judgement about the worth of a whole person. The successful applicant is entitled not merely to the role but to regard every later advantage as deserved. The unsuccessful applicant is expected to explain what is wrong with herself.

This enlargement is convenient for winners. It is also unnecessary for competent selection. We can want a surgeon who operates well without believing that her salary measures how much more valuable a human being she is than the person cleaning the ward.

Different kinds of work require different abilities. Some are scarce, difficult to acquire, or carry unusual responsibility. These facts can justify differences in reward. They do not produce a complete ranking of human worth. A society that forgets the distinction adds humiliation to inequality and then calls the result motivation.

\section*{What a wage tells us}

A wage is influenced by what employers need, what workers can offer, available alternatives, bargaining power, institutions, and the surrounding economy. It is not awarded by an impartial observer after comparing the moral effort of every working person.

Someone can perform demanding, useful work and receive little because many people are available to do it, because purchasers have limited money, or because the worker has few alternatives. Another can earn a great deal because he controls a scarce asset or occupies a position protected from competition. Effort alone does not explain the difference.

This does not mean that prices are meaningless or that an official could calculate the morally correct wage for every task. Prices convey information about scarcity and demand. They help people decide where to work, what to produce, and which alternatives to develop. Replacing them requires institutions with their own information problems and powers.

The more limited conclusion is that a market outcome needs more than its existence to become a moral justification. We can ask whether competition is open, whether participants face coercion or deception, whether costs are imposed on outsiders, and whether a rule protects a useful activity or a privileged position.

The person earning well can answer these questions without apologising for having a useful skill. The questions become harder when income depends on preventing others from developing or offering the same skill. A profession that uses unnecessary requirements to keep entrants out is not simply being rewarded for excellence. It is helping control the supply of competitors.

The argument should be made carefully in each case. Some requirements protect the public. Others protect incumbents. Calling every qualification a barrier can produce incompetence; calling every barrier a standard can preserve a comfortable monopoly. We need to know what the requirement tests and whether a less restrictive route could establish the same competence.

\section*{The fortune that was built honestly}

Suppose a woman builds a successful firm. She identifies a need, takes risks, employs people, and produces something customers value. She obeys the rules, pays what she owes, and accumulates substantial wealth. At the end of her life she wants to leave it to her children.

An argument against the transfer must take her claims seriously. She did not work only to satisfy a public plan. She had projects and attachments of her own. The prospect of helping her children may have been one reason she continued when the work was difficult.

The entitlement approach associated with Robert Nozick asks about the history of acquisition and transfer rather than judging a distribution solely by the pattern it produces. If holdings were justly acquired and transferred, a later inequality is not by itself proof of injustice.\footnote{Robert Nozick, \emph{Anarchy, State, and Utopia}, Basic Books, 1974, chapter 7, on entitlement, acquisition, transfer, and rectification.} This is a serious challenge to the view that any preferred distribution may simply be imposed whenever actual holdings diverge from it.

It does not spare us questions about the history. What counts as just acquisition? Were costs forced on people outside the transaction? Were earlier injustices involved? What institutions make ownership enforceable? And what rules about taxation and transfer were part of the system in which the fortune was created?

The founder's achievement also differs from her children's receipt. They may have worked in the firm or contributed in other ways. Or they may have done nothing to create it. Her right to direct a transfer and their claim to have earned it are not the same claim.

This is why the slogan that money has already been taxed does not settle inheritance policy. It identifies an earlier tax on an earlier event or holder. Whether a later transfer should be taxed depends on a further argument about the tax system, the purposes it serves, and its effects. The mere number of times money has figured in taxable transactions cannot decide every future transaction.

An inheritance tax can still be poorly designed, excessive, or easily avoided. Those objections deserve evidence. A serious defence of family property should engage with the actual rule rather than treat every proposed limit as a denial that the founder achieved anything.

\section*{What the tax is supposed to accomplish}

Inheritance taxation can pursue several purposes. It can raise revenue, reduce the concentration of wealth, limit advantages recipients did not create, or finance opportunities for people receiving little from their own families. These purposes overlap but are not identical.

A tax designed mainly to raise revenue may look different from one intended to restrain very large transfers. A rule that raises little because extensive exemptions protect the largest fortunes may still impose inconvenient costs on smaller estates. Its name does not tell us whether it serves its announced purpose.

The OECD's 2021 examination of inheritance taxation discusses both the case for the tax and the importance of design, including exemptions, avoidance, business assets, and liquidity. It treats these as practical questions rather than reasons either to abandon taxation entirely or to assume that any tax bearing the name is effective.\footnote{OECD, \emph{Inheritance Taxation in OECD Countries}, OECD Tax Policy Studies, 2021, doi:10.1787/e2879a7d-en. \href{https://www.oecd.org/en/publications/inheritance-taxation-in-oecd-countries_e2879a7d-en.html}{Report}.}

Liquidity is easy to illustrate. Suppose an inheritance consists chiefly of a working business, with little cash available. A substantial immediate payment could require borrowing, selling assets, selling part of the business, or changing its operation. The economic effect depends on the options. It cannot be inferred from the paper value alone.

Deferral or payment by instalments may address some such difficulties. Broad exemptions can address them too, but may create routes through which wealth is reorganised to avoid tax. The rule must therefore distinguish a real problem in paying from a convenient label attached to assets that do not need special treatment.

Gifts made before death also matter. So do valuation, trusts, residence, and enforcement. A high headline rate accompanied by easy avoidance can reward access to advice rather than reduce inherited advantage. A simpler, more enforceable rule with a lower rate may perform differently. The question requires examination of actual behaviour, not approval of the impressive number.

The same discipline applies to claims that a tax will destroy investment or save equality. People respond to rules in varied ways. Some work mainly for income, some for family, some for status, interest, or duty. A change can affect their choices without determining them uniformly. Forecasts about those effects should expose their assumptions.

\section*{The work that might not be done}

Any account of redistribution must consider production. The resources being distributed are not all lying in a storehouse unaffected by what people expect to happen next.

A person may decide whether to expand a firm, train for a difficult occupation, save, or retire partly in response to the reward she expects. A tax or regulation can change that decision. Ignoring the possibility does not make an egalitarian proposal more generous. It can leave less available for the people it intends to help.

But the incentive argument is frequently extended beyond its evidence. A wealthy person may describe every preferred concession as necessary for enterprise. A protected industry may call competition a threat to investment. A recipient of inherited assets may claim the productive virtue of the ancestor who created them while devoting his own energy to protecting an easy income.

Concentrated advantage can weaken incentives too. A young person who expects a closed profession to reject outsiders has less reason to prepare for it. A new firm confronting a market protected for incumbents may never enter. A secure heir may have little need to improve a business that political access shields from failure.

We should therefore examine incentives for more than the people already at the top. A system may encourage one family to preserve its wealth while discouraging many others from developing capacities that would enlarge the economy. The defence of incentives becomes selective when it treats only the first response as economically important.

A useful policy seeks to preserve rewards for creating value while limiting rewards for excluding competitors, obtaining exemptions, or shifting costs onto people unable to object. These activities can coexist within one firm or fortune. That makes judgement harder; it does not make the distinction unimportant.

\section*{What a common inheritance might provide}

Parents are not the only people from whom we inherit. Roads, public health, accumulated scientific knowledge, legal institutions, and useful standards were developed by people now dead or unknown to us. Private effort works within a world partly made by others.

This does not establish that the state owns every result of private effort. It does challenge the picture of a fortune produced entirely outside society and later subjected to an unrelated public demand. Some contribution to maintaining the shared conditions can be justified as part of the arrangement that makes private activity possible.

The revenue can be used well or badly. Invoking the common inheritance does not certify the budget. We should ask whether the institutions financed by it provide reliable services, widen access, and remain accountable. Waste does not become harmless because the money was collected for a fair purpose.

One possible use is to provide young adults with some capital of their own. An initial grant could help pay for training, equipment, a move, or a period between jobs. It would not equalise childhoods. It could reduce the dependence of people who receive little from family.

The details matter. A grant can be absorbed by higher prices if the relevant supply is constrained. Restrictions on use can prevent some waste while creating intrusive supervision. An unrestricted payment respects choice but permits choices others will dislike. Public provision of education or housing may serve some purposes better than cash, while cash can accommodate needs an official scheme misses.

These are design questions, not reasons to expect a single perfect instrument. The standard should be whether the arrangement materially improves people's capacity to make and revise a life without creating a worse dependence on administrators or private gatekeepers.

\section*{Losing without being ruined}

The importance of inherited wealth increases when one unsuccessful decision can have consequences from which a person cannot recover.

Two people leave secure jobs to pursue a new occupation. One can return to a family home and receive help. The other faces unpaid rent almost immediately. Even if their abilities are identical, the second may rationally avoid the change. A society that praises the first person's courage while ignoring the difference in possible losses has misunderstood the decision.

Reliable basic services and opportunities to retrain can make experimentation available to more people. Security can support effort by making a failed attempt survivable. It can also be badly designed in ways that discourage work or make additional earnings yield little practical gain. Again, the effect depends on the arrangement.

The aim is not a promise that every ambition will be financed indefinitely. People must eventually face constraints and make choices. The question is whether the consequences of a reasonable failure are proportionate, and whether a second chance depends almost entirely on family wealth.

A person who has lost a competition still needs somewhere to live, medical care, and some prospect of useful work. Providing these does not deny the winner's achievement. It refuses to make the result a complete judgement of whose life deserves support.

This is especially important when public rhetoric insists that everyone can succeed if they try. Some positions are scarce. Some abilities are unequally distributed. Illness, timing, family obligations, and luck affect outcomes. Encouragement becomes cruelty when it demands that every unsuccessful person interpret these conditions as proof of a defective will.

We can expect responsibility from people without claiming that they chose every obstacle they face. We can also recognise obstacles without assuming that no action remains possible. Practical assistance should help discover what remains possible rather than attach a permanent identity of failure or victimhood.

\section*{The attraction of a final account}

Both the successful and the unsuccessful may want a complete explanation that makes their position morally intelligible. The winner wants effort to account for everything. The loser may want circumstances to account for everything. Neither desire provides the evidence needed to explain a particular life.

We rarely know how much would have changed under another childhood, a different teacher, another illness, or a decision made a year earlier. This uncertainty does not stop us recognising achievement. It should stop us pretending that the recognition settles every claim the achiever can make on the rest of society.

It should also restrain the pleasure of exposing unearned advantage. A person who had help may still have worked extraordinarily hard. Showing that her parents paid the rent does not show that her research, business, or art required no talent. The point of mentioning the rent is to give a fuller account, not to erase what she did.

There is an ordinary form of honesty available here. A person can say that she accomplished something and that help, circumstance, and luck made the accomplishment possible. She need not reduce herself to a passive product of forces. She need not claim to have created herself either.

Public policy requires a similar restraint. It can reward skill, tax transfers, provide services, and limit the purchase of influence without announcing that the resulting distribution is perfectly deserved. Some inequalities will remain because removing them would damage valuable freedoms or cost more than the improvement could justify. Others will remain because reform has failed or power has resisted it. Those are different explanations and should not be collapsed into an endorsement of the result.

The best-qualified applicant should get the job when the qualification is relevant and the selection fair. Afterwards she should be able to live among people who did not get it without imagining that the appointment revealed their lesser worth. A position gives her work to do. It need not give her a philosophy of human rank.

\chapter{Being Good in Company}

Near the end of the Soviet Union, I visited the Lebedev Physical Institute in Moscow. I brought a bottle of red wine from the West. My colleagues produced another bottle, brought back from a conference in Georgia.

Alcohol was difficult to obtain during Gorbachev's anti-alcohol campaign. The bottles were modest gifts, but obtaining them had required more than walking into a nearby shop. A director thanked me and said something I have remembered ever since.

``It is good that you brought a bottle of red wine. It would have been better if you had brought your whole country with you.''

The remark was funny, friendly, and impossible to fulfil. It also made the limitations of my generosity apparent. I could supply a bottle. I could not supply the arrangements through which ordinary goods became routinely available where I lived.

The campaign itself had serious purposes. Alcohol abuse damaged health and family life, and restrictions could reduce that harm. Research has associated the campaign with substantial mortality changes. Restricting legal supply also affected access, revenue, and behaviour in ways a simple account of good intentions would miss.\footnote{Jay Bhattacharya, Christina Gathmann, and Grant Miller, ``The Gorbachev Anti-Alcohol Campaign and Russia's Mortality Crisis'', \emph{American Economic Journal: Applied Economics} 5, no. 2 (2013), pp. 232--260, doi:10.1257/app.5.2.232. The visit and conversation are the author's recollection from the original manuscript.} A gift of wine at that moment did not constitute a policy on any of these matters. It was a gift among colleagues.

Remembering the incident, I can make myself the fortunate visitor who understood the limits of charity. At the time I was simply someone who had brought wine. The director supplied the better sentence.

\section*{The people we become in our stories}

We tell stories about our lives partly to make them intelligible. We choose events that seem to explain what changed us, what we valued, and why we acted as we did. The selection need not be dishonest to be selective.

An act remembered as courage may also have involved anger. A sacrifice may have been freely chosen and later resented. A gift may have expressed affection and a wish to be liked. There is no reason to expect our motives to arrive one at a time.

The stories become misleading when the part we want to play determines which consequences can be admitted. Someone who thinks of herself as the generous member of a family may find it difficult to hear that her gifts create obligations recipients dislike. A defender of free expression may be outraged by criticism of his own speech. A person proud of facing unpleasant facts may repeatedly choose facts that make other people look foolish.

The role can survive these contradictions because each is given a separate explanation. This time the recipient was ungrateful. This criticism was malicious. This contempt was deserved. Any explanation may be right in a particular case. A pattern in which the cherished role never needs revision deserves a closer look.

The same occurs in organisations. A charity can think of itself as the voice of people it has stopped listening to. A political movement can treat the harm caused by its members as a temporary obstacle on a morally sound journey. A firm can make responsibility part of its brand and then regard scrutiny as damage to that responsibility.

The danger is not that people or organisations have identities. Some continuity of purpose is useful. The danger is that preserving the identity becomes more important than noticing whether the conduct still supports it.

\section*{Why the audience matters}

A gift often has witnesses. The witnesses can encourage more giving, remind others of a need, and help establish that promises are being kept. There is nothing inherently corrupt about visible generosity.

Suppose neighbours organise meals for someone recovering from an illness. Publishing a rota helps ensure that food arrives. Thanking those who volunteer recognises work and may encourage others. The public aspect contributes directly to the practical result.

It would be silly to dismiss the meals because some volunteers enjoy being appreciated. A person can want gratitude and still cook a useful dinner. We should be more interested in whether the meals arrive, meet the recipient's needs, and distribute the work reasonably than in proving that each cook has perfectly pure motives.

The arrangement can nevertheless change. Suppose praise is attached mainly to announcing participation. Several people volunteer for the first week, when everyone is interested. Later weeks receive little attention. The same two neighbours fill the gaps while the original volunteers continue to enjoy a reputation for having helped.

Now the public account is becoming detached from the work. Correcting it need not involve accusing every absentee of hypocrisy. The organiser can make the remaining tasks visible, ask for definite commitments, and stop presenting an initial offer as a completed contribution.

The problem is harder when the people controlling the publicity also benefit from the flattering account. An institution may prefer the photograph of the launch to the record of what happened three months later. A donor may be offered a ceremony because ceremonies are easier to organise than a careful evaluation of results.

Public recognition should therefore be connected to work that can be checked. Otherwise the people best at appearing to contribute may acquire authority over those doing the less visible part.

\section*{The unpleasant person who helps}

We can make the opposite mistake by allowing our dislike of a person to erase the value of what he does.

A donor may be pompous, politically objectionable, and eager to have his name displayed. The money may still pay for something needed. Accepting that fact does not oblige an institution to accept every condition he attaches or to pretend that his other conduct is admirable.

Likewise, a helpful person can behave badly elsewhere. A public contribution does not provide a general exemption from scrutiny. We should be capable of thanking someone for one thing and holding him responsible for another.

This is harder than it sounds because reputation tends to become a single description. Once a person is good, accusations seem unfair. Once a person is bad, credit seems disloyal. The urge to keep the description coherent can distort the assessment in either direction.

Institutions need procedures that resist this simplification. A complaint about mistreatment should not fail because the accused raises money successfully. A contract should not be granted because its supplier supports a popular cause. A donation should not purchase access to decisions unrelated to the gift.

Nor should misconduct automatically make every association with the person contaminating. The question is what continuing the association does, what obligations exist, and which standards should apply. Public life becomes needlessly destructive when people are expected to erase every complicated connection in order to demonstrate that they have understood a wrong.

One benefit of attending to particular acts is that we need not decide the final moral character of everyone involved. We can accept the useful meal, refuse an improper condition, investigate a complaint, and leave the complete account of the person unfinished.

\section*{Cooperation that remembers}

The importance of reputation becomes clear when people expect to deal with each other again. A shopkeeper has reason to treat a customer fairly if the customer can return, tell others, or take business elsewhere. A neighbour has reason to keep a promise if later cooperation will depend on being trusted.

Robert Axelrod's well-known computer tournaments explored strategies in a repeated Prisoner's Dilemma, a simplified game in which each player can cooperate or defect. The strategy called TIT FOR TAT began by cooperating and then copied the other player's previous action. Its strong performance in the tournaments helped draw attention to the combination of initial cooperation, a response to exploitation, and a return to cooperation when the other player resumed it.\footnote{Robert Axelrod, \emph{The Evolution of Cooperation}, Basic Books, 1984; Robert Axelrod and William D. Hamilton, ``The Evolution of Cooperation'', \emph{Science} 211 (1981), pp. 1390--1396, doi:10.1126/science.7466396. The argument concerns specified repeated-game settings, not a universal ranking of moral conduct.}

The result is often reduced to the cheerful claim that being nice wins. The rule itself says something more demanding. It does not keep cooperating regardless of what the other player does. It also does not make a single defection the justification for permanent hostility.

The model is deliberately limited. Real people can misunderstand actions, make mistakes, have unequal power, and leave relationships. Strictly copying a perceived wrong can prolong an accidental quarrel. A strategy successful in one population of opponents is not a universal commandment for families or states.

Its usefulness lies in separating parts of cooperation that moral language can blur. An initial willingness to help is one part. A response to repeated exploitation is another. Making renewed cooperation possible is another. A person can be deficient in any one of them.

Consider the neighbour who borrows tools. Lending once may be an ordinary courtesy. Lending again after several tools have been returned damaged may be imprudent. Refusing all future contact after one promptly acknowledged accident may destroy a relationship unnecessarily. The appropriate response depends on the pattern, the explanation, and whether the borrower changes his conduct.

A rule can make this clearer. The next loan may require agreement about return and repair. If the neighbour keeps the agreement, trust can recover through conduct. The lender need not decide whether the borrower is fundamentally a good person before offering another chance.

\section*{The good person who cannot refuse}

Some people find refusal especially difficult because their position within a group depends on being helpful. They organise, cover shifts, lend money, and listen. Others come to regard their availability as a feature of the world rather than an effort someone is making.

At first the person may enjoy being needed. Later she finds that even a modest refusal requires an explanation. A friend who asks whether she can help has already planned around the answer. Her withdrawal is treated as a new difficulty she has created.

There can be genuine needs on the other side. The friend may have no convenient alternative. This does not make the helper's time inexhaustible. If the group depends on a contribution, it should recognise and distribute the work rather than praise the person into continuing indefinitely.

The helper has a responsibility too. If she repeatedly agrees while privately resenting the request, others may not know the arrangement is failing. They may be selfish, or they may be taking her answer seriously. Her wish to preserve a generous image can contribute to the misunderstanding.

This is not a reason to blame an exhausted person for everything that has happened. It is a reason to make refusal possible before exhaustion supplies it in a harsher form. A clear limit can protect a relationship that an eventual explosion would damage.

The limit need not be accompanied by a new identity as someone who has learned to put herself first. She may still care greatly about the people asking. She may simply be unable or unwilling to do this particular thing. Turning the refusal into a complete philosophy creates another role to maintain.

\section*{When the group rewards the wrong sacrifice}

A group can also make sacrifice competitive. Members demonstrate commitment by working longer, accepting worse conditions, giving more, or tolerating treatment they would condemn elsewhere. The person who proposes a limit appears less devoted than the person who agrees to one more demand.

This can happen in admirable causes because the cause supplies a reason to overlook the immediate cost. A volunteer organisation needs to help vulnerable people. A research team needs to finish important work. A campaign faces a serious opponent. There is always a reason that this week is exceptional.

The burden may fall unevenly. People with savings can accept unpaid work longer than people supporting a household. People without caring responsibilities can attend more late meetings. Those differences may then be interpreted as differences in commitment.

The group loses capable members and learns the wrong lesson from their departure. They did not care enough, it says. In fact, it may have designed participation around a narrow kind of life and called that life morally superior.

The practical correction begins by asking what work is necessary and what merely demonstrates devotion. A meeting that produces nothing need not become a test of belonging. A task requiring skill should not be allocated according to who can afford to perform it free. A person who needs rest should not have to discredit the cause in order to defend the need.

The cause itself may benefit from these corrections. Exhaustion, resentment, and a narrowing membership can damage work that depends on judgement and continuity. The appearance of wholehearted commitment is not a reliable measure of how well an organisation can keep its promises.

\section*{What the bottle could do}

The bottle in Moscow could be shared. It could express friendship. It could make an evening more pleasant. Those were worthwhile possibilities even though it could not repair an economy.

The director's joke enlarges the thought. The availability of the gift depended on arrangements beyond the giver: production, distribution, payment, trust, and ordinary routines that made replenishment possible. Bringing another bottle would help another evening. It would not automatically reproduce those arrangements.

This is not a reason to despise small gifts. Most affection is expressed through small acts that do not solve large problems. It is a reason to know what kind of act we are performing. The mistake begins when the good feeling attached to a gift is treated as evidence that the same action, expanded indefinitely, supplies a workable policy.

There may be ways to provide more, help build capacity, or change rules that produce scarcity. Each requires further knowledge and decisions. The recipient's gratitude cannot answer those questions on the donor's behalf.

It also need not become the donor's lasting possession. A person helped today may disagree with us tomorrow. A colleague who accepts wine does not accept our account of his country. A friend who receives care does not owe us a flattering role in every future version of the story.

We can give something useful and allow the relationship to remain a relationship between separate people. That involves less control than generosity sometimes seeks.

Years later, the remembered remark remains sharper than any lesson I might have offered at the table. I cannot make the visit into proof that I understood the society around me. I can remember that people who lived there understood things a visitor could not bring in a bottle. Their joke was part of what they gave back.

\chapter{The Power to Call Yourself Good}

Imagine a donor visiting a workshop her money helped establish. One of its former trainees now runs a small business. He employs two people, has little time for meetings, and declines an invitation to appear in the donor's next film. He is grateful for the help. He would rather spend the afternoon working.

She is disappointed. His independence was among the project's stated aims. Yet the independent person is less available than the grateful beneficiary. He no longer wants to explain his past difficulties to strangers, and he has opinions about the programme that would complicate the film.

The donor may accept this gracefully. Or she may begin to find him arrogant. Success has gone to his head, she says. He seems to have forgotten where he came from.

Something worth examining has happened if the same independence changes from an achievement into a character defect when it inconveniences the benefactor. The gift may have contained a condition that appeared in no agreement: you may improve your life, but you must continue to make my place in it visible.

This is one way that helping can acquire an interest in dependence. It need not begin in fraud. A useful act, warmly intended, can establish a relationship whose giver later struggles to relinquish.

\section*{Nietzsche's unwelcome question}

Nietzsche refuses to accept moral praise at its own valuation. If people call themselves humble, compassionate, or selfless, he asks what form of power becomes available through those descriptions. The answer may be less flattering than the vocabulary.

In the first essay of \emph{On the Genealogy of Morality}, he presents \emph{ressentiment} as a force that creates values through a hostile reversal. People unable to prevail directly make the opponent's strength evil and their own condition good. The victory takes place in the interpretation of the contest. His polemical account of Christianity links moral authority to this transformation; the later essays examine guilt and the ascetic ideal.\footnote{Friedrich Nietzsche, \href{https://www.gutenberg.org/files/52319/52319-h/52319-h.htm}{\emph{On the Genealogy of Morality}}, 1887, especially first essay, sections 10--14, and the second and third essays. The linked older translation uses the title \emph{The Genealogy of Morals}. Nietzsche's historical and psychological claims are presented as arguments to examine, not as established diagnoses of particular contemporary groups.}

The argument is unsettling because it questions the apparent renunciation of power. Someone unable to command through force may command through guilt. Someone who cannot win a competition may acquire the authority to declare winning discreditable. The language of humility can place its speaker above the people it judges proud.

This is not a reliable history of every value or a diagnosis to apply to everyone who asks for help. A complaint can be justified. Weak people can identify actual wrongs. The interesting challenge is to discover whether the moral description explains a practice or protects its beneficiary from examination.

An administrator who speaks for the vulnerable may obtain a position, a budget, and a right to question other people's conduct. A critic should ask whether the people represented can question hers. If they cannot dismiss their representative, dispute her account, or say that they no longer want her service, the representation has begun to resemble ownership.

The institution's language may make this difficult to notice. Its clients are empowered; its programmes are emancipatory; its leadership listens. Meanwhile, a client who rejects its preferred account is described as confused, manipulated, or not yet ready. The organisation can hear the people it serves only when they say what it expected.

There is no need to infer secret malice. A sincerely convinced helper can be especially obstinate because disagreement appears to threaten the recipient's welfare. The wish to control is experienced as concern. The resulting control is no less real for that.

\section*{The advantage of being injured}

A person who has suffered a wrong deserves an accurate account of it and, where possible, a remedy. Suffering does not confer a general competence on every subject or make every later act just.

These distinctions are often difficult to maintain in a campaign. The injured person becomes its representative. His account attracts attention, his presence authenticates the cause, and his disagreement with a policy is treated as decisive because he has experienced the harm. Experience supplies knowledge that outsiders can lack. It does not supply every kind of knowledge needed to decide what follows.

Someone who has survived a failing service may know exactly where it humiliated him. He may know less about the effects of a proposed replacement on other users. A bereaved parent can describe a loss with authority no official possesses. That authority does not automatically settle the design of a national rule.

Saying so can sound cold because public sympathy has joined two things: acknowledgment of the experience and acceptance of the proposed conclusion. A person who doubts the conclusion appears to doubt the suffering. The campaign can then avoid a practical objection by returning to an experience nobody disputed.

There is a further danger for the person placed in this role. Recovery, ambivalence, or an inconvenient opinion can make him less useful. He may feel that he has to remain publicly injured in the approved way. A cause organised around giving him a voice can leave him repeating a version of himself he has outgrown.

An institution should make room for the possibility that those it helped will cease to need its account. The former trainee can become an employer with ordinary faults, ambitions, and disagreements. He should not have to preserve a history of helplessness as his main qualification for being heard.

None of this cancels a valid claim for redress. It limits the authority that can be extracted from it. An injury should not become a passport allowing its holder, or more often its interpreter, to cross every boundary in the lives of others.

\section*{Resentment is not an explanation by itself}

The language of resentment can be used just as dishonestly as the language of injury. A comfortable person hears a complaint about unfair treatment and diagnoses envy. Having named a discreditable emotion, he considers the complaint answered.

But envy and injustice can coexist. A worker may resent a manager's salary and also have been denied wages owed. A rejected applicant may dislike the successful candidate and also have evidence that the competition was rigged. Discovering the emotion does not adjudicate the claim.

The same is true of the historical origin of a value. An idea can emerge through an ugly struggle and acquire defensible reasons later. Conversely, a principle with a noble origin can be used cruelly. We should examine the reasons and consequences rather than imagine that a genealogy supplies an automatic verdict.

This places a limit on the use of Nietzsche. His suspicion can expose a neglected possibility. It cannot substitute for establishing that the possibility occurred. Calling compassion a will to power whenever it restrains someone powerful is merely to distribute moral suspicion in favour of the strong.

Nor should strength be confused with whatever enables a person to win. A bully can prevail over someone who has shown more courage, restraint, or practical ability. An inherited monopoly can defeat a capable competitor. A prisoner may display a strength of character absent from the official who controls the door. The visible distribution of command does not tell us everything worth admiring.

A critic who borrows Nietzsche's contempt without his unsettling questions can end by applauding the nearest large fist. This is no escape from moral sentimentalism. It is sentimentalism about strength, complete with a flattering picture of the kind of person the critic imagines himself to be.

\section*{Compassion that notices the person}

The injured stranger in the parable of the Good Samaritan has more immediate needs than a satisfactory theory of morality. He has been attacked and left beside the road. The Samaritan attends to his wounds, transports him to an inn, pays for care, and undertakes to meet further expense.\footnote{Luke 10:25--37, the parable of the Good Samaritan.}

The story challenges the refusal to recognise a claim because the person in need belongs to the wrong category. Its practical details also deserve attention. The helper does things. He makes arrangements. He puts resources behind a promise.

A reader can draw demanding obligations from that example. What the reader cannot honestly extract is a complete administrative programme already supplied by the parable. It does not decide how a hospital should be funded, which public authority should act, or how to resolve several simultaneous claims on limited capacity. Those questions require additional thought.

Equally, its use of a particular person's resources does not prove that organised public assistance is illegitimate. Pooling resources can make help reliable where waiting for a passing benefactor would leave people abandoned. The question is whether the organisation preserves attention to the actual need and accepts responsibility for delivery.

Christian moral language also contains its own criticism of performance. The warning in Matthew against giving in order to be seen by others directly challenges a use of charity for public elevation.\footnote{Matthew 6:1--4, on almsgiving and the desire for an audience.} Nietzsche's picture does not exhaust everything within the tradition he attacks.

That observation is not a request to exempt churches from his questions. An institution can preserve a severe warning in its sacred text while ignoring it in its fundraising, hierarchy, and treatment of dissent. Indeed, the contrast can make the failure sharper. The criticism is available inside the institution's own declared commitments.

Compassion deserves examination by what it permits the other person to become. Does it relieve a need, restore options, and accept an eventual departure? Or does it require a continuing demonstration of inferiority so that the helper can remain necessary?

There will be people whose need does not end. Severe disability, illness, or age can require enduring assistance. Their continued dependence is not evidence that care has failed. Independence is not a moral entrance requirement for receiving help. The relevant freedom may be the ability to express a preference, refuse an indignity, or retain a relationship that is not controlled by the provider.

A society that recognises only the successful graduate of assistance will neglect people who cannot perform that success. Their lives are not defective versions of the donor's preferred ending.

\section*{When excellence is treated as an accusation}

There is a related use of moral language against people who do something unusually well. Their achievement changes the comparison. Colleagues who were comfortable with one another's performance may now feel exposed.

Imagine a workplace where a new employee solves recurring problems that others have learned to tolerate. She may be tactless about it. She may need to understand responsibilities her improvements overlook. Those are matters to discuss. Instead, she is told that she is making everyone else look bad.

The objection is revealing. It treats the appearance of the group as a more urgent concern than the work. If her achievement is real, others might learn from it. That would require admitting that something had been done poorly. It is easier to make excessive competence a breach of manners.

Concern for cooperation can then become a means of enforcing mediocrity. The talented person is called selfish for refusing to conceal what she can do. The established group calls its own refusal to improve solidarity.

The reverse abuse is familiar too. An employee may invoke excellence to excuse humiliation, refuse routine duties, or demand that everyone accommodate his preferences. A useful contribution can be real and still fail to justify those demands. Protecting excellence should mean allowing work to be done and recognised, not creating a caste whose behaviour is beyond review.

What matters is keeping the achievement separate from the status claims built around it. The group should not suppress a better result to avoid discomfort. The achiever should not convert a better result into sovereignty over colleagues.

There is room to admire someone without submitting to him. There is room to protect equal standing without pretending that every performance is equal. Public life becomes needlessly dishonest when it offers only worship or levelling as responses to distinction.

\section*{A sacrifice may serve the person watching}

Visible self-denial can establish authority even when it helps little. The exhausted campaigner, the conspicuously austere official, or the donor who announces how much has been given up asks the audience to infer commitment from cost.

Cost can be evidence of commitment. It is not evidence that the commitment is well directed. Someone can spend years pursuing a mistake. The investment may make correction more difficult, because acknowledging error would remove the public meaning of the sacrifice.

Others can become trapped by the comparison. A colleague who refuses a pointless late meeting appears less devoted than the person who stays. A parent who protects time with her children seems less serious than the worker who is always available. The institution has turned an expenditure of time into a moral ranking without first establishing what the expenditure accomplishes.

There is also a temptation to make suffering authenticate a proposal. If a measure is painful, it must be serious. If someone asks whether a less painful method would work, the question sounds like weakness. The inconvenience has become part of the desired display.

A sensible aim is usually to obtain the worthwhile result with less unnecessary suffering. Sacrifice may be unavoidable, but reducing its cost is ordinarily an improvement. An institution that becomes suspicious when the same good can be achieved more easily should ask whether the good was its only purpose.

This question applies to austerity as readily as to generosity. Cutting a service can demonstrate hardness while saving little. Demanding more unpaid effort can demonstrate dedication while ruining the people needed to continue the work. Neither hardship nor kindness proves its usefulness by its appearance.

\section*{The benefactor without a final word}

Return to the donor and the former trainee. Their disagreement need not end the relationship. She can value the workshop's achievement and accept that it has produced a person with whom she cannot arrange every interaction. He can acknowledge the help without agreeing to become its permanent advertisement.

If the project requires evaluation, he may have relevant obligations under an agreement. Those should be stated. Ordinary accountability differs from an indefinite emotional debt. The donor is entitled to examine whether money was used as promised. She is not entitled to decide which identity the recipient must inhabit afterwards.

Perhaps he really has become vain. People who succeed sometimes do. The appropriate response is the ordinary one available between adults: disagree with the vanity, limit a request, or decline further association. It need not involve reclaiming moral ownership of his life.

There is a small melancholy in becoming unnecessary to someone we helped. A familiar role ends. Gratitude may remain, but it no longer organises the other person's day. A gift that achieves its purpose can therefore deprive its giver of something she enjoyed.

Allowing that loss is among the less theatrical forms of generosity. There may be no film of it. The former trainee goes back to his work, and the donor has to find something else to do with the afternoon.

\chapter{The Consequences of a Good Intention}

Imagine several farms drawing water from the same reservoir during a dry summer. Each farmer can save part of a crop by taking more water now. If they all do so, the reservoir may be exhausted before the summer ends.

One farmer decides to show restraint. She takes less than she is allowed. Another uses the water left available. Her sacrifice has reduced her own harvest without preserving the reserve she intended to protect.

It would be wrong to conclude that restraint is useless. It would also be wrong to praise her decision as though its purpose had been achieved. What she intended, what she did, and what happened are three different parts of the account.

She might have influenced the others by setting an example. She might have begun a negotiation, exposed a defective rule, or preserved some water despite their behaviour. Those possibilities need examination. The agreeable fact that she meant well cannot decide which occurred.

Public morality often stops at the point where this examination should begin. The actor has made a sacrifice, defended a principle, or declared a limit. To ask what followed sounds like an attempt to deny the virtue of the act. Yet the crop, the reservoir, and the other users remain outside the declaration, responding to what people actually do.

\section*{The actor cannot be the only subject}

Max Weber's distinction between an ethic of conviction and an ethic of responsibility concerns this difficulty. In \emph{Politics as a Vocation}, he contrasts fidelity to a guiding commitment with answerability for foreseeable consequences. He does not simply equate conviction with irresponsibility or responsibility with having no principles. Towards the end of the lecture he treats the two as complementary demands on the person capable of political work.\footnote{Max Weber, \href{https://de.wikisource.org/wiki/Politik_als_Beruf}{\emph{Politik als Beruf}}, 1919, especially pp. 56--57 and 64--65 of the original edition. The distinction is commonly translated as the ethics of conviction and responsibility.}

The distinction is useful when an official defends a measure entirely by naming what it expresses. We stood for compassion. We refused complicity. We put safety first. Such statements tell us something about the account the official wishes to give of the decision. They have not yet told us what the decision did to the people affected.

If those people have become less safe, less able to obtain help, or more dependent on an abusive intermediary, the result deserves more than a repetition of the intention. The official cannot remain the main beneficiary of a policy whose practical effects have disappointed everyone else.

Responsibility also concerns what could reasonably have been foreseen. Nobody knows every consequence in advance. An unforeseeable event can defeat a competent plan. Judging solely by the eventual result would confuse bad luck with negligence and good luck with wisdom.

The relevant inquiry reconstructs the choice: what information was available, which warnings were credible, what alternatives existed, and what the decision-maker had reason to expect. It also asks what happened after the consequences became clearer. An initially defensible choice can become indefensible when maintained against accumulating evidence.

This is a more demanding standard than either praising intentions or condemning every failure. It requires us to understand a decision before assigning responsibility for it. That work is often less satisfying than finding a villain, including when a villain is eventually found.

\section*{Shared water is not an endless gift}

The farmers face a particular kind of problem. Water taken by one user is unavailable to the others, and keeping every user from taking too much requires an arrangement that can actually be maintained. Each person's incentive can conflict with the preservation of the shared resource.

This differs from sharing an idea. If one farmer learns a useful way to reduce leakage and teaches it to the others, his own knowledge need not diminish. It also differs from ordinary exchange, in which two parties may both obtain something they value more than what they surrender. The word sharing covers very different practical situations.

It matters which one we are discussing. An appeal appropriate to spreading knowledge may fail when the resource is physically depleted through use. A defence of mutually beneficial exchange may fail when part of the cost is imposed on people who never entered the transaction.

The reservoir does not become inexhaustible because the users have a common identity. Nor does it become impossible to govern because each user has an incentive to take water. Between those conclusions lies a field of institutional possibilities.

Elinor Ostrom's work challenged the assumption that shared resources must be governed either by privatisation or by a single distant authority. She studied arrangements in which users developed and maintained rules fitted to their circumstances. Defined rights, monitoring, sanctions, ways to resolve disputes, and participation in making rules could matter to their durability. Her argument was not that local communities always succeed or that one design fits every resource.\footnote{Elinor Ostrom, \emph{Governing the Commons}, 1990, and \href{https://www.nobelprize.org/prizes/economic-sciences/2009/ostrom/lecture/}{``Beyond Markets and States: Polycentric Governance of Complex Economic Systems''}, Nobel Prize lecture, 8 December 2009. The reservoir and farms in this chapter are invented examples, not a retelling of one of her case studies.}

In the imagined case, the farmers might agree on allocations that change with the reservoir level. They would need a way to measure withdrawals, detect breaches, and deal with genuine emergencies. A rule without monitoring could reward the person willing to ignore it. Monitoring without a fair procedure could give a local official opportunities to punish enemies.

Their knowledge would matter. They may understand the fields, crops, and seasonal demands better than someone arriving from the capital. Their unequal power would matter too. The largest farm might dominate the meeting, while tenants or downstream users have little voice. Calling an arrangement communal does not establish that everyone affected helped shape it.

Successful cooperation therefore needs more than a pleasing view of the community. It requires people to settle details that are easy to describe as unneighbourly until the water runs out.

\section*{What happens after someone breaks the rule?}

Suppose the farmers reach an agreement and one exceeds his allocation. The others can excuse him, penalise him, or reconsider the agreement. Each response can be justified under some circumstances. None is sensible without knowing what occurred.

A meter may have failed. An emergency may have arisen. The farmer may have deliberately taken extra water because the penalty is lower than the expected value of the saved crop. Treating every breach as an accident invites exploitation. Treating every accident as defiance can destroy cooperation that was working.

A graduated response gives the group more than two choices. It can establish the facts, require correction, and increase consequences if violations continue. The person sanctioned also needs a way to challenge an error. Otherwise enforcement can become the private power of whoever controls the records.

The humane part of the arrangement includes these limits on punishment. It also includes protecting people who complied. If the group repeatedly forgives a deliberate excess while others lose crops through restraint, its mercy is being financed by the obedient.

Eventually the compliant farmers may stop complying. Their change need not show that they were secretly selfish all along. A rule that applies only to those willing to be bound by it has ceased to offer them the cooperation they accepted.

This is one reason a limit can protect generosity. People may be more willing to share when they can see that persistent exploitation will be checked. Asking for unlimited patience can destroy the ordinary patience on which the arrangement depended.

There is no need to transform the violator into a permanent enemy. If he repairs the breach and observes the agreement, renewed cooperation may serve everyone. A rule that cannot allow return can make a single conflict more expensive than it needs to be.

\section*{The effects beyond the gesture}

The first farmer's unilateral restraint might still have value. It could make a risk visible or demonstrate that using less water is feasible. A public commitment can help others coordinate if they were each waiting for assurance that someone else would begin.

But an example needs people able and willing to follow it. If the other users simply take the available water, repeating the same sacrifice more intensely does not solve the problem. The farmer needs a different means: an agreement, enforceable limits, better information, or an improvement that reduces the demand without destroying the crop.

This distinction applies to consumer campaigns. Refusing a product can reduce demand, signal disapproval, or support an alternative. It can also shift production to another buyer or leave the most vulnerable workers bearing a cost that the owners avoid. The result depends on the market, available substitutes, bargaining power, and the campaign's scale.

A campaign should therefore explain the path from the requested act to the intended improvement. Who is expected to change behaviour, and why? What prevents a substitute arrangement from reproducing the harm elsewhere? What will be observed to discover whether the effort is working?

These questions do not require every participant to become an economist. They do require the organisation soliciting participation to offer more than a moral portrait of the conscientious purchaser. Buying or refusing to buy is an action in a system that other people can respond to. Their responses belong in the argument.

Symbolic action can be worthwhile even when its immediate material effect is small. It can express solidarity, refuse participation, or preserve a public objection. The claim should then be stated at that scale. A symbolic refusal does not become a successful remedy merely because describing it more modestly would make recruitment harder.

\section*{There are limits to calculation too}

Once consequences receive attention, another temptation appears. A clever official can defend almost any action by predicting a sufficiently large benefit. The harm is immediate and specific; the promised advantage is distant and difficult to test.

An institution may ask a few people to endure severe burdens for a future it describes with confidence. Those people may have no way to challenge the assumptions, and the officials may leave before the promised result can be assessed. Calling the calculation responsible gives it a respectability its evidence has not earned.

Rights and procedural limits constrain this power. They give people claims that cannot simply disappear inside an official's estimate of aggregate advantage. The constraints also reflect what we have reason to fear from institutions allowed to exempt themselves whenever they predict a good result.

This does not mean that consequences become irrelevant whenever a right is invoked. Rights can conflict, their scope can be disputed, and their enforcement has effects. It means that a projected benefit must be considered alongside the kind of authority required to obtain it and the protections being withdrawn.

In the reservoir case, an allocation that maximises total crop value might leave a small farmer with no viable livelihood. The group could choose to preserve more farms at the cost of a lower immediate total. It could provide compensation or alter planting before the next season. What it cannot do honestly is treat the largest monetary harvest as though it were the only possible meaning of a good outcome.

Values enter the choice. Counting does not remove them. It helps show what they cost under particular conditions.

The person who claims to have abandoned principles in favour of results has usually retained principles about which results count. Leaving those principles unstated makes them harder to examine; it does not make the decision free of them.

\section*{The time in which an obligation must be met}

A promise to enlarge capacity can be sound and still fail the person who needs help before the enlargement is complete.

The farmers might build additional storage, change crops, repair channels, or reduce losses. These measures could improve later summers. They do not necessarily supply water this afternoon. An argument about long-term abundance cannot decide a present allocation until the means connecting the two have been established.

The same applies to trained staff, homes, transport, and care. More can often be provided. People must be trained, buildings completed, equipment maintained, and routines made reliable. These processes take time even when the money has been approved and everyone intends to cooperate.

Meanwhile, someone allocates what exists. If political leaders refuse to acknowledge that task, it moves to the point of delivery. A clerk, teacher, or nurse has to say no under conditions the announcement ignored. The leader can then condemn the refusal as though it had not been made necessary by the gap in the promise.

This is an especially convenient arrangement for public virtue. The promise has a named author. The refusal belongs to an anonymous system. The person at the desk is expected to reconcile them with patience and a smile.

An honest commitment includes an account of the transition. Who receives the service first? What happens to those who wait? Which temporary losses are accepted? Who can change the plan if the new capacity arrives late?

The answers may expose an unattractive choice. Concealing the choice does not improve it. It merely ensures that the people most affected encounter it without the explanation those responsible owed them.

\section*{Responsibility after a bad result}

Suppose the farmers agree on a reasonable allocation and the summer turns out far drier than expected. Crops are lost despite compliance. The agreement has failed to achieve everything its participants hoped.

The later inquiry should not pretend that everyone knew the outcome in advance. It should examine whether the forecast was used appropriately, whether uncertainty was taken seriously, and whether action changed when the situation worsened. A poor harvest is not by itself proof that the original allocation was foolish.

Equally, the existence of uncertainty cannot protect a decision from all criticism. If the plan depended on a single favourable forecast while credible warnings were dismissed, the uncertainty was mishandled. If measurements showed the reserve falling faster than expected and no one had authority to respond, the arrangement contained a practical failure.

Accepting responsibility can involve changing the rule, compensating an avoidable injury, admitting a misleading promise, or leaving a position. It should not be reduced to appearing solemn at a press conference. The people who lost something need to know what acknowledgment changes.

There may still be a remainder that no admission resolves. A small farm can close after generations in one family. Its owners may understand why the allocation was reasonable and continue to grieve what followed. Their grief does not have to prove that someone acted wrongly in order to deserve recognition.

Moral argument often finds this unsatisfactory. It wants either a guilty actor or a vindicated outcome. Sometimes there is neither. People made a defensible choice within bad conditions, and somebody suffered a real loss.

\section*{What the good intention must endure}

The farmer who first restrained her use may learn that her act helped less than she hoped. That discovery need not make her abandon concern for the reservoir. It may give the concern a more effective form.

She can work on an agreement rather than demand that everyone admire her example. She can accept monitoring for her own farm as well as for the neighbour's. She can support a fair penalty when it falls on an ally. She can acknowledge an improvement proposed by someone whose earlier conduct she resented.

These acts offer less immediate satisfaction than a conspicuous sacrifice. They require cooperation with people who may never share her account of the problem. They also give her intention a chance to become something more durable than an entry in her own story.

The aim remains a good one. The means become answerable to the world in which the aim is pursued. This is where the next chapter begins: with someone who wishes to be good, a shop from which she hopes to provide help, and a growing collection of people who have plans for both.

\chapter{The Shop That Could Not Say No}

Shen Te buys a tobacco shop. She has been poor, and now she has a little money. She hopes to support herself without continuing to sell her body, and to help some of the people who are still as poor as she was. It is a modest hope. She does not propose to save the world. She would like to keep a shop and be kind.

Bertolt Brecht gives this hope very little time to settle. People arrive who need food, cigarettes, shelter, payment. Some have helped Shen Te before. Some have nowhere else to go. Some simply recognise a person who will find it difficult to turn them away. The shop fills. Its stock diminishes. The people who depend upon its remaining open behave as though keeping it open were somebody else's problem.

We need not decide that all these people are wicked to see the difficulty. A hungry person has a reason to ask for food. Ten hungry people have ten reasons. The reasons do not produce another sack of rice.

The gods who enabled Shen Te to buy the shop have been looking for a good person. Their search has not gone well. People who could accommodate them find excuses. Shen Te, who can barely accommodate herself, offers them a bed. They reward her, and the reward becomes a test that their command to be good has not prepared her to pass. Money gives her something to distribute. It does not tell her when to stop.

\section*{The cousin arrives}

Shen Te's solution is theatrical and painfully plausible. She invents a cousin, Shui Ta, dresses as a man, and returns with a different manner. The cousin can say what the good woman cannot. He tells people to leave. He bargains. He challenges demands instead of feeling obliged to satisfy them merely because they have been made.

The disguise changes no physical resource. There is still the same shop, the same stock, the same rent. What changes is the expectation that every claimant will be received with agreement. A person who was inexhaustibly available has acquired someone authorised to disappoint people on her behalf.

We recognise this arrangement outside the theatre. A sympathetic professional refers the awkward request to an administrator. A generous parent tells the child to ask the stricter parent. A charity's public appeal promises compassion while its staff apply eligibility rules. Sometimes the division of labour is sensible. Sometimes it lets one person enjoy being loved while another absorbs the anger required to keep their common undertaking solvent.

Shen Te must play both parts. She cannot afford a separate employee to protect her against her promises.

At first, we may welcome Shui Ta. Someone has finally noticed that the shop will close if its contents are treated as an inexhaustible common supply. People who praise Shen Te's generosity are not necessarily prepared to replenish it. Their indignation at the cousin can contain both genuine distress and annoyance that an easy source of goods has become difficult.

But Brecht does not leave the cousin at the door, performing the useful service of refusing the hundredth unreasonable request. Shen Te falls in love with Yang Sun, an unemployed pilot. Having saved him from suicide, she hopes to help him return to flying. He is willing to use her devotion. Her willingness to see what she hopes for, rather than what his conduct suggests, threatens the little security she has acquired.

Pregnancy makes the problem still more urgent. She now imagines a child whose survival depends upon her. The cousin appears more often. The shop becomes a larger tobacco business. People who once lived upon Shen Te's gifts find themselves working under Shui Ta's authority. The business survives and grows, while the means of its survival become increasingly harsh.

The woman who wanted to shelter people has constructed an enterprise that exploits them. This is not an incidental excess that can be removed from the story to leave a cheerful lesson about personal boundaries. It is part of the story's accusation.

When Shen Te remains absent, Shui Ta is suspected of having done away with her. The gods return as judges. She reveals the disguise and explains that she has been both people. The gods are relieved to discover that their good person still exists. They do not repair the conditions under which she must live. They leave her with their demand, and the play leaves the audience with the unresolved difficulty.\footnote{Bertolt Brecht, \emph{Der gute Mensch von Sezuan}, written with the collaboration of Ruth Berlau and Margarete Steffin; first performed in 1943. Published in English under several titles, including \emph{The Good Person of Szechwan}. The discussion follows the play's action rather than quoting a particular translation. \href{https://www.bloomsbury.com/uk/good-person-of-szechwan-9781408100073/}{English edition, Methuen Drama}.}

\section*{What the play establishes}

A play can expose a conflict without proving that every society must reproduce it. Brecht chooses the pressures, the available alternatives, and the people who arrive at the shop. We would be foolish to treat his stage as a statistical sample. We would be equally foolish to dismiss the conflict because it has been staged.

There are at least two distinct failures here. Shen Te cannot maintain a workable limit on what others take. Shui Ta acquires power over other people and uses it in ways that serve the enterprise at their expense. The second failure does not make the first imaginary. The first does not excuse the second.

This distinction matters because readers can recruit the play for opposite forms of self-congratulation. The sentimental reader blames the world for making kindness impossible, as though kindness required no judgement about competing claims. The hard businessman congratulates himself for understanding reality, as though every harsh decision in his own business had been forced upon him by the need to survive. Both can watch the play and recognise only the foolishness of the other.

To see more, we have to follow an actual decision. If a shopkeeper gives away the money reserved for rent, the landlord may close the shop. If she keeps the rent money, someone may remain hungry. She cannot make both consequences disappear by describing her intention. The choice has a cost, and a sufficiently serious account of it will say who bears that cost.

Now change the decision. The shopkeeper can pay her workers more and still meet the rent, but discovers that desperate workers will accept less. Paying less increases her profit. This is no longer the same problem. Perhaps she offers other reasons: investment, risk, a reserve against bad months. Those reasons may deserve examination. They do not acquire validity merely because an earlier refusal was necessary.

The word \emph{survival} is particularly useful to people who would prefer not to identify the point at which survival became expansion, and expansion became personal enrichment. A firm can face a genuine threat in January and exploit the memory of it in June. An institution can preserve emergency powers after the emergency has passed. Someone who once needed to become hard can grow attached to the freedom that hardness gives him.

It is uncomfortable to admit that a necessary act may help form an unnecessary habit. We prefer to think that character is settled before circumstances arrive, so that good people use power well and bad people misuse it. Brecht makes the relation less convenient. The role that protects Shen Te also changes what she can do to other people.

\section*{A limit that can be explained}

Suppose a small organisation provides beds for people who would otherwise sleep outside. It has twenty beds. On a cold night, twenty-three people arrive. The staff may find temporary space, call other organisations, or seek emergency accommodation. These efforts matter. They can change what is possible that night.

If no additional safe place can be found, however, calling the limit heartless does not supply a twenty-first bed. Someone must decide what to do. A rule based on immediate physical danger will favour some people over others. A queue favours those who arrive first. Personal discretion may notice circumstances a rule misses, but may also favour the articulate, familiar, or attractive. Every arrangement excludes somebody.

The need to decide does not establish that any decision is acceptable. Staff cannot convert finite capacity into a general licence for contempt. They can explain the rule, record refusals, report unmet need, and challenge the circumstances that keep the capacity so small. They can also conceal refusals, blame rejected applicants, and use the existence of twenty beds as evidence that they have fulfilled every obligation. Both institutions may publish photographs of their occupied rooms.

A defensible limit remains open to an account of what it costs. An evasive limit is presented as though those outside it had ceased to exist.

This is why a budget is necessary but not morally conclusive. It establishes what can be spent under a given arrangement. It does not, on its own, establish that the arrangement deserves preservation. A shelter may be unable to enlarge itself. A wealthy city that funds it may have more choices. Following the money can reveal the difference between a constraint encountered by the person at the door and a constraint chosen by someone who never has to stand there.

The reverse evasion is also possible. To say that the city could provide more does not make the night worker responsible for providing what the city has withheld. Holding the nearest sympathetic person answerable for every failure above her is an efficient way to exhaust the people who still help.

Shen Te's visitors do something of this kind. The suffering in the town is real, and the little shop becomes its most accessible address. Accessibility is mistaken for sufficient capacity. Her willingness to listen becomes a reason to keep asking her.

\section*{Good intentions need administration}

Institutions that promise mercy have always had to confront the unappealing work of maintaining themselves. A religious community needs food, buildings, rules of membership, and some means of deciding what happens when members injure one another. A teaching about forgiveness does not settle a dispute about who may withdraw money from its account. A command to welcome strangers does not specify what should happen when a guest threatens another guest.

It is easy to expose the distance between an institution's elevated language and its mundane procedures. The distance is sometimes evidence of hypocrisy. A religious authority that preaches humility while protecting an abuser is not merely attending to unavoidable administrative detail. It is using the survival of the institution as an excuse to betray people for whose protection it claims authority.

Yet an institution without procedures can betray people too. Informality often benefits whoever is already influential. If no one writes down who may make a decision, the decision may be made by the person whom everyone is afraid to contradict. If all disputes must be settled through loving mutual understanding, the less powerful participant may be pressured to forgive until the stronger one is satisfied.

The dull rule can protect the person whom warmth fails. A recorded payment prevents the generous superior from later remembering a gift as a debt. A grievance procedure can give someone a route that does not require winning the affection of the accused person's friends. A clear end to a volunteer's shift can protect her from a supposedly compassionate manager who regards exhaustion as proof of devotion.

None of these arrangements guarantees justice. Each can become cumbersome or be manipulated. Their value lies in the specific abuse they make harder. To condemn them because they do not look generous is to judge a lock by its lack of hospitality while ignoring the person who needs to sleep behind it.

The harder question comes when the administration begins to serve itself. Funds raised for the work preserve offices whose contribution to the work is obscure. Criticism becomes disloyalty. Staff learn that protecting the organisation's reputation matters more than reporting a failure to its beneficiaries. The institution still speaks in the name of those it helps, but those people have little power over what is said in their name.

Here the cousin has ceased to be an occasional expedient. He has become the permanent management, and Shen Te appears chiefly in the publicity.

\section*{Who is allowed to refuse?}

There is an inequality concealed in some praise of kindness. A powerful person's refusal may be called firmness. A dependent person's refusal may be called ingratitude. The same family can celebrate a son's determination to pursue his career and condemn a daughter's reluctance to interrupt hers for care. The relevant difference is not necessarily the seriousness of the need. It may be the success with which a particular person has been assigned the duty to meet it.

Someone trying to change this arrangement may initially seem less good. She no longer answers every call. She insists on an agreed contribution from others. She gives enough notice to make another arrangement possible and then actually leaves. People who experienced her previous availability as a natural condition can describe the change as a decline in her character.

She need not invent a male cousin, but the temptation to borrow another authority is understandable. A doctor's instruction, a new employer, or a partner's objection can make a refusal socially acceptable where her own statement that she is tired would not. The excuse protects her, at the cost of confirming that she is not entitled to decide for herself.

There are also people who learn the language of boundaries as a means of escaping ordinary reciprocity. They accept help readily and describe requests to return it as infringements upon their autonomy. They have found a vocabulary in which consideration for others is always somebody else's unhealthy expectation. Nothing in the word \emph{boundary} settles whether the refusal is fair.

We have to return to the relationship: what was undertaken, what each person can do, what alternatives exist, and what happens when the promise is broken. This is slower than sorting people into the generous and the selfish. It also allows us to notice when yesterday's victim of an unfair demand becomes tomorrow's author of one.

\section*{What the gods avoid}

The gods' departure is cruel because it preserves their satisfaction at Shen Te's expense. They have obtained the evidence they wanted. A good person exists. Her inability to live according to their demand becomes a difficulty she must manage after they have gone.

There is a familiar human version of this departure. An adviser recommends a generous course and does not remain to help carry it out. A politician announces an entitlement and leaves the agency to ration inadequate resources. A board praises the staff for doing more with less until the staff cannot do more. When the arrangement fails, the people who made the promise have another engagement.

It would be convenient if the answer were merely to keep them in the room. Sometimes that helps. Making someone follow a decision through its consequences can puncture the pleasure of issuing it. But presence alone does not overcome a conflict of interest, and a person can watch suffering without revising a cherished demand.

What is missing is the possibility that the demand itself must change, or that its author must provide the means to meet it. An instruction that cannot be questioned and need not be financed grants its issuer a remarkable freedom. Someone else will supply the labour, accept the loss, or be blamed for the failure.

Shen Te's difficulty should therefore leave us suspicious of immaculate advice. It should also leave us suspicious of the person who claims that every cruelty is the unavoidable price of keeping the shop open. We need to see the accounts, the wages, the alternatives, and the people sent away. The theatre ends before that inquiry is completed. Outside it, there is still tomorrow's opening time, and someone must decide how the shop will be run.

\chapter{The Person Who Is Not in the Room}

The donated boxes arrive after the flood. Imagine them stacked under a tarpaulin beside a community hall: coats, shoes, toys, tins, bedding, things that seemed useful when someone packed them at home. Volunteers have collected them, borrowed a van, bought fuel, and driven for hours. They have done more than express sympathy.

At the destination, another volunteer must find somewhere dry to put the boxes. Someone must open them, sort their contents, check what can be used, and dispose of what cannot. The hall already contains families who have lost their homes. The people working there need information about medicines, missing relatives, power, and water. They may have plenty of coats and no appropriate way to store donated food.

The arriving volunteers expect gratitude. The receiving volunteer feels another task being added to a day that has ceased to have manageable limits. Neither need be a bad person. They are standing at different ends of the same gift.

Disaster researchers Mary Nelan, Samantha Penta, and Tricia Wachtendorf studied this mismatch after Hurricane Sandy and the Oklahoma tornadoes of 2013. Donors and organisers could prefer goods because goods felt personal and their delivery made help visible. Recipients and relief workers had to cope with the resulting distribution of useful, excessive, and unsuitable items. Their research describes how unwanted donations can produce what relief workers call a second disaster.\footnote{Mary M. Nelan, Samantha Penta, and Tricia Wachtendorf, ``Paved with Good Intentions: A Social Construction Approach to Alignment in Disaster Donations,'' \emph{International Journal of Mass Emergencies and Disasters} 37(2), 2019, pp. 174--196. See the \href{https://hazards.colorado.edu/research-counts/the-disaster-cycle/avoiding-the-second-disaster-of-unwanted-donations}{research account published by the Natural Hazards Center}.}

The observation is small enough to be understood and large enough to be embarrassing. A sincere wish to help does not establish what help is needed. The person whose response we are most eager to demonstrate may be the person we have been least careful to consult.

\section*{The completion that comes too early}

For the donor, the work can feel complete when the van leaves. The things are no longer in the garage. They are going to people in need. There is a photograph of the loaded vehicle and a satisfying account of collective effort.

For the recipient, the work is beginning. Transport has moved the goods without necessarily making them usable. A sack of mixed clothing is not yet a coat in the right size on a person who needs one. A machine without the right spare part is not yet working equipment. An unfamiliar medicine is not made appropriate by the fact that it was expensive elsewhere.

It is tempting to treat these as mere logistical details, to be discussed after the moral importance of giving has been established. But the details determine whether the gift helps. If arranging it consumes more scarce labour than its contents save, the moral description needs to change with the result.

This does not require a donor to become an expert in emergency supply chains. It requires accepting that the people handling the emergency may know more about the immediate need, and that their answer can be inconvenient. They may want money rather than goods, fewer volunteers rather than more, or a specific delivery next week rather than an impressive delivery today. Following that answer deprives the donor of some control over the appearance of the act.

Distrust of charities can make this difficult. People reasonably want to know whether money will be wasted. A material donation seems to settle the matter: there is the blanket, and there is the person receiving it. Yet the picture excludes the cost of delivering, storing, and distributing the blanket. It also excludes the other item that the recipient might have chosen instead. Tangibility gives one kind of reassurance while leaving other questions unanswered.

The relevant comparison is between available ways of meeting the need. It is not between a visible object and an imaginary organisation without overheads. Somebody must organise even the most direct act of relief. Refusing to pay for that work does not remove it; it may leave it to someone already overwhelmed.

There are circumstances in which goods are urgently required and money cannot summon them locally. There are also circumstances in which purchasing nearby is faster and supports functioning suppliers. The difference depends upon what is available, where it is, and how it can move. A universal preference for goods or cash would repeat the original mistake of deciding before listening.

\section*{Where the old shirt goes}

An ordinary clothing donation involves a longer chain and a quieter version of the same uncertainty. A bag enters a collection point. Some garments can be sold locally. Others enter commercial sorting and export. Some are worn again; others become material for other uses or waste that somebody must manage.

The European Environment Agency describes a trade whose outcomes are far less certain than the comforting image of a garment passing directly from one wardrobe to a grateful stranger. Quality, demand, sorting capacity, and the eventual treatment of unusable material all matter. The fact that an item was collected for reuse does not prove that it was reused.\footnote{European Environment Agency, \href{https://www.eea.europa.eu/en/analysis/publications/eu-exports-of-used-textiles}{\emph{EU Exports of Used Textiles in Europe's Circular Economy}}, 2023. The subsequent discussion distinguishes possible effects on traders, consumers, production, and waste disposal; it does not assume one effect in every receiving country.}

There is nothing inherently discreditable about selling donated clothing. A sale can finance charitable work, provide an affordable garment, and pay the person who sorts or repairs it. The buyer is not cheated merely because the original owner gave it away. Nor is the original owner automatically entitled to dictate an entire subsequent economy because she surrendered a shirt.

The difficulty appears when the word \emph{donation} prevents scrutiny of the whole transaction. A low-quality consignment may leave a trader with little worth selling and a large disposal problem. A receiving town may bear waste costs absent from the exporter's account. Conversely, a ban intended to protect domestic manufacturers can deprive low-income buyers of affordable clothing and sellers of a livelihood. Whether domestic production then expands is another question, not a reward that follows automatically from the nobility of the intention.

We do not learn enough by assigning the entire trade a benevolent or predatory character. We need to know which goods can actually be used, who pays for the rejects, what alternatives buyers have, and what a proposed change does to their prices and employment. These questions can make a familiar little act of tidying seem absurdly complicated. The complication was already there. The collection bin allowed the donor not to see it.

There is a particular discourtesy in assuming that poverty makes a person grateful for anything. Need may make choice painfully narrow; it does not abolish preference, climate, size, taste, or dignity. If we would regard an object as unusable in our own home, a vague reference to people less fortunate does not tell us where it becomes useful.

The better donation may be less dramatic: a requested item in good condition, a payment that permits a purchase, or the decision to buy less in the first place. It may produce no story worth telling at dinner. That is a difficulty for the story, not for the person whose cupboard or street is spared our rubbish.

\section*{The tenant we can interview}

Consider another scene. A tenant receives notice of a rent increase she cannot pay. She has lived in the flat for years. Her work, friends, and medical care are nearby. A policy limits the increase, and she stays. Here is a recognisable beneficiary who can explain exactly what the decision means.

The person who might have rented another flat in the same neighbourhood next year is harder to show. She has no address there yet. If fewer flats are offered for rent, she may search longer, live farther away, or never arrive. Her loss has no moment as visible as the saved tenancy.

This difference in visibility matters politically. The existing tenant has a specific threatened home. The prospective tenant has a possibility. A government can appear humane by counting only the first, while an opponent can appear economically sophisticated by discussing supply as though forced moves were an insignificant nuisance. Neither has described the whole problem.

A study of San Francisco's 1994 expansion of rent control by Rebecca Diamond, Tim McQuade, and Franklin Qian illustrates the conflict. The policy reduced displacement among covered tenants. The affected landlords also reduced the supply of rental housing by about fifteen per cent, through routes including sales to owner-occupants and redevelopment. This was a finding about the affected properties and a particular policy, not a claim that fifteen per cent of every city's housing disappears under any rent regulation.\footnote{Rebecca Diamond, Tim McQuade, and Franklin Qian, \href{https://www.aeaweb.org/articles?id=10.1257/aer.20181289}{``The Effects of Rent Control Expansion on Tenants, Landlords, and Inequality: Evidence from San Francisco''}, \emph{American Economic Review} 109(9), 2019, pp. 3365--3394, doi:10.1257/aer.20181289.}

That last distinction is essential. Controls differ in their treatment of new construction, maintenance, vacancy, inflation, and conversion. Housing markets differ too. A study cannot be converted into a universal number by dropping the description of what was studied.

Still, the mechanism is intelligible. If a rule changes the return from letting a property while leaving other uses available, owners may change its use. Calling the owners greedy does not prevent the change. Closing a particular route of avoidance may help, but then its effects must be followed in turn. A policy's beneficiaries are not its only participants.

Nor does saying that prices reflect scarcity answer how the scarcity was created or whether every available response is acceptable. Restrictions on building, limited infrastructure, concentrated ownership, local political resistance, financing conditions, and the location of employment can all matter. Telling the displaced tenant to appreciate the information contained in a price is an especially bloodless substitute for housing policy.

The practical question is how to protect people during adjustment while making room for others to live. That can involve support for tenants, more construction, different rules for development, and public provision. The combination has to be paid for and administered. It cannot be replaced by the belief that a legal promise creates the resources necessary to honour it.

The missing person is not always a future tenant. She may be the cleaner commuting two hours because the people who praise the neighbourhood's inclusive character oppose another building in it. She may be a young adult remaining in a damaging household because every affordable tenancy is already occupied. The existing residents' attachment to their surroundings is real. So is the attachment others might form if they were allowed to enter them.

\section*{The field outside the picture}

A field of crops offers an attractive picture of an alternative to oil. Plants absorb carbon as they grow. Fuel made from them releases carbon when burned. Unlike an oil deposit, the crop can be grown again. It is easy to see why the word \emph{renewable} sounds like the conclusion of the argument.

But a field is also land that could have been used differently. If a crop previously sold for food is diverted to fuel, people do not automatically cease wanting the food. Prices, consumption, yields, and production elsewhere may change. Another field may be brought into use. Clearing vegetation or disturbing soil can release stored carbon before the first litre of the new fuel is burned.

The effect can therefore occur far from the plant whose emissions are being counted. A domestic production target may influence land use in another country through a sequence of market responses. Those responses are harder to estimate than the fuel leaving a factory, but difficulty of measurement does not establish an effect of zero.

A prominent 2008 study by Timothy Searchinger and colleagues argued that emissions from indirect land-use change could reverse the expected greenhouse benefit of particular crop-based biofuels. Its numerical results depended on modelling assumptions and prompted substantial debate.\footnote{Timothy Searchinger and colleagues, \href{https://searchinger.princeton.edu/publications/use-us-croplands-biofuels-increases-greenhouse-gases-through-emissions-land}{``Use of U.S. Croplands for Biofuels Increases Greenhouse Gases through Emissions from Land-Use Change''}, \emph{Science} 319, 2008, pp. 1238--1240, doi:10.1126/science.1151861. The discussion identifies the mechanism without adopting a universal estimate for all fuels.} The enduring question is what happens when the accounting follows the displaced production beyond the original field.

This does not give us one verdict on every biofuel. Wastes, residues, crops, and production methods can have different consequences. Yields can improve. By-products may replace other products. Land can change use in ways that are more or less damaging. The comparison needs to specify the feedstock, process, alternative fuel, and period over which the effects are counted.

Food prices matter alongside carbon. A rise that benefits a crop seller can hurt a household buying most of what it eats. Energy security may supply a separate reason for domestic production. These effects deserve to be stated separately. A benefit on one measure should not be used to make a cost on another disappear.

The farmer supplying the fuel may have done exactly what public policy encouraged. Criticism should follow the rule and its effects rather than invent a villain at the nearest visible point. Equally, a policy cannot defend itself merely by showing the useful income it provides to that farmer. The people paying more for food, and the land affected elsewhere, belong in the account too.

The green field remains green in the photograph. That is why the photograph cannot finish the comparison. The difficult part has taken place outside it.

\section*{The people a successful reform leaves behind}

Even a reform that works overall can create a concentrated loss. A new process can reduce the cost of a good while destroying a particular livelihood. A cleaner technology can improve air quality while making specialised equipment and skills less valuable. A more efficient service can close a local office around which an elderly person's routine has formed.

To acknowledge the loss is not to establish that the reform should be stopped. It establishes that its success is a comparison, not a state in which nobody has been hurt. There are cases in which compensation, retraining, a slower transition, or a continuing minimum service is warranted. There are also demands for protection that would impose large continuing costs upon others for a small group's benefit. We cannot decide which is which from the intensity of the group's attachment alone.

The people gaining may be numerous, dispersed, and barely conscious of the gain. The people losing may be few, identifiable, and desperate. In other cases, the gain is concentrated in an organised industry and the cost spread across customers who each lose too little to protest effectively. Political visibility can favour either side, depending upon how the decision is arranged.

This is why simply siding with the person who tells the most affecting story is unreliable. Stories help us understand a loss that a total conceals. They can also conceal the total. A moving account of one threatened job does not tell us how many other jobs a protective measure may prevent. A national gain does not tell us whether its beneficiaries could readily absorb the cost of assisting the people it displaces.

The proper response is not to forbid stories. It is to let other people into them. The worker whose factory closes deserves to be heard, and so does the buyer whose household budget depends on affordable goods. The conservation proposal should include the person who lives beside the protected animal. An energy policy should include the family whose bill it raises. A charity should include the people receiving the things that donors enjoy giving.

\section*{Following a decision far enough}

No individual can trace every consequence of every purchase or vote. A demand to do so would produce paralysis or a pretence of omniscience. Most of the time we rely on institutions, habits, and a limited understanding of the circumstances. The question is whether we allow credible evidence of a neglected consequence to disturb the account we prefer.

There is a difference between failing to anticipate a cost and refusing to acknowledge it once it appears. There is a further difference between acknowledging it and making its bearer responsible for remaining quiet so that the public can continue to believe in the policy. The last step turns someone else's injury into a public-relations inconvenience.

We can make a more modest demand of ourselves. Before praising an action, follow it one stage beyond the photograph. Who opens the boxes? Who maintains the equipment? Who supplies the promised service? Who is left searching after the beneficiaries have been counted? We will not discover everything. We may discover the person whose absence made the result look uncomplicated.

Sometimes we will still make the same decision. We may send the money, impose the regulation, close the plant, or protect the animal. Acknowledging a cost need not reverse the judgement. It changes the way we are entitled to describe what we have done, and may change what we owe afterwards.

The more melancholy lesson is that some losses will remain after the best available choice. There may be no arrangement in which every person keeps a familiar home, a livelihood, a landscape, and the future once expected. We can try to prevent avoidable harm and share unavoidable burdens without inventing a final reconciliation. The person who lost something should not also have to endorse our account of the occasion as an unqualified success.

\part{What the World Will Allow}

\chapter{A Climate That Has Never Stood Still}

An old photograph shows a glacier reaching down a valley. A new photograph, taken from the same place, shows bare ground where the ice used to be. The comparison is striking. Even without knowing much about glaciers, a viewer can see that something substantial has changed.

Two people can look at these pictures and immediately begin different arguments. One sees proof that industrial society is destroying the Earth. The other says that glaciers have advanced and retreated before, and treats the observation as though it disposed of the matter. Each moves too quickly. The photographs show a change in a particular glacier. They do not, by themselves, establish its causes, the contributions of different causes, or the best response.

Answering those questions does not make the change less real. It makes the photographs the beginning of an inquiry rather than an occasion for displaying convictions already held.

The same discipline is needed at the scale of the planet. The climate has changed throughout Earth's history. Human beings are now changing it as well. Neither statement cancels the other. An argument that cannot accommodate both has become an obstacle to understanding what is happening.

\section*{Before there were factories}

Long before human beings burned coal in factories, the Earth experienced climates very different from those of our lives. Large ice sheets expanded and withdrew. Seas stood at different levels. Regions now dry were wetter; familiar habitats did not always occupy their present places. The comfortable image of a naturally stable world, disturbed only when people arrived, has no adequate basis in this history.

During the recent succession of glacial and interglacial periods, changes in Earth's orbit and the tilt and orientation of its axis altered where and when sunlight reached the surface. Summer conditions in northern latitudes mattered because snow left after one winter could survive the next summer and accumulate. Expanding ice reflected more sunlight. Changes in oceans, vegetation, and greenhouse gases amplified and complicated the response.

The important idea is that an initial change can set other changes in motion. The process need not have one cause operating alone. Nor must the factor that starts a transition remain its largest influence. Carbon dioxide can increase as a consequence of warming and then contribute to further warming. In another situation, adding carbon dioxide can begin a warming response. Calling it a feedback in the first case does not prevent it from being a cause in the second.

There are unresolved questions about the timing and dynamics of glacial cycles. A broad explanation can be well supported without every feature of the history being settled. That is a common condition of useful knowledge. It gives us reason to specify our confidence, not to pretend that explanation has failed because it is incomplete.

The past also makes the word \emph{natural} an unreliable reassurance. A naturally occurring drought can kill a crop. A volcanic eruption can disrupt a growing season. No human responsibility for the cause is needed for human beings to have an urgent interest in protection from the result.

It follows that eliminating human influence, even if such a thing were possible, would not eliminate climate risk. It does not follow that adding a new influence is harmless. A house can burn after a lightning strike; that is no defence of careless wiring. The history of other causes does not answer a question about the cause currently being investigated.

\section*{The difficulty of remembering a climate}

Written history adds human experience to the record, but it does not provide a complete set of measurements. A chronicler tells us about a terrible winter because it was terrible. A farmer records an early harvest because it affected work or income. An account may be exact about local suffering while saying little about conditions elsewhere.

Natural records help extend and test these observations. Tree rings, ice cores, sediments, corals, and other materials preserve evidence related to past conditions. They are not interchangeable thermometers. A tree's growth can respond to moisture as well as temperature. A record from one mountain does not describe an entire continent. Reconstructing climate requires understanding what each record can indicate and comparing different kinds of evidence.

This is why a phrase such as \emph{the Medieval Warm Period} needs care. It can be useful for discussing regional conditions over a broad interval. It becomes misleading when it suggests that every part of the world warmed at the same time by the same amount. Raphael Neukom and colleagues, examining the last two millennia, found that familiar preindustrial warm and cold periods lacked the near-global coherence of recent warming.\footnote{Raphael Neukom and colleagues, \href{https://doi.org/10.1038/s41586-019-1401-2}{``No Evidence for Globally Coherent Warm and Cold Periods over the Preindustrial Common Era''}, \emph{Nature} 571, 2019, pp. 550--554.}

That finding does not erase warm medieval conditions in a particular region. It changes the geographical claim we can make from them. A history of farming in Greenland may be relevant to Greenland. To use it as a global thermometer requires evidence that the farming history cannot supply on its own.

The Little Ice Age raises a similar difficulty. The name gathers together cold conditions that differed across places and decades. In parts of Europe, colder seasons and advancing glaciers threatened livelihoods. We should neither turn these experiences into a uniformly frozen planet nor dismiss them because the pattern was not uniform.

Even the boundaries of shorter climatic episodes are debated. Research linking sixth-century volcanic eruptions with prolonged cooling helped establish the term \emph{Late Antique Little Ice Age}. Subsequent criticism questioned how long and how widely that cooling persisted. The disagreement concerns the interpretation and reach of evidence, not whether every uncomfortable historical observation should be admitted or excluded according to its usefulness in present politics.\footnote{Ulf Büntgen and colleagues, ``Cooling and Societal Change during the Late Antique Little Ice Age from 536 to around 660 AD,'' \emph{Nature Geoscience} 9, 2016, pp. 231--236, doi:10.1038/ngeo2652; see also \href{https://www.nature.com/articles/ngeo2926}{``Limited Late Antique Cooling''}, 2017, and the \href{https://www.nature.com/articles/ngeo2927}{authors' reply}.}

These histories can make us more careful about modern claims. They cannot settle those claims by analogy alone. The fact that societies endured environmental change tells us that adaptation is possible. It does not tell us what adaptation will cost at a different rate of change, with different populations, coastlines, infrastructure, and political conditions.

\section*{A poor harvest is not a complete history}

There is another temptation in climate history: finding a temperature change near a political collapse and treating the first as an explanation of the second. A state experiences drought, prices rise, conflict follows, and a complicated history becomes a simple warning about what climate does to civilisation.

Weather and climate plainly affect food, water, disease conditions, and movement. But people do not encounter these effects without institutions. Grain can be stored or exported. A shortage can be relieved through trade or worsened by a blockade. Taxation can leave cultivators with reserves or remove them. A government can distribute assistance or use the emergency to reward its supporters. These choices affect the consequences of the same failed harvest.

Suppose two districts receive similarly inadequate rain. One has access to stored food, passable roads, and money to import supplies. The other has damaged transport and a population already close to hunger. Their outcomes can differ enormously. The rainfall has not become irrelevant in the first district. Its human effect has been altered by what people previously built and what they now do.

This complicates both nostalgia and catastrophe. A society may have lived with lower emissions while exposing its members to far more immediate environmental danger. Another may possess remarkable protective technology and still allow people to die because access to it is unequal or institutions fail. There is no reason to identify low energy use with safety, or wealth with immunity.

Historical explanation becomes more useful when it specifies a route from the physical change to the human outcome. Which crop failed? Which region supplied food? What happened to transport and prices? Who could move? Who could not? A climate label attached to a century is too coarse to answer these questions.

The same questions will matter when we consider our own future. They are more demanding than imagining that a rising global temperature translates directly into one predetermined social condition.

\section*{What greenhouse gases do}

The basic physical account can be stated without turning the reader into a specialist. The Earth receives energy from the Sun. Some sunlight is reflected; some is absorbed. The Earth also emits energy towards space, chiefly as infrared radiation. Over time, a persistent difference between incoming and outgoing energy changes the amount of heat in the Earth system.

Greenhouse gases absorb and emit infrared radiation at particular wavelengths. Increasing their concentration changes the way energy escapes to space. With more carbon dioxide, the planet initially loses less energy than it otherwise would under the same conditions. Warming, together with other changes in the system, increases outgoing radiation and moves the balance towards a new state.

The Sun need not become hotter for this to happen. Nor does the atmosphere need to become a thinner screen. The greenhouse effect is also distinct from the problem of stratospheric ozone depletion. A simple explanation should simplify the mechanism, not replace it with a different one.\footnote{NASA, \href{https://science.nasa.gov/earth/climate-change/global-warming/}{\emph{Global Warming}}, and \href{https://science.nasa.gov/climate-change/faq/is-the-ozone-hole-causing-climate-change/}{\emph{Is the Ozone Hole Causing Climate Change?}}. The explanation here describes the energy balance; a glass greenhouse also limits convective heat loss, so the familiar name should not be treated as an exact mechanical analogy.}

The subsequent response includes changes in water vapour, clouds, snow, ice, and the oceans. These affect how much warming occurs and how it is distributed. The complexity of the response is one reason that estimates have ranges. It does not make the original absorption of infrared radiation a matter of political preference.

The atmosphere is not the whole system. The ocean stores a great deal of heat, and different parts of it exchange heat on different timescales. A cool season in one place does not tell us whether the planet as a whole has accumulated energy. Nor does an especially hot afternoon establish the size of the long-term global trend.

The rise in atmospheric carbon dioxide is measured, and its connection with human activity is supported by more than coincidence. The quantities emitted, changes in the carbon's isotopic composition, and other observations help identify fossil carbon as the main source of the increase. This is a claim about physical evidence that can be inspected, not an inference from whether industrialisation looks morally attractive.\footnote{NOAA Global Monitoring Laboratory, \href{https://gml.noaa.gov/infodata/faq_cat-3.html}{carbon-cycle questions and answers} and \href{https://www.gml.noaa.gov/outreach/isotopes/mixing.html}{explanation of carbon isotopes and sources}.}

\section*{What is different now}

The relevant modern question is therefore not whether climate can change without us. It is how much of the observed change is attributable to our activities, alongside natural influences and internal variability.

In its 2021 assessment, the IPCC estimated human-caused warming from 1850--1900 to 2010--2019 at about 1.07\textdegree{}C, with a likely range of 0.8\textdegree{}C to 1.3\textdegree{}C. Greenhouse gases produced warming partly offset by other human influences, especially aerosols. Natural drivers and internal variability did not explain the long-term global rise. These numbers concern specified periods; they are not a temperature reading for the day this book is opened.\footnote{IPCC, \emph{Climate Change 2021: The Physical Science Basis}, \href{https://www.ipcc.ch/site/assets/uploads/2021/09/Doc_4_Rev_2_AR6_WGI_Approved_Summary_for_Policymarkers.pdf}{Summary for Policymakers}, especially A.1 and Figure SPM.2.}

An average increase of around a degree may sound small to someone accustomed to much larger changes between morning and afternoon. The comparison is misleading. A daily change at one location and a persistent shift in a global average describe different things. A change in the average can alter the frequency of extremes and occur alongside much larger changes in particular regions or seasons.

Still, the average is not a complete description. Rainfall, soil moisture, sea level, ocean conditions, and the timing of seasonal changes matter independently to people trying to grow food or maintain a city. A single global number is useful because it summarises an important feature of the system. It becomes unhelpful when asked to stand in for every local consequence.

For this reason, attributing a particular disaster requires more work than pointing to global warming. A flood involves rainfall, ground conditions, river management, drainage, exposure, and vulnerability. A changing climate may have increased the probability or intensity of the rainfall. Building in a floodplain may have increased the damage. Neglected drainage may have made a local result worse. These causes can operate together.

The possibility of multiple causes frustrates political argument because participants often want exclusive responsibility. If climate contributed, one side wants every other cause to disappear. If poor planning contributed, the other wants climate to disappear. The people cleaning mud out of their homes need a better account of which risks can actually be reduced.

\section*{What a vanished glacier means}

We can now return to the photographs with more useful questions. How much mass has the glacier lost? Over what period? What changes in snowfall and temperature affected it? How does its behaviour compare with other glaciers? How long does the position of its front take to respond to a change in its mass?

A glacier's outline is not an instantaneous temperature gauge. Ice flows, gains snow, melts, and responds with delays. One glacier advancing does not refute a wider pattern of loss. A dramatic retreat at one site does not establish the exact rate everywhere else. The physical detail matters because the consequences differ by place.

A community using meltwater needs to know when water arrives, what other sources exist, and how future flows may change. An initial increase in melting can supply more water while depleting the stored ice that supplied it. A tourism business faces another set of consequences. A person who loved the sight of the glacier has suffered a loss that is neither a water account nor a business calculation.

That last attachment should not be mocked as though only priced things were real. Landscapes help form memory. A place where someone walked with a dead parent may matter for reasons that cannot be transferred to a different, more convenient view. The glacier's earlier advances and retreats do not make the present loss unreal to that person.

But attachment cannot establish that the chosen landscape represents a permanent natural entitlement. A photograph records a date. Wanting the place to remain as it was is intelligible; assuming that it would have done so without human activity requires evidence. We can grieve a change while being exact about its causes and realistic about what can still be preserved.

\section*{The science does not write the policy}

Nothing in this account tells us that every measure proposed in the name of climate is sensible. A physical finding cannot calculate a household's ability to replace its heating system without information about that household. It cannot decide a country's energy mix without knowing its resources, infrastructure, demand, and alternatives. It cannot make a technology undesirable merely because the wrong political constituency likes it.

Nor is criticism of a policy an answer to the physical findings. A poorly designed subsidy does not change the absorption properties of carbon dioxide. An activist's exaggeration does not put lost ice back into a valley. Discovering foolishness in public debate gives us a reason to criticise the foolishness; it does not grant exemption from the rest of the evidence.

This division of work is uncomfortable for anyone seeking a single allegiance. We may accept the evidence of human-caused warming and reject a particular transition timetable. We may support rapid deployment of a technology while disputing the claim that it solves every problem. We may judge a proposal inadequate because it does too little, too expensive because it does little for its cost, or unjust because its burdens have been placed on people unable to bear them.

These are arguments worth having. They require descriptions of consequences, feasible alternatives, and the values used to compare them. The climate record will constrain those arguments without completing them.

The changing world gives us neither permission to do anything nor a command to do one particular thing. It gives us conditions within which decisions have to be made. We cannot preserve every familiar feature of the landscape, and we need not add avoidable damage simply because some change is inevitable. Between those observations lies the difficult work that a photograph can make urgent but cannot perform for us.

\chapter{What a Forecast Can Tell Us}

Imagine a town deciding how high to build a flood barrier. The engineers cannot tell the council exactly which storms will occur over the next seventy years. They can study past water levels, the river, the ground, and possible changes in rainfall and sea level. They can calculate how different designs perform under different conditions. They can explain which failures are unlikely but costly, and which are more ordinary.

The councillors still have to decide. A higher barrier costs more and may cut the town off from its waterfront. A lower one may be overtopped sooner. Delaying the project preserves money today while accepting risk. No option is equivalent to having no uncertainty.

One councillor announces that the experts have predicted the future and must be obeyed. Another says that nobody can know the future, so the report is worthless. Both avoid the work that the report makes possible: deciding under stated uncertainties, with a record of the reasons.

Climate projections belong to this kind of work. Their value depends neither on infallibility nor on our willingness to be frightened. It depends on what they represent, how they have been tested, and whether the question we ask is one they can help answer.

\section*{The condition attached to the line}

A projection says what a model produces under specified conditions. If emissions follow one course, the estimated response differs from the response under another. The world cannot follow every course at once. Showing several does not mean the modeller expects several incompatible futures to occur.

This sounds obvious until a graph enters public argument. Its most alarming line becomes the future. The conditions printed beside it disappear in the retelling. An estimate of what could happen under a high-emissions pathway becomes a claim about what certainly will happen, sometimes attached to a demand for a policy that was never represented in the calculation.

The opposite mistake is to choose a low-emissions pathway and treat its existence on a graph as reassurance that the necessary changes are already under way. A model can describe a world with rapid technological deployment without demonstrating that governments will fund it, firms can supply it, or the relevant institutions will permit it. A possible route is not a completed road.

It helps to distinguish two questions. Given a set of atmospheric changes, how does the physical climate respond? And what atmospheric changes will human societies actually produce? Models of the physical response can be useful even when political and economic developments are difficult to predict. But a projection that combines both must be assessed with both in view.

The word \emph{scenario} is therefore not a small qualification that cautious scientists add to escape embarrassment. It identifies what sort of claim is being made. Removing it changes the claim.

A person reading the graph should be able to find the assumed emissions, the time period, the baseline, and the range of responses. If these are inaccessible or concealed beneath a headline, the presentation has made informed judgement harder even when the underlying work is sound.

\section*{A forecast old enough to examine}

James Hansen and his colleagues published a well-known set of climate simulations in 1988. Their paper used three scenarios for changes in trace gases and aerosols. Scenario A continued strong growth, B assumed slower growth, and C imposed a sharp curtailment so that net forcing stopped increasing after 2000. The paper was not one unconditional line announcing precisely what must happen.\footnote{James Hansen and colleagues, \href{https://www.giss.nasa.gov/pubs/abs/ha02700w.html}{``Global Climate Changes as Forecast by Goddard Institute for Space Studies Three-Dimensional Model''}, \emph{Journal of Geophysical Research} 93, 1988, pp. 9341--9364, doi:10.1029/JD093iD08p09341.}

To evaluate it later, we need to compare the temperatures projected with those observed. We also need to compare the assumed atmospheric influences with those that occurred. If a scenario assumes too much growth in a gas, its temperature projection can be too high even if the model's response to a given influence is reasonably accurate.

That distinction does not entitle anyone to declare the whole forecast correct after changing its inputs. The original combined projection and the model's physical performance answer related but different questions. A fair assessment reports the distinction rather than using whichever comparison is more flattering.

Zeke Hausfather and colleagues examined historical climate projections in a study published online in 2019. Most of the projections they studied were consistent with observations when differences between assumed and realised forcing were accounted for. Their analysis included the Hansen scenarios. It is evidence that earlier models contained useful information about warming, not that every original scenario became an accurate description of later human behaviour.\footnote{Zeke Hausfather, Henri F. Drake, Tristan Abbott, and Gavin A. Schmidt, \href{https://agupubs.onlinelibrary.wiley.com/doi/10.1029/2019GL085378}{``Evaluating the Performance of Past Climate Model Projections''}, \emph{Geophysical Research Letters} 47, 2020, e2019GL085378; first published online 4 December 2019.}

For the reader, the lesson is procedural. Keep the original forecast. Keep its assumptions. Choose a comparison suited to the claim. Explain the sources of error. Do not quietly substitute a newer forecast and then congratulate the old one; do not select an unfulfilled scenario and pretend it was the only one offered.

This procedure is less entertaining than declaring a famous scientist vindicated or exposed. It gives us something the entertaining verdict does not: a basis for deciding which aspects of the next projection deserve confidence.

\section*{An error that should have been stopped}

The Himalayan glacier claim was a different kind of failure. The IPCC's 2007 Working Group II report included a statement suggesting that the Himalayan glaciers could disappear by 2035 or sooner. The claim was poorly supported and should not have survived the assessment process. In January 2010, the IPCC acknowledged the problem and the failure to apply its standards properly.\footnote{IPCC, \emph{Climate Change 2007: Impacts, Adaptation and Vulnerability}, section 10.6.2; IPCC, \href{https://archive.ipcc.ch/pdf/presentations/himalaya-statement-20january2010.pdf}{statement on the melting of Himalayan glaciers, 20 January 2010}. The statement acknowledges poorly substantiated estimates and failures in applying review standards.}

There is no need to wait until 2035 to identify this error. A claim can be unjustified when published, even if its proposed test lies decades ahead. In this case the source and review were already open to examination.

The important institutional question is how such a claim acquired the authority of a major assessment. What evidence was used? What review took place? What happened to objections? How were corrections communicated? A statement does not become reliable because it appears among many reliable statements. Readers are entitled to expect the assessment's procedures to discriminate between them.

The convenient defence is that the glaciers are losing ice anyway, so the mistaken date changes nothing important. It changes the credibility of the particular claim and our assessment of the process that transmitted it. If a date is vivid enough to help persuade the public, it is important enough to correct visibly.

The convenient attack is that one error discredits every finding in the report. That saves the critic from examining the evidence behind the other findings. A mistake about the fate of particular glaciers does not refute measurements of atmospheric composition or the radiative properties of greenhouse gases. Those claims stand or fall through their own evidence.

We should be able to criticise the assessment body without being recruited into either defence. Its work is too consequential to be shielded from scrutiny, and too extensive to be fairly judged by selecting a single defective paragraph. Serious criticism can demand a correction, investigate the process, and still rely on findings that remain well supported.

\section*{What happens when the date passes}

Public campaigns favour deadlines because a date gives urgency a shape. A distant risk becomes a countdown. The tactic can motivate attention, but it creates obligations that the speaker may later prefer to forget.

Some deadlines describe a political goal. Others estimate when cumulative emissions are likely to exhaust a carbon budget associated with a particular temperature limit and probability. Others concern a projected event in a particular region. Still others are rhetorical statements about the last chance to save the world. These are not interchangeable claims.

If a target is missed, the promised action has not occurred. If an event prediction fails, the prediction needs examination. If a carbon budget estimate changes, we need to know whether emissions, observations, assumptions, or scientific understanding changed. Simply replacing the old date with a new one explains none of these things.

There is also a special difficulty when a warning helps prevent what it warned against. A successful evacuation means that predicted deaths may not occur; it does not prove that the storm was harmless. A projection based on continued emissions can differ from reality because emissions were reduced. To distinguish prevention from a mistaken warning, we need a record of the assumed conditions and the actions actually taken. Otherwise every failure can be claimed as a success, and every success can be dismissed as proof that concern was unnecessary.

Climate campaigning sometimes adds another damaging implication: if the date is missed, nothing useful remains to be done. That turns a reason for urgency into a future reason for despair. A temperature objective can be exceeded while substantial differences remain between subsequent courses. More warming can bring more damage. Avoiding some of it still matters.

This does not mean every change is smoothly reversible. Ice sheets, ecosystems, and other parts of the climate system can respond slowly or cross thresholds. Returning a global average to an earlier value would not simply restore every lost feature. The IPCC's assessment of overshoot explicitly recognises additional risks and difficulties in bringing temperatures back down.\footnote{IPCC, \emph{Climate Change 2023: Synthesis Report}, \href{https://www.ipcc.ch/site/assets/uploads/2023/03/Doc5_Adopted_AR6_SYR_Longer_Report.pdf}{Longer Report}, section 3.3. The possibility of reducing warming after an overshoot does not imply that all its consequences can be reversed.}

The world does not offer a single date before which all is recoverable and after which all effort is futile. Such a date would simplify public communication at the cost of misdescribing the problem.

\section*{One year and a lasting change}

A related confusion concerns temperature thresholds. An individual year can be warmer or cooler because of natural variability superimposed on the long-term trend. A year above 1.5\textdegree{}C relative to a preindustrial baseline is not identical to a sustained climate at that warming level. The Paris temperature goals concern the long term.\footnote{Met Office, \href{https://www.metoffice.gov.uk/blog/2023/why-1-5c}{\emph{Why 1.5\textdegree{}C?}}, 2023, explaining the distinction between annual variability and long-term warming.}

The distinction is not a trick for postponing an embarrassing announcement. It is necessary to compare like with like. If we judge a long-term target by a single warm year, consistency would require judging it reversed by a subsequent cooler year. That would turn the target into a report on annual fluctuations.

Neither should the distinction be used to make an unusually warm year seem irrelevant. It occurs within a changing climate and can bring serious effects. What it establishes about a long-term threshold must be stated correctly. The human costs of heat do not wait for an averaging convention; the convention still matters when describing the trend.

Public communication often treats qualifications as weaknesses. Yet a reader given the definition can understand it. Confusion is more likely when an institution announces a dramatic result and supplies the meaning only after objections arrive. Trust would be better served by explaining the measure before asking the public to react to it.

\section*{Not all uncertainty belongs in one range}

Suppose a projection gives a range for sea-level rise. That range may reflect uncertainty about future emissions, climate response, and the behaviour of ice sheets. A coastal authority then needs local information: land movement, tides, waves, storm surge, and the height and condition of existing defences. A global range cannot do all this work.

There is further uncertainty about exposure. How many people will live near the water? Which buildings will be there? Will drainage be improved or neglected? Will a barrier be maintained? An economic damage estimate must make assumptions about these matters as well as about the physical hazard.

A precise monetary figure at the end of the exercise can conceal a long sequence of conditional judgements. The number may be useful for comparing options within a study. It should not be presented as though someone has already counted the future repairs.

Ranges also differ in how they are constructed. A spread across model runs is not automatically a complete probability distribution for reality. Models can share assumptions or omit the same process. A low-probability outcome may be poorly quantified rather than confidently negligible. The responsible response is to explain what the range includes and what it does not.

We should ask the same of reassuring estimates. People often become demanding statisticians when the result is alarming and accept a favourable central estimate without inspecting its uncertainty. The risk does not become more knowable because we prefer the lower number.

\section*{Decisions that do not require a perfect forecast}

The flood barrier returns us to practical life. The town can do some useful things without knowing the exact water level in 2090. It can maintain drainage, avoid placing critical electrical equipment at the lowest point of a building, preserve evacuation routes, and stop treating every undeveloped floodplain as an invitation to construct homes. The value of these actions depends on local circumstances, but many can be assessed across several plausible futures.

Other investments can be designed for later adjustment. A structure might permit additional height. A planning rule might contain review dates tied to observed changes. Reserving space today can make a future modification possible. These choices have costs; flexibility is not free. But their value can be compared with the cost of committing to one forecast and discovering too late that it was wrong.

There are limits to waiting for more information. Construction takes time. Some changes are difficult to reverse. An uncertain but serious risk can justify action before its precise magnitude is known, especially when delay removes options. Equally, a declaration of urgency does not establish that the first proposed action is effective. We need to compare the consequences of delay with the consequences of acting badly.

This is where uncertainty should discipline policy rather than halt it. A robust plan works tolerably under more than one credible future. It identifies where it would fail, what observations would justify revision, and who is responsible for making that revision. A fragile plan depends on favourable assumptions while presenting itself as the only responsible choice.

The contrast is visible in ordinary maintenance. A town can build an impressive defence and fail to fund its inspection. The ribbon-cutting is counted as completion; the less visible obligation continues for decades. No improvement in climate modelling can compensate for a gate that no longer closes.

\section*{Who answers for the forecast?}

Accountability needs to match the claim. Scientists should preserve methods and explain uncertainty. Officials should say which findings they used and how values and budgets affected the decision. Advocates should correct specific claims rather than retreat to the general importance of the cause whenever challenged.

The public also has work to do. A criticism should quote what was actually said, distinguish a scenario from a prediction, and use the period and measure appropriate to the comparison. Mocking an invented version of a forecast is no more sceptical than defending an exaggerated one.

None of this requires reverence. An authority that demands trust while resisting inspection deserves pressure. A commentator who turns every uncertainty into permission to ignore a risk deserves it too. What matters is whether the pressure improves the account we have of the world.

A good forecast leaves room for events to teach us something. A bad public argument absorbs every event without changing. When warming is faster, it proves the argument; when slower, the danger is merely concealed. Or, in the opposite account, every cool interval refutes warming while every warm interval is dismissed as an exception. These positions can survive indefinitely because they have stopped risking a specific error.

The engineers in our town cannot offer that luxury. Their work will meet water. The council cannot eliminate uncertainty before spending money, and the residents cannot recover all the years lost if the decision fails. They can insist upon a clear account now, preserve the ability to revise it, and attend to the unglamorous work of maintenance. That is less consoling than knowing the future. It is also something they can actually do.

\chapter{The Cost of Changing the Energy Supply}

A refrigerator does not care which political movement favours the electricity that powers it. It needs the supply to continue. The same is true of a water pump, a hospital lift, a railway signal, and the machinery in a factory. People can disagree about the energy system while depending upon it in remarkably similar ways.

This dependence is easy to overlook where electricity is normally available. The socket makes an enormous arrangement appear to be a feature of the wall. Mines, generators, fuel contracts, reservoirs, transmission lines, transformers, control systems, repairs, and human shifts disappear behind the small act of switching something on.

Replacing parts of that arrangement is necessary if we are to reduce the emissions it produces. The difficulty is that the services cannot simply pause while the replacement is built. Homes still need heat, factories still need power, and food still needs refrigeration. A transition is a sequence of working arrangements, not only a picture of the final one.

The most revealing question for an energy proposal is often very plain: what will supply the service on the day after this change is made?

\section*{What lower consumption can mean}

Two households reduce the energy they use for heating. One has insulated its home and remains comfortable. The other cannot afford the bill and sits in a cold room. The meter records a reduction in both places. It does not record the difference in their lives.

Efficiency is valuable because it can provide a service with fewer inputs. Deprivation reduces the service. There are also voluntary changes in preference: someone may choose a cooler room or a shorter journey without suffering a serious loss. A sensible account distinguishes these possibilities rather than assigning moral value to the direction of the meter.

The distinction becomes more consequential where people have little energy access. Light, pumping, refrigeration, and mechanical power can change what a household or business can do. The international agencies tracking energy access continue to report large populations without electricity or clean cooking. These deficits concern ordinary capabilities, not extravagant consumption.\footnote{IEA, IRENA, United Nations Statistics Division, World Bank, and WHO, \href{https://www.iea.org/reports/tracking-sdg7-the-energy-progress-report-2025}{\emph{Tracking SDG7: The Energy Progress Report 2025}}.}

A wealthy visitor can admire a low-energy way of life during a brief stay and leave before its constraints become burdensome. The resident cannot depart from the labour of fetching fuel, the smoke of cooking, or the loss when food spoils. To romanticise these conditions from a comfortable home is to make someone else's restricted choices into an aesthetic preference.

There is no need to insist that development must reproduce every technology once used by rich countries. A village may gain electricity most practically from a local renewable system. A city may need a much larger network. The question is what can provide reliable services at a cost people can afford, now and as their needs grow. Climate policy becomes credible to those users when it helps answer that question.

Telling them that the world cannot afford their development is a political demand of extraordinary severity. If cleaner alternatives are wanted, financing and making them workable matters more than applauding the restraint of people who have few alternatives at all.

\section*{Cheap electricity and a complete supply}

Renewable generation has become substantially cheaper. IRENA's assessment of projects commissioned in 2024 found that ninety-one per cent of new utility-scale renewable capacity produced electricity at a lower levelised cost than the cheapest newly installed fossil-fuel alternative. That is a significant finding about new generation.\footnote{IRENA, \href{https://www.irena.org/Digital-Report/Renewable-Power-Generation-Costs-in-2024}{\emph{Renewable Power Generation Costs in 2024}}, 2025. The comparison concerns newly commissioned utility-scale renewable capacity and the levelised cost of new fossil-fuel generation; it is not a comparison with every existing plant's operating cost.}

It is not a finding that every electricity system can already be replaced at no additional cost. Levelised cost spreads a plant's costs over its expected lifetime output. It does not by itself tell us the value of that output at each hour, or supply all the costs of delivering reliable power to a particular customer.

A solar plant can produce inexpensive electricity at midday. A household still needs some electricity after sunset. If many solar plants produce at similar times, adding another may provide less value during those hours unless demand, storage, or connections change. Wind has a different pattern, but its output is also variable. The engineering problem is how these patterns fit together with the rest of the system.

Critics can misuse this distinction too. Saying that solar requires other equipment does not establish that it is uneconomic. Every generating technology requires a wider system. A thermal plant needs fuel and cooling and can fail unexpectedly. Hydropower depends on water conditions. Transmission and reserves were not invented to make wind power look viable. They are part of supplying electricity.

The relevant comparison is between complete arrangements providing a specified quality of service. How much generation is needed? What storage, flexible demand, connections, and firm capacity are required? What happens during an extended period of low renewable output? What does the alternative cost, including fuel, maintenance, pollution, and emissions?

The IEA's work on integrating wind and solar emphasises that the challenges change as their share grows. Early additions can often be accommodated with limited changes; higher shares require more substantial changes in flexibility, stability, operation, and investment.\footnote{IEA, \href{https://www.iea.org/reports/integrating-solar-and-wind}{\emph{Integrating Solar and Wind}}, 2024, especially the discussion of stages of integration and changing requirements for system flexibility.} Success at one stage gives useful experience. It does not remove the need to plan the next stage.

Numbers become misleading when each side chooses the boundary that flatters its technology. One quotes the cost of electricity leaving a new plant. Another adds every conceivable network cost to that plant while treating the existing system's infrastructure as free. Both can produce a persuasive comparison by making the omitted costs difficult to see.

\section*{A machine has to exist before it can help}

Energy transitions also involve physical production. A plan may call for generators, cables, batteries, transformers, heat pumps, charging points, or industrial equipment. These things require materials, factories, workers, permits, transport, and connections. A target expressed in gigawatts does not reserve a transformer at a factory.

Some constraints can be eased quickly. Others take years. Skilled workers cannot be trained by changing a line in a national strategy. A mine or transmission corridor cannot be assumed available merely because a scenario needs it. Suppliers may hesitate to expand if policy changes unpredictably or orders lack financing.

This is not a reason to regard the existing system as permanent. It is a reason to treat deployment as serious work. The timetable must include bottlenecks. Where demand for equipment will rise sharply, the plan needs to ask whether manufacturing can expand, how quality will be maintained, and which projects should have priority if supply falls short.

Money alone cannot immediately remove every bottleneck. More bids for the same scarce equipment may first raise its price. A subsidy can stimulate supply over time, but the response depends on firms believing that the demand will last. A policy that repeatedly stops and restarts can leave both workers and manufacturers reluctant to commit.

There are choices about timing within a household as well. Replacing an old heating system at the end of its life is different from discarding one installed recently. A landlord's incentives differ from a tenant's. A heat pump's performance depends on design and installation, and sometimes on improvements to the building. Announcing a preferred device is easier than arranging competent installations at scale.

The household that cannot finance the change may agree completely about the climate. Its lack of money is not a scientific disagreement. A government that mistakes the two will respond to a practical obstacle with a lecture and then wonder why its advice is resented.

\section*{Technologies without a moral pedigree}

An energy technology should not have to be emotionally congenial to be considered. Nuclear power, fossil generation with carbon capture, large hydroelectric projects, batteries, transmission lines, and different forms of renewable generation each raise particular questions. Their risks, costs, construction times, and performance deserve comparison in the system where they would operate.

Nuclear power can provide substantial low-emissions electricity with a different operating profile from wind and solar. That does not settle the cost or timetable of a proposed new project, the management of waste, or the quality of the institutions responsible for it. Keeping an existing plant running safely and commissioning a new design are different decisions. Treating both as a single vote on whether one likes nuclear power prevents useful judgement.

Carbon capture likewise names a family of processes, not a guarantee. A project needs to show what proportion is captured, what energy the process uses, what emissions remain elsewhere in the supply chain, and how storage will be monitored. A promised future capture system does not reduce today's emissions. A working system should not be dismissed merely because its industrial setting is unfashionable.

Hydrogen must be produced using energy and then stored or moved. Its usefulness depends on the application and the alternatives. Turning electricity into hydrogen and later back into electricity loses energy at each stage. That may still be worthwhile where storage or a particular industrial process requires it. The losses remain part of the calculation.

The attractive word \emph{natural} tells us little about a dam's effect on a river or a biofuel's effect on land use. The unattractive word \emph{industrial} tells us little about the emissions avoided by a reliable machine. What matters is the operation and its consequences. An allegiance to a technology can become as resistant to evidence as an allegiance against it.

\section*{The fuel someone else buys}

There is a further problem when a country's energy choices operate within world markets. If demand for a fuel falls in one place, the global price may fall and encourage some additional consumption elsewhere. The final reduction in consumption can therefore be smaller than the first buyer's reduction.

Hans-Werner Sinn developed a related argument about expectations. If resource owners anticipate much tighter future restrictions, they may have an incentive to extract and sell sooner. This is the mechanism associated with his ``green paradox.''\footnote{Hans-Werner Sinn, \href{https://www.ifo.de/en/publications/2008/working-paper/das-grune-paradoxon-warum-man-das-angebot-bei-der-klimapolitik}{\emph{Das grüne Paradoxon: Warum man das Angebot bei der Klimapolitik nicht vergessen darf}}, ifo Working Paper, 2008; \emph{The Green Paradox: A Supply-Side Approach to Global Warming}, MIT Press, 2012.} It is a warning about possible supply responses, not a law that every climate policy increases emissions.

How much of a demand reduction is offset depends on producers, consumers, substitutes, timing, and the scale of the change. Some resources may become uneconomic to extract. Cleaner substitutes can reduce demand persistently. Policy coordination can alter both supply and demand. It would be just as careless to assume complete offset as to assume none.

The practical implication is that an emissions claim needs to follow the response outside the jurisdiction announcing it. Suppose a country closes a steel plant and imports the same quantity of steel. Its territorial industrial emissions fall. The global result depends on where the replacement steel is made, how it is made, and the additional transport. A national success can coexist with little global gain or even a loss.

Conversely, a country producing equipment that lowers emissions elsewhere may contribute more than its territorial account alone reveals. Accounting systems have purposes, and territorial accounts are useful for assigning responsibility. They should not be mistaken for a complete description of the atmospheric result.

Border measures, product standards, and international agreements can attempt to address these differences. Each introduces its own problems of measurement, enforcement, trade, and distribution. The point is not to find a policy without complications. It is to stop congratulating ourselves before checking the complications that could defeat its purpose.

\section*{Who can afford to set an example?}

An expensive experiment may still be worthwhile if it develops a technology others can use more cheaply later. Early markets can support learning, manufacturing scale, and infrastructure. The benefit may extend beyond the first country's emissions reduction.

That argument requires a plausible account of learning. It cannot justify every subsidy indefinitely. What is expected to become cheaper? Through what process? How will progress be assessed? What would lead the government to change or stop the support? A programme that survives every failure because it represents leadership has acquired protection from the very evidence it was supposed to produce.

There is also a difference between demonstrating an attractive alternative and demonstrating a willingness to bear costs that others will not bear. A policy that works only where public budgets are exceptionally large may be difficult to imitate. A visibly disorderly transition can discourage adoption even if some technical components work well.

Leading by example is therefore an empirical claim about influence. It is not established by the leader's intention. The example must make a workable path more available or persuasive to others. If it produces a costly result that strengthens opponents elsewhere, the global effect may differ from the speech announcing it.

This is especially important when leaders ask ordinary households to finance the experiment. People can reasonably contribute to shared learning. They are entitled to know that this is what they are doing, what protection they have against excessive loss, and why the chosen experiment is worth the cost. Calling every objection selfish avoids that explanation.

\section*{The bill that arrives before the benefit}

Climate policy often asks for costs now in exchange for benefits spread across places and years. That can be justified, but the timing matters politically and materially. A household cannot pay today's electricity bill with its share of a future reduction in global damage.

Carbon pricing can encourage reductions where they cost less, while generating revenue that can be returned or used for other purposes. Yet the incidence depends on actual behaviour and available substitutes. A person with a short commute and good public transport has different options from a worker travelling between dispersed sites. A payment designed around an average household may leave some households under severe pressure.

Subsidies have distributional problems too. If a subsidy requires an initial purchase, it may favour people who already have capital, suitable property, or access to credit. Renters may finance benefits used chiefly by owners. These effects are not reasons to reject every subsidy; they are reasons to examine who can actually use it.

The design of assistance affects whether people experience the change as a feasible undertaking or an order issued by those insulated from its consequences. Where a policy creates a serious loss, admitting the loss is the beginning of a credible response. Telling the loser that the national total is favourable adds insult without adding resources.

An energy transition also needs enough political durability to continue through elections, price shocks, and disappointments. A policy that depends on everyone remaining enthusiastic is poorly designed for a long project. Clear rules, visible results, reliable services, and an intelligible distribution of costs matter because they help a programme survive periods when enthusiasm does not.

\section*{Keeping the lights on through the argument}

There will be no final energy arrangement that ends the need for judgement. Equipment ages. Demand changes. New technologies appear. Supply chains become vulnerable. Environmental effects that received little attention become clearer. A workable system requires continuing decisions and maintenance.

The aim should be to reduce damaging emissions while providing the services people need, with an honest account of cost and risk. This is demanding enough without requiring every component to symbolise virtue. We should be prepared to keep a useful machine, replace a familiar one, abandon a disappointing project, and adopt a technology associated with people we dislike.

There is a kind of melancholy in the loss of an industrial landscape or an occupation. A community built around a mine or plant does not simply transfer its memories to the next source of electricity. Those attachments deserve attention, although they cannot indefinitely determine what everyone else must buy or endure. Change may be necessary without being experienced as liberation by the people through whose lives it passes.

For them, the transition should mean more than an explanation that history has chosen another industry. It should include practical means to work, live, and remain connected to others. These obligations belong in the cost of the change, not in a compassionate paragraph appended after the budget is settled.

Meanwhile, the refrigerator must keep running. This prosaic requirement is a useful restraint on political imagination. It asks us to show how the new arrangement will work, hour by hour and year by year, while recognising that keeping the old one also has consequences. The public deserves a plan that can survive that question.

\chapter{Living with the Damage}

Imagine an older woman alone in a flat during a heatwave. The rooms remain hot through the night. She has been told to keep cool, but the advice contains no means of cooling the flat. The nearest public building with air conditioning is several streets away. Getting there would require an effort she is unsure she can manage.

Her immediate problem is not solved by winning an argument about the cause of the heatwave. Nor can a reduction in emissions today cool her bedroom tonight. She needs practical protection: a usable cool place, help reaching it, a check from someone who knows her circumstances, or changes to the building and its equipment.

There is something indecent about treating such measures as distractions from the real issue. The woman is part of the real issue. A politics concerned with future suffering must remain capable of noticing present suffering that does not wait for the preferred solution.

This is the starting point for adaptation. It is also the point at which a different evasion becomes possible: helping people endure the consequences can be used as an excuse to do nothing about their causes. We need to examine both, without forcing a person in danger to wait while we choose a side.

\section*{Protection is ordinary work}

People have always adapted to weather and climate. Roofs keep out rain. Clothing reduces exposure. Water is stored for periods when it is scarce. Crops and planting times reflect local conditions. None of this implies that people approve of droughts or regard storms as harmless.

Climate change makes some existing arrangements less adequate. A building suited to one range of summer conditions may overheat under another. A drainage system designed around earlier rainfall can be overwhelmed. A crop that once performed reliably may become vulnerable to different temperatures, pests, or water shortages. Adaptation asks what needs to change and who can make it change.

The answer is often institutionally dull. Someone must maintain a drain, inspect a bridge, keep a public building accessible, or identify residents who need assistance during heat. The work is less visible than announcing a long-term target. Its success may consist in something not happening: a room does not become dangerously hot, a road remains passable, a person does not become isolated.

WHO's guidance on heat--health action plans treats protection as a combination of warning, preparedness, reduced exposure, functioning health systems, and attention to people at greater risk. An alert becomes useful when it is connected to actions people and institutions can actually take.\footnote{WHO Regional Office for Europe, \href{https://www.who.int/europe/publications/i/item/9789289062930}{\emph{Heat--Health Action Plans: Guidance}, second edition}, 2026. The example opening this chapter is fictional and illustrates the gap between advice and the means to act upon it.}

Issuing advice without considering the recipient's options can be an efficient way for an authority to appear active. Someone living in a badly designed rented flat may have little power to change it. A worker paid by the task may be unable to stop during dangerous heat without losing essential income. The information matters, but so do the conditions that determine whether it can be used.

This is why adaptation cannot be reduced to individual prudence. People can take precautions, yet buildings, work rules, transport, public spaces, and health services are jointly organised. A warning system that assumes every resident has a car, money, and a person to call will fail some of the people most in need of it.

\section*{Cooling without an argument about character}

Air conditioning is an instructive case because it makes the tension visible. Cooling uses energy and may add emissions, depending on the electricity and equipment. It also protects people from heat. A blanket moral judgement on its use skips the questions that matter: how much cooling is needed, how efficiently it can be provided, and who lacks access.

Buildings can reduce heat exposure through shade, insulation suited to local conditions, ventilation when outside conditions permit, reflective materials, and design. Trees and public shade can help outdoors, although trees need suitable space, water, and time to grow. These measures can reduce demand for mechanical cooling. They will not remove the need for it in every building or climate.

A practical programme can improve efficiency and clean the electricity supply while extending protection. There is no virtue in making a vulnerable person endure dangerous heat to preserve an appearance of restraint. Nor is there sense in ignoring wasteful design because a machine can compensate for it at continuing cost.

The same principle applies to water. Reducing leaks may be cheaper and less damaging than finding another source. Recycling, storage, changes in use, and desalination each have a place under particular conditions. Each requires resources and introduces consequences. Naming a technology does not establish its suitability for a watershed.

A reservoir may alter an ecosystem and displace a community. Groundwater pumping can relieve an immediate shortage while depleting the supply on which later users depend. A cheap tariff can protect essential access while encouraging waste if poorly designed; a high price can restrain waste while excluding people from a necessity. The task is to distinguish uses and effects rather than assume that abundance or scarcity supplies its own policy.

\section*{When a defence changes the risk}

A flood wall can protect existing homes. It can also encourage more building behind it because the area now appears safe. If a larger flood eventually exceeds the defence, more people and property may be exposed than before. This does not establish that the wall was a mistake. It establishes that the change in behaviour belongs in the risk assessment.

A similar problem arises when a measure protects one place by increasing pressure elsewhere. Water diverted from one district has to go somewhere. A shoreline intervention can affect neighbouring stretches of coast. An adaptation financed by a wealthy community can leave a poorer neighbour with new burdens.

The IPCC uses the term \emph{maladaptation} for responses that produce adverse consequences, such as increasing vulnerability or shifting risks to other people or places. Its assessment also distinguishes limits that may be overcome through additional resources or institutional changes from limits for which no effective adaptation is available to meet the relevant objective.\footnote{IPCC, \emph{Climate Change 2022: Impacts, Adaptation and Vulnerability}, \href{https://www.ipcc.ch/report/ar6/wg2/chapter/technical-summary/?lang=en}{Technical Summary}, especially the discussions of adaptation limits and maladaptation.}

The terms are useful if they lead to specific questions. Who is safer after the project? Who is more exposed? What behaviour does the project encourage? What happens if it fails? Does the investment make later alternatives harder? A project can be locally successful and still deserve criticism for the costs it exports.

Maintenance belongs in this account. Building a defence without financing its upkeep is a way of postponing part of the decision. The structure's existence creates confidence before its long-term reliability has been secured. An elected official may leave office with a completed project while successors inherit an obligation that no longer attracts votes.

Adaptation is therefore not simply the application of engineering to an external hazard. It changes a social arrangement. Ownership, insurance, planning, and expectations can change with it. The protection must be assessed together with the life that grows around it.

\section*{The place that cannot be kept}

Some communities will confront a more painful choice: whether to continue defending a place or help people move. The phrase \emph{planned retreat} is administratively tidy. For a resident it may mean leaving a home, neighbours, graves, a business, and a view known since childhood.

Compensation can help someone obtain another house. It cannot reproduce the same life elsewhere. Even a generous payment does not make the loss imaginary. People asked to leave are entitled to a process that acknowledges this, provides reliable information, and offers choices before an emergency removes them.

There are difficult questions about who should pay. If owners bought with knowledge of a serious hazard, should others finance repeated reconstruction? If risks were concealed or public authorities encouraged development, how should that affect the claim? What about tenants who did not choose the location's long-term risk but have no capital with which to move? A policy concerned only with property values can overlook people whose main asset is their place in a community.

These questions cannot be settled by declaring every threatened home sacred. Defending one location can consume resources that would protect more people elsewhere. Nor can they be settled by describing residents as irrational because their attachments exceed the market price. A cost calculation can inform a decision without explaining everything the decision takes away.

There are also nonhuman losses that cannot be repaired by moving people. A habitat may disappear or change faster than its species can adjust. A seawall protecting a town may leave no room for a coastal habitat to move inland. Adaptation to protect one valued thing can make the loss of another more likely.

We should be suspicious of accounts that turn every such conflict into a future opportunity. Some adaptations are improvements. Others are the least damaging way to manage a loss. Calling the second an improvement because we dislike the word \emph{loss} adds no protection.

\section*{Why adaptation does not settle emissions}

The ability to protect people from some consequences is not a reason to impose avoidable additional consequences. A hospital can treat an injury; that does not make preventing the injury pointless. Likewise, the fact that a city can adapt to some warming does not establish that adaptation to indefinitely greater warming will remain affordable or effective.

The reverse inference is also false. The importance of reducing emissions does not remove the need to prepare for changes already occurring or risks that would exist even in a stable global climate. Some parts of the system respond over long periods. People will need protection while mitigation takes effect, and from hazards that mitigation does not eliminate.

The two approaches affect one another. Lower future warming can reduce the scale of adaptation required. Better adaptation can reduce immediate harm and preserve the capacity to act. Badly designed adaptation can increase emissions or fix people in vulnerable places. Badly designed mitigation can leave them unable to afford necessary protection. A policy needs to examine these interactions rather than assign all virtue to one side.

There is no universal budget ratio that settles the matter for every place. A city facing repeated present flooding may have urgent protective work. A national government also influences emissions through infrastructure and energy policy. A poorer country may require external support for both. The allocation depends on hazards, capabilities, costs, and obligations.

People who describe adaptation as surrender can overlook the present victim. People who describe mitigation as unnecessary can assign a growing burden of adaptation to later victims. Both can appear practical while placing the most difficult part of their argument outside the room.

\section*{Can we remove what we have added?}

Carbon dioxide removal attempts to take carbon dioxide from the atmosphere and store it. The broad category includes biological approaches and engineered ones. It is distinct from capturing emissions at an industrial source before they enter the atmosphere.

The accounting question is whether carbon has been removed, how much remains removed, and for how long. Growing a forest can store carbon, but fire, disease, harvesting, or later land use can release it. Engineered removal requires energy, equipment, storage, and verification. The costs and permanence differ across methods.

The land also has other uses. A large removal programme may compete with food production, habitat, or the people who already live there. A model that allocates land to carbon removal has not obtained the residents' agreement or established that the ecological consequences are acceptable.

None of this makes removal inherently fraudulent. It makes the difference between a credible activity and a convenient promise important. Future removal can become an excuse for continued emissions if its feasibility is assumed rather than established. A company should not receive the same recognition for a speculative future tonne as for a verified and durable removal already achieved.

There are legitimate reasons to develop the capacity, including dealing with residual emissions that are difficult to eliminate and potentially reducing atmospheric carbon after an overshoot. Those reasons are stronger when the programme is treated as a demanding material undertaking. It is least credible when presented as permission to postpone every difficult change.

\section*{The temptation to turn down the sunlight}

Other proposals would alter the amount of sunlight absorbed by the Earth. Some seek to increase reflection by introducing particles into the stratosphere or changing cloud properties. The intended effect is cooling. These approaches do not remove atmospheric carbon dioxide and do not solve all the effects of its accumulation, including ocean acidification.

Their possible advantages and risks are different from those of emissions reduction or carbon removal. A change in the global average could occur alongside uneven regional effects. Precipitation could change. A programme maintained for years would raise questions about continuity, interruption, and the consequences of stopping while greenhouse-gas concentrations remained high.

The political difficulty is at least as serious as the technical one. Who decides the desired temperature? Who accepts responsibility if one region attributes a drought to the intervention? What kind of agreement would bind states that expect different benefits and harms? An action described as protecting the planet could become a profound exercise of power over people who had no effective say in it.

In 2021, the US National Academies recommended a governed research programme to improve understanding of solar geoengineering. It did not treat a recommendation to investigate as permission to deploy. The report put research governance alongside physical research because the questions cannot responsibly be separated.\footnote{National Academies of Sciences, Engineering, and Medicine, \href{https://www.nationalacademies.org/read/25762}{\emph{Reflecting Sunlight: Recommendations for Solar Geoengineering Research and Research Governance}}, 2021, doi:10.17226/25762.}

It would be careless to treat all research as though it were deployment. Modelling, observation, laboratory work, and an outdoor intervention have different implications. It would also be careless to assume that research is politically innocent. Funding can create constituencies, institutions, and expectations that favour further development. The scale and rules of a programme matter.

A blanket refusal to learn can leave decisions to the less scrupulous or to an emergency in which ignorance is especially dangerous. An unrestricted research programme can create momentum towards an intervention before legitimate authority exists to make it. The task is to govern inquiry in ways that make its results available without pretending that knowledge automatically grants permission.

We may eventually judge a proposed intervention too risky. We may judge some limited research useful. These are decisions about identifiable activities. Describing the whole field as salvation or hubris is easier, but it does not tell us which experiment should proceed or who should supervise it.

\section*{What success would leave unfinished}

A successful response to climate change would still leave losses. Some would result from changes already under way; others from the measures adopted in response. A landscape would change, an occupation would contract, a community might move, a species might not survive. No honest programme should promise that sufficient commitment can make every valued thing compatible with every other.

This recognition can produce bitterness. People may ask why they should act if damage cannot be entirely prevented. The answer is that differences in damage remain differences in lives. Protecting a person from a particular harm does not become pointless because another harm remains. The older woman in the overheated flat does not require a theory of final planetary success before benefiting from a cooler room.

There is also relief in giving up the demand that every measure redeem the whole situation. A useful project can be useful within limits. A warning can be acted upon without turning every day into a state of alarm. A mistake can be corrected without declaring that all earlier concern was fraudulent.

None of this permits indifference to what cannot be repaired. It permits a form of attention less absorbed in announcing its own seriousness. There is work to be done on a building, a water system, a warning network, a source of electricity. There are people who need help making an unavoidable move. The success of that work may never look like the restoration of the world we remember. It can still leave people with more room to live in the world that remains.

\chapter{A Country Changes}

Germany recorded 654,241 births in 2025. The total fertility rate was 1.32 children per woman. These are sober figures, but public discussion rarely leaves them sober for long. They become evidence that the country is disappearing, that young people have become selfish, that immigration is indispensable, or that immigration has already destroyed what it was meant to preserve.\footnote{Federal Statistical Office of Germany (Destatis), \href{https://www.destatis.de/EN/Themes/Society-Environment/Population/Births/_node.html}{\emph{Births}}, final results for 2025, released in 2026. The total fertility rate is a period measure, distinct from completed cohort fertility.}

Before accepting any of these conclusions, we should understand what the numbers say. The birth count is a count of births during a year. The total fertility rate combines the birth rates observed at different maternal ages during that year. It estimates how many children a woman would have if she experienced those age-specific rates throughout her reproductive life. It is not a completed biography of the women alive in 2025.

Timing matters. If births are postponed, the annual rate can fall even when some postponed births later occur. If fewer women are in the ages at which births are common, the total number of babies can fall without an equivalent change in individual fertility. Neither complication makes a persistently low rate unimportant. It prevents us from claiming more precision than the measure provides.

Demography is often frightening because its arithmetic is real and its consequences arrive slowly. A country cannot create twenty-five-year-old workers next year by encouraging births this year. Yet the same slowness makes demographic numbers attractive to prophets: the eventual outcome lies far enough away for a confident speaker to avoid an early reckoning.

\section*{People already born}

Some demographic changes are more predictable than others. A large cohort approaching retirement is already alive. Its approximate size is known. Deaths, migration, health, and retirement choices will affect the outcome, but the cohort cannot disappear from the calculation because a government would prefer a younger population.

Germany's sixteenth coordinated population projection, published in December 2025, found across its variants that roughly a quarter of the population would be at least sixty-seven by 2035.\footnote{Destatis, \href{https://www.destatis.de/EN/Press/2025/12/PE25_446_12.html}{\emph{By 2035, One Quarter of Germany's Population Will Be Aged 67 or Over}}, 11 December 2025, presenting the sixteenth coordinated population projection. Projection results depend on their stated assumptions.} The convergence of different variants over the near term is part of the point. Some of the change has already been built into the age structure.

This creates practical obligations. Pensions must be financed. Health and care services need workers and facilities. Housing, transport, and ordinary services need to remain usable as more residents become old. These needs do not become manageable because the people reaching retirement worked hard, nor do they become undeserved because younger people face a difficult bill.

Long life is itself an achievement. To describe older people only as a dependency burden is to forget what the earlier generations hoped medicine and prosperity would permit. The difficulty lies in arranging a society that can sustain the achievement. A longer life does not automatically come with additional years of good health, adequate savings, or a person available to provide care.

The ratio between age groups is useful, but it is not the complete economic ratio. Some people of working age do not work; some older people do. Hours, productivity, wages, health, and the distribution of work all matter. A younger population can be poor and underemployed. An older population can be productive and well organised. Age structure constrains the possibilities without specifying the entire result.

\section*{A smaller generation}

Persistently low fertility changes the size of future generations. To see the direction, imagine an artificial population in which each generation produces only seventy per cent as many future parents as the preceding one, with no migration and no other change. One hundred becomes seventy, then forty-nine, then about thirty-four. The compounding is real within the example.

It is not a forecast for a modern country. Generations overlap. Mortality changes. Migration occurs. Fertility differs across time and groups. The age distribution can produce population momentum. A simple sequence teaches what a sustained replacement shortfall can do; it does not tell us that the assumptions will remain fixed for a century.

This distinction matters when a calculation is used to announce national extinction. A rate observed today becomes permanent in the argument, every group is assumed to retain its current behaviour, and categories are carried unchanged through future marriages and generations. The arithmetic may be correct while the description of people is implausibly frozen.

There is still a serious problem to address without the extinction story. Fewer children today can mean fewer entrants to the labour market later. School networks may contract in one region while care needs expand. Infrastructure built for a larger population can become expensive per resident. A smaller tax base may have to maintain roads, water systems, and public buildings spread across the same area.

Contraction can bring benefits too, such as reduced pressure on some land and housing. But these benefits do not necessarily occur where people need them. An empty house in a village with few jobs is not a cheap home near a thriving city's workplaces. National totals conceal the geography in which lives must be arranged.

\section*{Why the state cannot order a future}

A government concerned about births can reduce obstacles to family life. Affordable housing, dependable childcare, predictable employment, and a workable relationship between paid work and parenting can make desired children easier to have. These measures can be justified by the lives they improve even when their eventual effect on fertility is uncertain.

The government cannot assume that removing an obstacle will produce a particular number of babies. People differ in what they want. Some do not want children. Others want them but do not find a partner, cannot conceive, become ill, or revise their plans. A national target encounters private relationships that cannot be managed like orders to a factory.

The temptation to blame women is especially revealing. A society may encourage education and employment, make housing costly, organise careers around uninterrupted availability, and then complain that women have postponed motherhood. It has arranged conflicting demands and assigned the resulting tension to one sex as a defect of character.

Men's decisions, work patterns, willingness to care, and reliability as partners belong in the account. So do institutions that assume someone at home will absorb every interruption. A nursery that closes unpredictably is not merely an inconvenience if a parent's job depends on being present elsewhere. A promised right to childcare is insufficient when the place does not exist.

None of this establishes that a better policy can return fertility to replacement level. It identifies avoidable impediments and distributes responsibility more honestly. A state can make family life less difficult without claiming ownership of the result.

Children should also be wanted for more than their future contributions. To speak of them chiefly as pension financiers is to adopt a strangely impoverished account of the future one hopes to preserve. The point of sustaining a society is that people can live in it, not that it can manufacture enough people to service its accounts.

\section*{What migration changes}

Migration affects population size and age structure more quickly than births. In 2025, Germany had about 235,000 more arrivals than departures across its borders. That net figure came from much larger flows in both directions.\footnote{Destatis, \href{https://www.destatis.de/DE/Presse/Pressemitteilungen/2026/06/PD26_184_12411.html}{\emph{Nettozuwanderung 2025 mit 235 000 Personen deutlich gesunken}}, June 2026: approximately 1.48 million arrivals and 1.25 million departures across Germany's borders.} People leave as well as arrive, including citizens and non-citizens. A country does not simply add each year's arrivals to a permanent stock.

New residents can work, establish businesses, raise children, and take part in public life. They can also need housing, education, language instruction, care, and income support. Which effects dominate depends on who arrives, when, where, and under what conditions. A national head count cannot supply the answer.

Migrants also age. This does not make their contribution temporary in a dismissive sense; every person's working life is temporary. It does mean that immigration cannot abolish ageing through a single numerical addition. The long-term arrangement still needs to finance childhood, illness, unemployment, and retirement.

Nor are migrants a standard product that can be ordered in whatever quantity the pension system requires. A receiving country competes with other destinations, depends on people's willingness to remain, and encounters their independent plans. Family life, qualifications, language, and the conditions of work affect those decisions.

The sending country belongs in the account as well. Recruiting a trained nurse can help a hospital here while making staffing harder elsewhere. Remittances and opportunities can benefit the place of origin, but the effects should be examined rather than assumed to cancel. The shortage does not become morally irrelevant when it crosses the border in the other direction.

Migration policy therefore cannot be reduced to a demand for more or fewer people. It needs to distinguish purposes and capacities. That is the subject of the next chapter. Here the demographic point is enough: migration can alter a trajectory without magically completing the institutions needed to make that alteration work.

\section*{Continuity is more than a total}

A resident worried about change may not be worrying about population size at all. The shops have different names. The language heard in the street has changed. A familiar association has closed. A school teaches children whose parents arrived from many places. The resident experiences a loss of familiarity that a favourable employment statistic does not answer.

It is patronising to insist that such a person must welcome every change. Places matter through habits, shared references, and accumulated memory. Someone can miss these without believing that strangers have fewer human rights. The feeling should be heard accurately before being classified as either wisdom or hatred.

But a feeling of dispossession does not establish who caused every change or what should be done about it. The old shop may have closed because children of its owners moved away, customers bought online, rent rose, or the owner retired. An immigrant business may occupy a premises that would otherwise stand empty. Several processes can occur together, and the most visible newcomer may receive blame for all of them.

Nor does familiarity confer unlimited authority over the future inhabitants of a place. People born there grow up with different preferences. They marry people their parents would not have chosen, abandon religious practices, change occupations, and move. A society preserves itself through people who do not simply repeat their predecessors.

The difficult question is which forms of continuity deserve public protection. A common language for schools and public institutions serves a practical purpose. Equal legal standing, security, reliable administration, and the freedom to dissent are not merely decorative customs. They can be required and defended without demanding identical food, dress, ancestry, or private belief.

Some disagreements genuinely concern the public order. A demand that religious authority override civil law, or that women accept restrictions inconsistent with equal rights, cannot be dismissed as an innocent difference in taste. Equally, foreign parentage does not tell us that a person holds such a demand. We need to identify conduct and commitments rather than let ancestry perform the judgement.

\section*{Categories that outlive their usefulness}

Terms such as \emph{foreign citizen}, \emph{foreign-born}, \emph{immigrant}, and \emph{person with an immigration background} answer different questions. A person can be a citizen and have been born abroad. A child can be born locally and have foreign-born parents. Naturalisation changes citizenship statistics without moving anyone. A definition based on parental birthplace can classify people whose daily lives and language are entirely local.

For some analyses, these distinctions are indispensable. They can help identify barriers in education or employment and assess whether those barriers diminish across generations. But a category designed to reveal a disadvantage can become misleading when treated as a permanent description of cultural allegiance.

If a child's parentage is counted as evidence that the country is losing itself regardless of the child's language, conduct, or citizenship, the argument has made belonging impossible in advance. Integration can never succeed under a definition that refuses to count its success. Conversely, if every statistical difference is suppressed for fear of encouraging prejudice, institutions lose information needed to address persistent failure.

The right category depends on the question. A school needs to know a child's command of the teaching language. A labour-market study needs qualifications and work history. An electoral analysis needs citizenship and participation. None gains precision by using a single broad label for everything.

Careful categories can also make criticism sharper. Saying that a particular programme leaves a specified group without adequate language after several years is more actionable than declaring that integration has failed everywhere. The specific claim identifies a result that authorities can be required to explain.

\section*{The temptation of irreversible decline}

There are changes that cannot be undone. The childhoods already spent in inadequate schools cannot be supplied again. A closed institution may be difficult to rebuild. People who left a declining region may never return. Acknowledging these losses is part of taking demography seriously.

It does not establish that all future outcomes have been decided. The claim that a country has passed a point of no return is often left without a definition of what return means. A return to which year, population, economy, or social arrangement? If the standard is literal restoration of a past society, every country is beyond return. That says little about whether its schools, housing, employment, and public order can improve.

Fatalism can be gratifying because it relieves the speaker of proposing a workable response. Every improvement is dismissed as temporary; every failure becomes confirmation. The person who claims to face reality may have constructed a belief that reality can no longer challenge.

Official optimism can be equally evasive. A government points to successful people while refusing to discuss patterns of dependence or institutional overload. It calls criticism a failure of attitude and expects those living with the consequences to adopt a better one. Both the optimist and the prophet of decline can avoid the particular work that would make a difference.

A more demanding account asks which outcomes can still change. Employment participation can change. So can the quality of training, the time needed to recognise qualifications, housing supply, the organisation of care, and the rules of admission. Some reforms will disappoint. Some will improve matters without restoring what was lost. Their value should be judged against feasible alternatives, not a vanished country made perfect by memory.

\section*{A future with fewer consolations}

Ageing societies will face choices that cannot all be pleasant. Working longer may be feasible for some occupations and punishing in others. Higher contributions burden workers; lower benefits burden retirees. More productivity can help, but it cannot be ordered by optimism. Care still involves time with particular people, much of it resistant to the gains possible in manufacturing.

Migration can help meet some needs while creating demands that must be funded. Better conditions for families can improve life without producing the births a projection assumes. Regional contraction can be managed with more or less care, but some institutions will close and some places will become quieter.

It is possible to state all this without describing older people as an affliction, children as economic inputs, or newcomers as interchangeable substitutes. Demography counts people so that consequences can be understood. The counting should not make the lives disappear.

The country will not remain the country of anyone's childhood. That loss is universal enough to deserve less theatrical handling than it receives. What remains open is whether change is governed competently, whether obligations are made explicit, and whether people can still build lives that they recognise as their own. These questions are difficult. Unlike a declaration that everything is already over, they leave something to answer for.

\chapter{People Are Not Interchangeable Numbers}

A hospital recruits a nurse from abroad. A university admits an international student. A family is reunited after years apart. A person asks for protection from persecution. Someone else arrives hoping to find work without the required permission. All may appear in a public discussion under the word \emph{immigration}.

The word is broad enough to describe a movement and too broad to settle the questions arising from it. The hospital needs a qualification recognised and a person able to communicate safely with patients. The student needs a place to study and somewhere to live. The protection claim needs a fair decision. The family raises obligations that are not identical to an employer's need for labour.

An argument that treats these people as equivalent will predict badly and govern badly. The error occurs whether they are presented as the answer to every shortage or as a uniform burden. Both accounts replace distinct lives and legal circumstances with a convenient total.

We can begin more usefully by asking why admission is being considered, what it permits, and what institutions must be able to deliver afterwards.

\section*{The reason for admission matters}

Labour migration and humanitarian protection have different purposes. A labour route may select for skills, a job, earnings, or other attributes expected to benefit the receiving economy. Protection addresses danger and the inability to obtain effective protection at home. A person can have both valuable skills and a protection need, but neither establishes the other.

The Refugee Convention gives a specific meaning to refugee status and sets out obligations, including the central protection against return to persecution. Other forms of protection and human-rights obligations also matter. A fair procedure is needed precisely because a claim and an established status are not the same thing.\footnote{UNHCR, \href{https://www.unhcr.org/africa/news/stories/refugees-and-migrants-frequently-asked-questions-faqs}{\emph{``Refugees'' and ``Migrants'': Frequently Asked Questions}}, and the 1951 Refugee Convention and 1967 Protocol.}

It is evasive to defend humanitarian admissions by promising that every person will quickly become an economic asset. That subjects protection to a test it was not designed to meet and creates an expectation that may be false. A person who needs refuge may be old, ill, traumatised, or poorly educated. The case for protecting that person should be stated honestly, including the likely cost.

It is equally evasive to treat every wish for a better life as a protection claim. The wish can be entirely understandable without creating the same legal entitlement. Maintaining the distinction is part of preserving a system in which protection can be decided credibly.

These distinctions can be difficult in individual cases. Economic deprivation may be connected with persecution or conflict. Documents may be missing. Accounts may be incomplete, mistaken, or deliberately false. Difficulty is a reason to provide competent assessment, not to decide every case according to the most comforting general story.

A state that leaves people waiting for years produces harms of its own. Genuine claims remain unresolved; weak claims acquire time without a decision; people plan lives around uncertain expectations; and public confidence deteriorates. Speed without fairness can make grave errors. Fairness without workable timeliness can become a promise that the procedure itself defeats.

\section*{Capacity after the border}

Admission is not completed when a person crosses a border. Somewhere there must be a bed, a school place where appropriate, access to necessary services, and a way through the legal and administrative process. These resources are located in particular municipalities. They are not produced by dividing a national land area by the number of arrivals.

An empty building may lack transport, staff, or work nearby. A vacant hospital bed is not usable without the people needed to operate the ward. A school classroom is more than a room. National discussions of capacity often count the visible object and omit the service that makes it useful.

The pace matters as well as the total. A town may accommodate a given number over several years and struggle with the same number arriving in weeks. Language teaching, housing, and administrative decisions cannot always expand at the speed of an emergency. A national policy that ignores these differences transfers its most difficult obligations to local staff and residents.

There are cases in which a sudden humanitarian emergency warrants accepting severe temporary pressure. The pressure should then be acknowledged and shared. It should not be declared imaginary to preserve the favourable appearance of the decision. The teacher facing an unmanageable class has not refuted the case for refuge; she has encountered an unfunded part of it.

Conversely, claiming that capacity is fixed can conceal a refusal to build it. A wealthy jurisdiction may insist that it has no room while choosing other uses for money and land. The claim should be tested concretely. Which resource is scarce? How fast can it be expanded? At what cost? What would be displaced? A serious limit can survive these questions. A slogan about limits often cannot.

\section*{Work takes time, but time is not enough}

Employment is one of the most important outcomes because it provides income, contact, language practice, and a way to contribute. Yet the relevant comparison depends on the person's starting point. A recruited engineer entering a specified job is not comparable to a newly arrived refugee awaiting a decision and beginning to learn German.

An IAB study published in 2024 followed refugees who arrived in Germany between 2013 and 2019. Employment rose with time in the country, reaching sixty-eight per cent after eight years. For the 2015 arrival cohort observed in 2022, employment was seventy-five per cent among men and thirty-one per cent among women.\footnote{IAB, \href{https://doku.iab.de/kurzber/2024/kb2024-10_en.pdf}{\emph{Labor Market Integration of Refugees: Improved Institutional Settings Promote Employment}}, IAB Brief Report 10/2024, using the IAB--BAMF--SOEP Survey of Refugees. The eight-year estimate and the 2015-cohort comparison refer to the groups and observation periods specified in the study.}

The first finding contradicts the claim that refugees remain permanently outside work. The second reveals a large gap that an overall average can conceal. Neither finding tells us that every employed person earns enough to support a household without assistance, works full-time, or uses their previous qualification. Employment is a necessary part of the account, not the whole account.

The progression also raises questions about the years before work begins. Some time is needed for language, training, recovery, and finding a position. Some delay is avoidable. Procedures can leave a person unable to work, unable to move where work exists, or unable to have a qualification assessed promptly. An institution can help create the dependence later used as evidence that the person lacks initiative.

But administrative obstacles do not explain every persistent difficulty. Education, literacy, health, care responsibilities, expectations within a household, and the availability of suitable jobs matter. A programme that assumes the same barrier for everyone can spend heavily without addressing the reason a particular person remains outside work.

Time alone is not a policy. What happens during it matters. A language course that someone can attend, childcare that fits its hours, contact with employers, and recognition of usable skills can change the trajectory. An indefinite sequence of appointments without a clear route may simply make the waiting more bureaucratically elaborate.

\section*{Women hidden in the average}

The difference between men's and women's employment deserves direct attention. It may reflect children and care, previous schooling, language access, health, or norms about whether women should work outside the home. These explanations can overlap. We should not choose the most flattering one before investigating.

A woman who wants employment may be constrained by a partner or family. If the state treats the household as a single satisfied unit because one man works, it may overlook her dependence. Services should be accessible to her as a person with her own rights and plans. Information reaching the household's most powerful member is not necessarily information available to everyone in it.

At the same time, a government cannot simply announce employment as a requirement while failing to provide childcare or suitable training. A demand that would be difficult for any parent does not become feasible because the parent is foreign-born. Clear expectations require workable means.

There is no need to disguise coercive family restrictions as cultural richness. There is also no basis for assuming that every woman from a particular background accepts those restrictions or lives under them. The practical task is to make independent participation possible and to respond when another person prevents it unlawfully.

Success here may create conflict within a family. A woman who learns the language, earns money, or exercises an independent choice can unsettle expectations. Integration should not be measured only by the absence of visible conflict. Sometimes conflict is evidence that a previously subordinate person has acquired room to act.

\section*{The worker on the night shift}

Foreign citizens made up about one fifth of people employed in nursing occupations in Germany in 2025, according to the Federal Employment Agency. Since 2022, employment growth in that field had come from foreign citizens.\footnote{Federal Employment Agency, \href{https://statistik.arbeitsagentur.de/DE/Statischer-Content/Statistiken/Themen-im-Fokus/Berufe/Online-KV/Pflege.html}{\emph{Arbeitsmarktsituation in Pflegeberufen}}, data for 2025. Citizenship is the category used in these employment figures.} These are people participating in a service that an ageing society needs.

The fact does not establish that every form of migration solves a nursing shortage. A person without the relevant training is not a nurse because a hospital has a vacancy. Qualifications, language, supervision, and working conditions matter. Recruitment cannot be used to avoid improving conditions that cause existing workers to leave.

Still, anyone describing migrants only as additional patients has omitted the staff member helping those patients through the night. The omission is particularly ugly when dependence on that worker is combined with an insistence that she can never fully belong.

Employers can also use an attractive story about labour shortages to obtain workers under conditions that deserve criticism. A shortage at a particular wage is not identical to an absolute absence of people capable of doing the work. Recruitment may be justified; it should not grant an employer permission to suppress wages, tie a person's security to compliance, or neglect training.

The migrant's bargaining position matters. If losing a job threatens residence, housing, or the repayment of recruitment costs, refusing abuse may be difficult. A policy designed to meet labour demand needs enforcement that protects workers rather than merely delivering them to employers.

The interests of the receiving society and the migrant can align. They do not align automatically in every contract. Treating migrants as economic units makes it easy to celebrate their contribution while ignoring the conditions under which it is obtained.

\section*{What the fiscal account includes}

Public debate often asks whether immigrants pay more in taxes than they receive in benefits. The question is legitimate, but an answer depends on what is counted and over what period. A child usually receives services before paying taxes. An adult's contribution changes with employment, earnings, family circumstances, health, and age. A snapshot and a lifetime estimate describe different things.

Assumptions about future employment and wages can dominate a long-term calculation. So can decisions about how to allocate the cost of public services used by everyone. A new resident may not require another unit of every public good, but rapid population growth can require expensive additions to housing, schools, transport, or administration. The marginal cost need not equal the average cost.

There are economic effects outside the government's accounts: production, prices, rents, business formation, and changes in opportunities for existing workers. A positive fiscal estimate does not prove that every resident benefits. A negative estimate does not tell us whether a humanitarian obligation should be abandoned. It tells us something about a cost that ought to be faced openly.

The same standards should apply to flattering and unfavourable studies. Which group is described? Which year or cohort? Are students, recruited workers, refugees, and dependants combined? Are results adjusted for age, and is that adjustment suitable for the question? What future behaviour is assumed?

The public deserves these distinctions because migration is not a one-time purchase with a single price. It is an alteration of a society's membership and obligations. Pretending the account is simple helps advocates on both sides and leaves administrators with the complexity.

\section*{Reciprocity without humiliation}

A welfare system can require reasonable participation from people able to participate. Learning the language, attending an agreed programme, seeking suitable work, and providing accurate information can be legitimate conditions attached to particular forms of support. Disability, illness, care responsibilities, and other relevant circumstances must be assessed rather than treated as convenient exceptions or universal excuses.

The point should be to sustain a workable relationship between support and contribution. It should not be to make a person perform gratitude. Public benefits are administered under rules; an official should not be free to demand deference beyond them. Equally, receiving support does not make accurate reporting optional or fraud a harmless response to need.

Clear obligations can protect public consent for generosity. Arbitrary demands and theatrical punishments can destroy it. An administration that alternates prolonged inaction with a sudden campaign against a few visible cases is neither reliably firm nor reliably humane.

The rules should be comprehensible and contestable. A person should know what is expected, what help is available, what consequence follows a failure, and how an error can be challenged. This is ordinary administrative competence. Its absence turns every case into a struggle over discretion.

Residents already contributing can reasonably object when officials dismiss every concern about abuse or capacity as prejudice. They can also be misled by claims that exceptional abuse describes the whole system. The antidote is credible reporting and enforcement, not a demand that everyone either trust completely or reject completely.

\section*{A decision that is carried out}

Migration systems lose credibility when the formal outcome of a procedure has little relationship to what happens afterwards. Recognition of protection should lead to the security and opportunities associated with that status. A final refusal should have a lawful consequence, subject to applicable safeguards and obstacles to return. Long-term limbo is not a neutral compromise.

Return can be practically and legally difficult. Identity, documentation, cooperation by other states, family circumstances, health, and risks on return can matter. A government should explain these obstacles rather than announce an effortless solution. It should also distinguish obstacles it cannot presently overcome from failures of its own administration.

No person should be returned to persecution merely to demonstrate control. Nor should a government treat a procedure as meaningful while quietly assuming that its decisions need never be implemented. Either position corrodes the distinction between admission under a rule and remaining because the rule has ceased to operate.

The public argument often rewards the harsh announcement and ignores execution. Actual administration requires staff, agreements, evidence, appeals, and time. If these are absent, repeated declarations of firmness become another variety of the symbolic politics they claim to despise.

\section*{Belonging after usefulness}

A country can select some entrants for their likely contribution. It should be careful about allowing the language of selection to govern every later relationship. People become ill, lose jobs, change occupations, have children, and grow old. A society that values a person only while economically convenient has offered a precarious kind of welcome.

Citizenship and secure membership involve more than a favourable annual balance. They create rights and obligations that cannot reasonably be recalculated whenever a person becomes costly. This is one reason admission and the path to membership require clear rules: they concern relationships intended to endure changes in circumstance.

The receiving society also deserves candour about what is being undertaken. A temporary labour arrangement can become permanent settlement. A protection situation can last much longer than expected. Families form. Children grow up locally. Refusing to recognise these developments does not prevent them; it makes later decisions more painful and less honest.

There is room for a policy that is selective about some admissions, serious about protection, demanding about lawful conduct, and committed to the people who become members. These are not contradictory intentions. They become difficult when capacity is denied, rules are left unclear, and human beings are used as evidence for a story chosen before they arrived.

The nurse, the student, the reunited family, and the person seeking refuge will not have identical futures. A government cannot promise that all will prosper. It can stop treating their differences as an inconvenience to its rhetoric and build a system capable of responding to the lives actually before it.

\chapter{What Living Together Requires}

A child cannot follow a lesson in a language she does not yet understand. There is nothing obscure about this. She may be intelligent, curious, and eager to learn. None of these qualities supplies the missing words at the moment the teacher explains a fraction or asks her to read a passage.

The school can help her learn the language. That takes time and appropriate instruction. If many children need such help and the school lacks staff, the problem affects the whole class. Denying the difficulty does not protect the newcomers. It leaves them sitting through lessons that pass them by and exposes them to the resentment of classmates whose own instruction suffers.

This is a useful place to begin an argument about integration because it removes much of the theatrical language. The question is not whether we approve of diversity. It is whether the children can understand the lesson, what they need if they cannot, and who will provide it.

\section*{A common language}

German in a German public school is more than a marker of belonging. It is a means of access to instruction, institutions, work, and disagreement. A person unable to explain a problem to an official is dependent on whoever translates it. A worker unable to understand a safety instruction faces a practical danger. A parent unable to speak with a teacher has less power to help a child.

Multilingualism can be an advantage. It does not abolish the need for a shared language in a particular setting. The question is not whether a child should forget a family language. It is whether the child acquires the language needed to participate fully in the institutions around her.

Schools should therefore measure actual competence and teach accordingly. An ancestry label is an inadequate substitute. Some children with immigrant parents speak excellent German; some long-established families provide little support for reading. A diagnostic assessment can identify a need that a social category misses.

The IQB's 2022 educational assessment documented substantial differences in German-language competencies associated with immigration background, alongside the importance of social and family circumstances.\footnote{Petra Stanat and colleagues, eds., \href{https://www.iqb.hu-berlin.de/media/documents/IQB_Bildungstrend2022_Berichtsband.pdf}{\emph{IQB-Bildungstrend 2022}}, 2023, especially the analysis of immigration-related and social disparities. These are group-level patterns in assessed competencies, not measures of a child's fixed potential.} These patterns justify attention and resources. They do not determine what any particular child can achieve.

Lowering expectations to make the results look more equal is a betrayal disguised as kindness. A certificate that conceals inadequate reading does not make the next stage easier. It moves the discovery to an employer, a training programme, or the young person's own attempt to manage adult life.

High expectations also require teaching capable of meeting the need. A school cannot be made effective by being instructed to stop making excuses while its staffing and time remain inadequate. The standard and the means belong together. Otherwise either the certificate or the demand becomes a fiction.

\section*{The classroom and the official account}

Schools can become places where adults protect a public story at children's expense. A teacher reports that language support is insufficient and is told to avoid negative descriptions. A parent reports intimidation and is advised not to inflame tensions. An administrator praises inclusion while children withdraw from activities in which they no longer feel safe.

These are possible failures, not a description of every school. They matter because the incentives are recognisable. Reporting a problem can threaten a school's reputation before anyone supplies the resources to solve it. Staff may learn to make the report less alarming. The official picture improves while the children's day remains unchanged.

A good system makes a specific difficulty reportable. It records language needs, attendance, attainment, and serious incidents in ways that permit action. It does not treat all variation as proof of discrimination or all disadvantage as proof of incapacity. It asks what instruction and conditions would improve the result, then checks whether they do.

Parents also have obligations. Attendance, sleep, reading, communication with the school, and support for lawful rules cannot all be supplied by teachers. A parent who rejects a female teacher's authority because she is a woman has created a conflict the school should address directly. Describing it only as a communication issue can hide the demand being made.

The child should not have to settle this conflict between adults. She needs the school to remain a place where the teacher's authority rests on the role and the law, and where her own opportunities do not depend on which restrictions her family would prefer to impose.

\section*{Safety is not an attitude}

The same need for precision applies to crime. People are entitled to public safety whether they are long-established residents or recent arrivals. An assault does not become less serious because discussing the offender's background would be politically inconvenient. Nor does an offender's background make an innocent person from the same group responsible.

Police data in Germany show differing suspect rates and overrepresentation of some non-German and migration-related categories in particular offences. The categories and limitations must be read with the figures. A suspect is not a convicted offender; police records concern offences known to police; citizenship is not birthplace or ethnicity. Non-residents can appear among non-German suspects, and offences specific to immigration law cannot be treated as though citizens have the same opportunity to commit them.\footnote{Federal Criminal Police Office (BKA), \emph{Kriminalität im Kontext von Zuwanderung: Fokus Fluchtmigration}, Bundeslagebild 2023, and \href{https://www.bka.de/DE/AktuelleInformationen/StatistikenLagebilder/PolizeilicheKriminalstatistik/PKS2024/PKSTabellen/BundTVNationalitaet/bundTVNationalitaet.html}{Police Crime Statistics 2024 tables by citizenship and residence status}. The reports' definitions and methodological cautions are integral to interpreting their figures.}

Age, sex, social circumstances, residence, and the composition of particular migration flows matter when explaining differences. These are not magical adjustments that make violence disappear. They help identify what comparison is meaningful and what interventions might reduce it.

Suppose one group contains many more young men than another. A difference in overall suspect rates may partly reflect that composition. If a receiving policy changes the number of young men in a locality, the resulting demand on policing and prevention can still be relevant. Explaining part of a rate does not establish that the change has no practical consequence.

Equally, a raw difference does not establish that origin itself is the cause. Deprivation, exposure to violence, peer groups, opportunity, reporting, and enforcement can contribute. A policy based on a false causal account may be harsh and ineffective at once. A useful explanation should help us distinguish what can be changed from what merely identifies the people in the table.

There is no conflict between prosecuting an offender and studying the conditions that make offending more likely. The victim needs justice and protection. Future victims benefit if preventable causes are addressed. Treating explanation as excuse can prevent learning; treating explanation as a reason to withhold accountability can abandon the victim.

\section*{Who is missing from the crime story?}

Migrants also suffer crime, including violence within their own communities. A public argument that counts them only as possible offenders can overlook people who need the state most urgently. A woman threatened by relatives, a shopkeeper extorted by a criminal network, or a young person intimidated for leaving a religion cannot be protected by flattering accounts of communal harmony.

Nor are they protected by general hostility towards the community in which they live. That can make reporting more dangerous and increase dependence on the very people causing the harm. They need credible institutions willing to respond to specific offences without demanding that victims either defend or renounce an entire background.

The state should be particularly wary of treating self-appointed community leaders as if they spoke for everyone. A leader may have influence, useful knowledge, and the confidence of many residents. He may also have an interest in keeping disputes internal and dissenters quiet. Consulting him is not the same as hearing those under his authority.

This is one of the less comfortable implications of equal citizenship. The person with the least standing inside a group must be able to seek protection outside it. Public authorities cannot purchase an appearance of tranquillity by allowing local power to become unchallengeable.

The standard should be conduct. Threats, coercion, assault, fraud, and intimidation need a response whoever commits them. The response should be predictable enough that people do not have to calculate whether their case will embarrass an official constituency before asking for help.

\section*{What cannot be negotiated away}

Germany's Basic Law protects dignity, equality, belief, expression, and democratic government under law. These commitments are not satisfied by people occupying the same territory. They establish conditions under which different people may live together without one acquiring arbitrary power over another.\footnote{\href{https://www.gesetze-im-internet.de/englisch_gg/}{Basic Law for the Federal Republic of Germany}, especially Articles 1, 3, 4, 5, and 20.}

Freedom of religion includes the freedom to practise a religion. It must also leave room for a person to change or abandon one. Equality between women and men cannot depend on a family's permission. A woman does not lose her standing because a husband describes obedience as tradition. A child is not an instrument for preserving an adult's authority over the next generation.

A religious offence does not grant permission to threaten a speaker. Antisemitic intimidation remains intimidation whether its author invokes nationalism, religion, or a foreign conflict. The person targeted should not have to wait while officials decide which explanation would be least embarrassing to publish.

These are demanding standards, and long-established citizens can violate them too. Their application cannot depend on the offender's ancestry. Indeed, a state weakens the standards when it presents them as the private property of one population rather than rules binding everyone.

What must be shared is therefore more exacting than politeness and less extensive than uniform private life. People can disagree profoundly about religion, sexuality, family, and politics. They cannot demand that those disagreements entitle them to coercive authority over others outside the law.

The state need not compel affection between neighbours. It must preserve the conditions under which people who do not like one another can nonetheless use the same institutions and invoke the same protection.

\section*{Tolerance and the person who wants to leave}

Some defences of cultural difference describe groups as if their members had freely chosen a complete package. In reality, people are born into households and communities before they can choose. They may value some practices and reject others. They may conceal disagreement to preserve housing, relationships, income, or safety.

A policy that treats the group's most conservative account of itself as authoritative can strengthen those who already control its members. The state may imagine it is being respectful while helping to silence the person who wants to live differently.

Consider an imagined young woman who wants to continue education against her family's wishes. An official response that merely encourages dialogue may help if the family is open to persuasion. It is inadequate if threats and confinement are involved. The response must change with the facts. Respect for family life cannot mean indifference to what happens inside it.

Exit also requires practical means. A formal right to leave an abusive home is less useful without safe accommodation, income, information, and protection. Legal equality and material dependence can coexist. The latter does not nullify the former, but it can make exercising the right dangerous.

This principle reaches beyond migration. Religious communities, political movements, wealthy families, and intimate relationships can all punish departure. The test is whether the individual can refuse without being subjected to unlawful coercion, not whether outsiders find the group's identity attractive.

It is possible to defend a community against unfair attack while opposing its internal abuses. The demand that one must choose between loyalty and criticism is often made by those who benefit from the abuses remaining internal.

\section*{Standards without a performance of gratitude}

Integration is sometimes imagined as a ceremony in which newcomers display enthusiasm for everything the receiving country says about itself. This is an excessive and rather insecure demand. Citizens themselves argue about their country, criticise its history, and dislike its government. Membership should not require newcomers to perform a reverence that others are free to withhold.

There are nonetheless practical commitments a society can require. People need enough language and institutional knowledge to participate where feasible. They must obey the law, accept the equal rights of others, and meet reasonable obligations attached to their status and support. Public disputes must be pursued through lawful channels rather than threats or private enforcement.

These requirements are compatible with criticism. A person can learn the language partly in order to argue that an official has acted unjustly. That is participation, not a failure to be grateful. The state should be pleased that its institutions can be used against its mistakes, even when particular officials are less pleased.

The distinction protects the receiving society too. If integration is measured mainly by ceremonial professions, a skilled performer can satisfy it while disregarding its substance. If it is measured by practical conduct and participation, authorities have something more useful to observe.

No test can inspect every private conviction. A liberal society should be cautious about trying. It can require compliance with public rules and act against credible threats. Demanding proof of inward agreement risks creating the conformity without conviction examined earlier in this book.

\section*{Neighbours and citizens}

Daily life needs habits that law cannot fully supply. People queue, keep noise within bounds, return borrowed things, tell the truth about small matters, and assume that a minor mistake will not become a conflict. These habits make life easier among strangers. Their erosion can feel profound even when no constitutional principle has formally changed.

Newcomers need to learn local expectations, and established residents sometimes need to explain expectations they have never had to articulate. Misunderstanding is possible. So is deliberate disregard. Treating every conflict as one or the other in advance prevents a proportionate response.

Some habits deserve defence because they protect shared use of space. Others are simply preferences that people with power have been accustomed to impose. A common rule about noise has a different purpose from a demand that every neighbour prepare familiar food or celebrate the same holiday. The reason for the rule should be available for discussion.

Trust grows through repeated experience of conduct. It can be damaged when misconduct has no consequence or when innocent people are treated as permanently suspect. Both teach that ordinary reciprocity is insufficient. One rewards disregard; the other tells the compliant person that compliance will never be enough.

Government can support trust by making rules clear and enforcement fair. It cannot manufacture friendship, and it should not promise that every difference will become a source of delight. Coexistence can be successful while remaining awkward, limited, and occasionally irritating. That is already true among people whose grandparents lived in the same town.

\section*{The cost of pretending}

There is a recognisable moment at which a public institution begins to lose trust. People report a specific problem and receive a statement about the moral importance of the policy associated with it. The statement may be true and still fail to answer them. Repeated often enough, it teaches them that the institution's first concern is preserving its account of itself.

This is particularly damaging where migration is concerned. A person who cannot obtain a school place or feels unsafe in a public space may conclude that only the most hostile political voices are willing to notice. Those voices can then attach the real complaint to explanations and remedies that do not follow from it.

Suppressing the complaint does not remove that opportunity. A competent response gives people a reason to seek help without accepting the hostile explanation. Name the shortage, publish credible information, enforce the rule, and explain what remains difficult. Where a policy has imposed costs, say so. Where a rumour is false, show why. This work lacks the immediate gratification of condemning a side, but it is how institutions earn the right to reassure.

The child in the classroom needs this kind of honesty. She needs language teaching, a safe room, capable staff, and standards that prepare her for what follows. She does not need to be exhibited as proof that diversity has succeeded or condemned as proof that the country has failed. Adults owe her an education in the society they have actually made. The argument about how they made it does not cancel that obligation.

\chapter{The Wolf Is Not a Dog}

The wolf in the photograph stands in snow. Its coat is thick, its ears alert, its attention directed towards something outside the picture. It is easy to admire. The photograph asks very little of us beyond admiration.

The animal outside the photograph has to eat. It may kill a deer or a sheep. The distinction matters greatly to people, less so to a hungry wolf. A wolf does not recognise the moral role assigned to it in a conservation campaign, and it cannot be persuaded to preserve that role by choosing only prey that will not inconvenience the campaign.

This is where an affection learned through pets can mislead. A companion animal lives within a relationship arranged around human care. A wild predator lives by a different set of conditions. Recognising its beauty does not require forgetting those conditions. Forgetting them is a poor way to protect either the animal or the people who must live near it.

The wolf is neither a misunderstood dog nor a criminal. It is a predator. Any serious policy has to begin there.

\section*{The animals that attract us}

Our concern for animals is selective. A named dog in a kitchen can command more attention than thousands of unnamed animals in agricultural production. A large mammal in danger can raise money that an equally threatened insect cannot. We need not despise affection to notice that it is an unreliable allocation system.

Agathe Colléony and colleagues studied choices in a zoo animal-adoption programme and found that charisma influenced support more than endangered status. The result identifies a difference between public appeal and conservation need.\footnote{Agathe Colléony and colleagues, \href{https://cris.technion.ac.il/en/publications/human-preferences-for-species-conservation-animal-charisma-trumps/}{``Human Preferences for Species Conservation: Animal Charisma Trumps Endangered Status''}, \emph{Biological Conservation} 206, 2017, pp. 263--269.} It does not show that all money raised through appealing animals is wasted.

A conspicuous species can draw attention to a habitat that also supports less conspicuous life. The public may learn about an ecosystem through the animal it first wanted to see. There is a sensible use for this attraction. The trouble begins when the attraction is presented as though it had already established every ecological claim made on the animal's behalf.

``Biodiversity'' can then become a dignified word attached to a preference formed before the evidence was considered. Which habitats will benefit? Which species? At what population size? Under what land use? These questions should remain open even when the animal on the poster is magnificent.

The reverse prejudice is possible too. A farmer's loss does not establish that the predator has no ecological value. A disliked animal can perform a function; an attractive one can cause harm. Our feelings identify things we care about. They do not remove the need to find out what those things do.

\section*{Returning to an inhabited landscape}

Wolves have expanded naturally across much of Europe. This differs from a programme that captures animals elsewhere and releases them. The European Commission distinguishes the natural return of wolves from specific bear restocking projects.\footnote{European Commission, \href{https://environment.ec.europa.eu/topics/nature-and-biodiversity/habitats-directive/large-carnivores/eu-large-carnivore-platform/facts-and-common-misconceptions_en}{\emph{Facts and Common Misconceptions}}, on large carnivores and their return in Europe.} The distinction matters when assigning responsibility for an intervention, although it does not settle how a returning population should be managed.

A naturally arriving animal can create a problem. A deliberately restored animal can have value. The verbs should describe what happened rather than decide the policy in advance.

The European landscape to which wolves return is not an empty wilderness. It contains farms, roads, settlements, forests managed for timber, hunting, tourism, and protected areas. People have changed prey populations and habitats, and their activities will continue to affect the predator's place within them.

The idea that the wolf will simply restore a complete natural balance is too convenient. Predators can affect prey numbers and behaviour, but the effects depend on local circumstances. In a heavily managed landscape, human actions can reinforce, replace, or counteract ecological effects. A mechanism observed in one place is not automatically reproduced in another.

It may still be worthwhile to restore a native species without promising a dramatic cascade of benefits. The return itself can be a conservation objective. Stating that objective plainly is better than inventing a claim that every element of the landscape will improve once the predator arrives.

There is also no single historical state waiting to be restored. Forest cover, grazing, settlement, and wildlife have changed across centuries. Choosing a baseline is partly choosing what we want to preserve or recover. Evidence can show what existed and what may follow. It cannot make all historical states compatible.

\section*{The sheep also feels pain}

A wolf attack is not only an entry in a compensation form. Sheep can be injured, scattered, or killed. People have to find them, attend to survivors, dispose of carcasses, and decide how to protect the remaining flock. A financial payment may cover some loss without undoing the event.

The sheep's domestication does not make its suffering less real. If sentience is part of the reason to care about the wolf, it cannot be withdrawn from the prey when inconvenient. Nor does the fact that farm animals may eventually be slaughtered make every earlier injury irrelevant. Human responsibility for livestock creates duties of care during their lives.

The European Commission's 2023 analysis estimated that wolves killed at least 65,500 livestock animals annually in the EU, with sheep and goats making up about seventy-three per cent. Reporting systems differ, so the figure should not be treated as a perfectly harmonised census.\footnote{European Commission, \href{https://environment.ec.europa.eu/system/files/2023-12/Council\%20Decision\%20-\%20wildlife\%20and\%20natural\%20habitats.pdf}{proposal concerning the wolf's status under the Bern Convention, 20 December 2023}, drawing on its 2023 analysis of the wolf in the EU. The livestock figures are estimates assembled from differing national reporting systems.}

At the scale of all European livestock, the losses may look small. On an exposed farm they may be severe. A continental percentage spreads the harm across animals never at risk and owners who suffered none. It can be useful for understanding the scale of a sector-wide threat and misleading when used to dismiss a concentrated local burden.

The advocate enjoying the wolf's return may never encounter the loss. The owner encountering it may have little interest in the return. A policy has to connect these two experiences more fairly than a continental average can.

A conservation advocate need not apologise for living in a city. Public goods often concern places far from the people who value them. But distance should make the advocate more attentive to the cost, not more confident that those who report it lack the proper attitude towards nature.

\section*{What prevention asks of a person}

Fencing, shepherding, night enclosures, and livestock-guarding dogs can reduce attacks in suitable circumstances. Their usefulness depends on correct design, maintenance, terrain, flock size, and local conditions. The LIFE WolfAlps EU project includes practical assistance with such measures precisely because coexistence requires work.\footnote{LIFE WolfAlps EU, \href{https://www.lifewolfalps.eu/en/activities/c/c1/}{Action C.1: Prevention and Support for Livestock Farmers}. The project describes assistance with fences, guarding dogs, compensation procedures, and intervention teams.}

``Put up a fence'' can sound like a complete answer to someone who has never installed or checked one on difficult ground. Gates, vegetation, uneven surfaces, electricity, weather, and movement of the flock matter. A measure that is feasible on one pasture may be costly or impractical on another.

Guarding dogs introduce their own responsibilities. They need care and competent handling, and their presence can affect walkers and other dogs. More shepherding requires people and wages. Night enclosures can change labour and grazing patterns. Prevention may be entirely worthwhile, but it is not a free addition to the existing routine.

If the public wants the conservation benefit, the public should be prepared to finance a substantial part of the conditions that make it possible. Paying only after a verified kill leaves the owner carrying the continuing cost of prevention. A system can appear generous on compensation while remaining ungenerous about the work needed to avoid the injury.

The owner also has responsibilities. Where effective and reasonably supported prevention is available, refusing it while demanding unlimited compensation may be difficult to justify. Public funding should be connected to workable expectations, reliable assessment, and assistance rather than an assumption that either side is always innocent.

The question should be whether the arrangement reduces harm fairly. It should not become a contest over which participant can force the other to acknowledge a moral defeat.

\section*{Compensation is not coexistence by itself}

Compensation has several purposes. It can share the financial burden, maintain confidence in public policy, and make losses reportable. Good procedures need to establish what happened without making a claimant navigate unreasonable delays or impossible proof requirements.

The limits should be admitted. Money does not prevent the animal's suffering. It may not replace breeding decisions, lost time, or the confidence to continue a way of working. If repeated attacks make extensive grazing unmanageable, an owner may leave the activity even when individual carcasses are paid for.

That withdrawal can change the landscape in other ways. Some open and semi-open habitats are maintained through grazing. Ending it can benefit some species and disadvantage others. An account that treats the disappearance of livestock farming as nature simply returning has chosen a preferred ecological outcome without naming the choice.

There is no reason every existing agricultural practice must continue indefinitely. Some may be environmentally damaging or dependent on support that could be better used elsewhere. But if predator policy contributes to their end, that consequence should be included in the assessment. It should not emerge only after the public has been told that no important conflict exists.

A fair programme may therefore need prevention support, rapid response, credible compensation, and decisions about particular dangerous or persistently damaging animals. Which measures are lawful and effective depends on the jurisdiction and circumstances. No general affection for wolves can perform that local work.

\section*{Killing and the refusal to decide}

Lethal management is the point at which a policy may lose the appearance of innocence. Many people can support a species' recovery and recoil from killing an individual of that species. The recoil is understandable. It does not settle what should happen when interests conflict.

Sometimes advocates of killing promise more than it can deliver. Removing animals without understanding the local problem may fail to reduce losses as expected. Sometimes opponents of killing present nonlethal measures as universally effective when the evidence is more conditional. Each position needs to show the likely consequences of its proposal rather than rely on the emotional meaning of the act.

There is a difference between preventing a species' disappearance, maintaining a viable population, and preserving every individual animal. A conservation policy can defend the first two without treating the third as an absolute requirement. If individual preservation is the objective, its costs and effects should be stated honestly too.

The word \emph{management} can soften what happens. It may include killing. That should be said. The phrase \emph{letting nature take its course} can also soften what happens. It may include livestock suffering, starvation, or risks transferred to people who did not choose them. Neither phrase removes the need to decide.

Refusing to intervene is an action within a landscape already shaped by human decisions. It may be the best action. It does not return us to a position outside responsibility. We remain responsible for the policy under which we refrain, and for the foreseeable consequences we choose to accept.

\section*{Fear without a fable}

People can fear a predator beyond the risk it presents. Stories, rare attacks, images, and local uncertainty affect the response. Dismissing fear with contempt is unlikely to improve understanding, especially when officials have been careless about other costs.

Risk communication should distinguish sighting, proximity, unusual behaviour, livestock attacks, and danger to people. These are not identical events. A wolf seen in a district is not proof of an imminent attack on a person. Equally, a general statement that attacks are rare should not be used to ignore a particular animal's concerning behaviour.

The authorities need a credible procedure for reports and response. People should know what information to provide and what will happen with it. An institution that appears determined to explain away every report invites citizens to trust the most alarming unofficial account instead.

Fear can also become politically useful to those who oppose any predator presence. A vivid exceptional event can be treated as the normal condition, while the proposed response is left unspecified. How many animals would be removed? From which area? With what effect on conservation obligations and future recolonisation? Anger at a real injury does not make these questions unnecessary.

The public deserves neither soothing evasions nor a return to fairy-tale monsters. The animal's actual behaviour is consequential enough without being recruited into a story about civilisation and barbarism.

\section*{Our own part in the scene}

It is easy to expose sentimentality in the person who loves wolves and eats lamb. It is also easy to expose it in the person who loves a pastoral landscape without asking how its animals are raised and killed. These observations may identify inconsistency. They do not finish the ethical argument.

Most people have attachments that are more particular than their stated principles. We know one animal and not another. We respond to visible suffering and overlook remote suffering. A useful criticism should widen attention rather than merely allow the critic to enjoy catching someone out.

The sheep owner can be asked about animal welfare without losing the right to object to attacks. The conservationist can be asked to pay for coexistence without being required to stop valuing the wolf. A person can oppose needless cruelty in both farming and wildlife management while admitting that no available arrangement prevents all animal suffering.

Predation is part of the natural world, but naturalness is not itself a moral instruction. We do not ordinarily let nature decide whether to treat a pet's infection or protect a child from exposure. Neither can we redesign every ecological relationship around the elimination of suffering without consequences we scarcely understand. Our capacity and knowledge set limits on what intervention can achieve.

That leaves a less satisfying position than either unrestricted protection or unrestricted elimination. We must specify what we seek to conserve, what losses we will accept, what prevention we will finance, and when we will intervene. The answers may differ across places. They will need revision as populations and conditions change.

The wolf in the snow remains beautiful. Knowing what it eats need not diminish that response. It should diminish our willingness to confuse the response with a complete policy. The photograph ends at its edges; the life around the animal does not. A sheep, a shepherd, a walker, a forest, and the less conspicuous creatures of the valley all belong in the decision we make after looking at it.

\chapter{What Happened to the Mice}

On 9 July 1968, John B. Calhoun placed four pairs of mice in an enclosure supplied with food, water, nesting material, and protected shelter. Predators could not enter. The mice could not leave. The population grew, growth slowed, and eventually reproduction failed. The surviving animals aged in an environment that still supplied the necessities intended to keep them alive.

The experiment became known as Universe 25. It offers a disturbing sequence of events. It also offers a temptation: to turn the mice into a miniature version of whichever human society the observer already believes is failing.

For one reader, the enclosure demonstrates the evils of overcrowding. For another, it proves that welfare and comfort destroy the will to live. A third sees young people withdrawing into private consumption; a fourth sees the collapse of family life. The mice become witnesses in arguments whose terms they could not possibly recognise.

Before drawing those conclusions, we should spend some time in the actual enclosure. Its history is unsettling enough without making it prove everything.

\section*{Food was abundant; the world was finite}

Universe 25 was roughly two and a half metres square, with a structured arrangement of feeding sites, water, and nesting compartments. Its design allowed a population larger than the one ultimately reached to obtain food and water. The experiment was intended to reduce several ordinary causes of mortality. It did not give the animals unlimited space or every condition they might need for a viable social life.

That distinction is often lost when the enclosure is called a paradise. \emph{Paradise} is a human judgement about the arrangement. A supply of food and protection from predators removes serious hazards, but it does not tell us how mice experience repeated encounters, access to territory, nesting, or social position.

An empty nesting compartment is not necessarily equivalent to a workable place to rear young. Routes of access and the behaviour of other animals matter. Nominal capacity tells us what the structure can physically contain; it does not fully describe what its inhabitants can use.

The inability to leave is especially important. Animals that cannot establish themselves locally may disperse in a less restricted environment. In Universe 25, moving away did not lead to another landscape. All future attempts to find a place occurred within the same structure and population.

Calhoun had worked on rodent populations before this experiment. His earlier descriptions of rats introduced the memorable phrase \emph{behavioural sink}. Popular retellings sometimes merge those studies with Universe 25, as if every reported observation came from the same set of mice. The distinction between experiments matters when we ask what was actually observed and under which conditions.

\section*{The first expansion}

The founding mice initially adjusted to the enclosure and began to reproduce. Calhoun described an early period of establishment, followed by rapid growth. The population doubled approximately every fifty-five days during its most rapid expansion, reaching around 620 weaned mice before the rate slowed.

At this stage, a reader can understand why the arrangement might appear successful. The animals were surviving and producing young. The population curve rose. Food had not become scarce. If the only measures were survival and numbers, the project had much to show.

But a rising total can conceal differences within a population. Access to nesting sites and social positions was not evenly distributed. Young mice entered a setting already organised by adults. The conditions faced by the founders were not simply repeated for each generation.

This is one of the experiment's useful observations. An environment can remain physically similar while its meaning for a new entrant changes. The first inhabitants have opportunities to establish themselves that later inhabitants do not encounter in the same way. The available resources have become part of an existing social arrangement.

We should resist translating this immediately into a story about property-owning parents and excluded young adults. The human comparison may be worth discussing, but it is another argument. For now, the point concerns the mice: the population's history changed the conditions experienced by later members.

\section*{The slowing}

As growth slowed, Calhoun described increasing disruption in social behaviour. Some males were subjected to repeated attacks and withdrew. Maternal behaviour deteriorated. Young were rejected or inadequately cared for, and fewer survived. Aggression and withdrawal could coexist within the same supplied environment.

The provision of food did not prevent these outcomes. That fact is important, but it does not establish that provision caused them. Several aspects of the arrangement changed together as the population grew: density, patterns of contact, competition over positions, and the developmental environment of the young. The experiment does not allow us to point to one of these with complete confidence and declare the others irrelevant.

The difference between density and experienced crowding is useful here. Density is a relation between numbers and area. Crowding involves the encounters, interruptions, and lack of escape created by the arrangement. Two spaces with the same number of animals can differ in access, separation, and the possibility of avoiding conflict.

The same distinction is familiar in less consequential settings. A room with many people can be comfortable if movement, acoustics, and expectations are well arranged. A smaller gathering can feel oppressive when someone blocks the exit or controls every interaction. This does not make a human room equivalent to a mouse colony. It shows why head count alone is an incomplete description of social space.

Calhoun interpreted the deterioration as a failure of the organisation through which the population reproduced itself. There was food for bodies, but the connected behaviours of courtship, nesting, care, and social interaction were becoming unreliable. The young inherited those conditions rather than the circumstances encountered by the founders.

\section*{An ending that took time}

The population reached roughly 2,200 around day 560. A few mice born up to about day 600 survived beyond weaning. Later pregnancies did not produce a surviving new generation. The distinction matters: the last successful birth was not the last conception.

The population then declined as its members aged. In a final editorial update dated 13 November 1972, Calhoun recorded twenty-seven survivors. His paper projected a later death of the last male; it did not report the completed physical extinction of every mouse as an event already observed.\footnote{John B. Calhoun, \href{https://journals.sagepub.com/doi/10.1177/00359157730661P202}{``Death Squared: The Explosive Growth and Demise of a Mouse Population''}, \emph{Proceedings of the Royal Society of Medicine} 66(1, part 2), 1973, pp. 80--88, doi:10.1177/00359157730661P202. See the population-curve caption for the November 1972 update and the main text for the distinction between last surviving birth and later conceptions. The account here separates Calhoun's observations from his broader human interpretation.}

The distinction between an observed count and a projected ending matters, particularly in a story used to warn against unwarranted optimism. Calhoun had documented reproductive failure and a dwindling group of survivors. The final physical ending still lay ahead.

The decline is disturbing because it was not an abrupt shortage that could be relieved by another delivery. The food remained. The population became smaller, yet ordinary reproduction did not resume in the way a simple account of overcrowding might lead us to expect. Something about the animals' development and social behaviour had outlasted the period of greatest numbers.

Calhoun also reported work by a colleague moving small groups from crowded populations to low-density settings. Reproductive and social behaviour remained seriously impaired. Those observations warrant interest. They do not, in the form reported there, establish that every possible intervention would fail or that the condition was genetically fixed.

The difference between \emph{did not recover under these conditions} and \emph{could never recover} is substantial. The second claim requires much more evidence. Human arguments about irreversibility often cross that distance without noticing it.

\section*{The animals called beautiful}

Some males attracted Calhoun's attention because they did not fight or court. Their coats were in good condition, without the wounds of the more embattled animals. Their activity was largely confined to eating, drinking, sleeping, and grooming. He called them the ``beautiful ones.''

The name made them memorable. It also helped turn a behavioural description into an accusation ready for use against human beings. A person who spends time on appearance, avoids conflict, or does not have children can be placed under the same label, and the resemblance is treated as an explanation.

This is precisely where care is needed. A mouse's failure to develop reproductive behaviour and a human decision not to become a parent are not the same event. Human beings can remain childless while caring for others, creating work, teaching, participating in institutions, or simply living a life they value. The absence of reproduction does not define the whole life.

Nor is visible conflict a measure of health. A scarred animal is not more admirable because another animal attacked it. A human being who avoids an abusive institution may be acting sensibly. We cannot convert withdrawal into pathology merely because an observer values participation in the institution being avoided.

The genuine question is what the withdrawal means in its setting. Is an animal capable of the behaviours needed to sustain its population? Is a person choosing solitude, trapped in it, or unable to form relationships they want? Similar appearances can conceal different conditions. The vivid label is least helpful when it makes us stop asking.

\section*{What was tested?}

An experiment gains explanatory power by distinguishing among possible causes. If several conditions change together, the result can demonstrate that a particular arrangement failed without isolating the decisive mechanism.

Universe 25 combined a laboratory strain, an enclosure, a particular spatial design, abundant supplies, restricted exit, and a developing social history. These are not minor qualifications attached to an otherwise universal law. They are the conditions under which the outcome occurred.

To establish that abundance itself caused collapse, we would need comparisons that separate abundance from confinement and other features. To establish that density alone caused it, we would need to examine arrangements with differing density and social access. To establish that a learned impairment was permanently irreversible, we would need much more than a dramatic account of unsuccessful transfers.

This does not make the observations worthless. They reveal a failure that a purely nutritional account would miss. They suggest questions about development, social organisation, and environmental design. They can guide further inquiry. But they do not exempt that further inquiry from being done.

The distinction between observation and interpretation matters even when the interpreter is the investigator. Calhoun's account deliberately reached towards human concerns and used religious language about two forms of death. These ambitions are part of the paper. The experimental record does not become a proof of every philosophical conclusion attached to it by its author.

A scientist can be perceptive about a pattern and expansive in interpreting its reach. Respect for the work requires us to distinguish the two rather than accept or reject the whole person.

\section*{The comfort explanation is too easy}

The popular lesson that comfort destroys life is attractive because it turns a complex outcome into a familiar moral warning. People have long suspected that ease makes others weak, particularly younger people whose difficulties take unfamiliar forms. The mice seem to give that suspicion scientific equipment.

But the animals did not live without pressure. They lived with pressures the provision of food did not remove. Calling their environment comfortable because calories were available assumes that calories define comfort. Repeated attack, disrupted care, and inability to escape are not indulgences.

The inference also treats deprivation as though it must be the cure for any failure occurring without deprivation. That does not follow. If a well-fed child is lonely, hunger is not the obvious remedy. If a workplace is materially comfortable and socially destructive, making the chairs worse will not necessarily improve it.

Material security can coexist with social failure because it is not the whole of social life. That is a narrower claim than saying security causes the failure. It is also a more useful one. It directs attention towards what is missing rather than inviting us to remove what is present.

The opposite mistake is to suppose that adequate provision must be sufficient and therefore any remaining difficulty is an individual's ingratitude. The experiment gives us reason to question that assumption too. An arrangement can meet a list of material needs and still be badly suited to the creatures living in it.

\section*{Why the analogy persists}

The story persists partly because the enclosure is visible. We can imagine its walls, the nesting boxes, the expanding population, and the well-fed survivors. Human social problems are harder to hold in one picture. The enclosure offers a manageable shape for anxieties that otherwise sprawl across work, family, housing, technology, and politics.

That can be intellectually useful. An analogy may help us notice a question we had neglected. Does a system provide meaningful ways to enter adult roles? Can people leave a destructive setting? Are young people learning the practices they will need? Have we mistaken the delivery of resources for the creation of a workable life?

The analogy becomes harmful when the picture supplies the answer before human evidence is examined. A falling birth rate is assigned to a behavioural collapse. A dense city becomes a doomed enclosure. Digital entertainment becomes grooming. Different phenomena are gathered under one story because they look alike from a sufficient distance.

The apparent explanatory power comes partly from discarding differences. People can speak about their situation, change institutions, invent occupations, organise care in new ways, and deliberately alter family size. Their choices occur within constraints, but the constraints are not identical to a metal wall around a laboratory population.

Human recovery can also draw on people and practices outside a failed setting. A school can recruit teachers, a town can receive new residents, and a person can encounter a relationship that changes what seems possible. None guarantees recovery. Each matters when judging whether the enclosure is an adequate model.

\section*{The question worth retaining}

We should retain the experiment's challenge to a simple promise: supply the material necessities and flourishing will follow automatically. The promise is insufficient. Relationships, development, access, and the organisation of space matter too.

We should discard the equally simple counter-promise: restore hardship and flourishing will return. Hardship can injure the capacities it is supposed to strengthen. It can destroy families, narrow choices, and leave people with less ability to help one another. A policy needs evidence of its effects, not a mouse-shaped justification for severity.

The declining colony remains a sorrowful thing to contemplate. The supplied bodies continued while the population's means of renewal failed. We do not need to turn it into a prophecy about ourselves to find that unsettling.

Its proper place in a human argument is as a prompt for closer attention. Where people have enough to eat and still cannot find a tolerable way to live, ask what their arrangements permit and prevent. Do not assume that full shelves make the question frivolous. Do not assume that the answer has already been found inside an enclosure built for another species.

\chapter{What Comfort Does Not Explain}

Imagine a man on the first Monday of retirement. For years he has wanted relief from the alarm clock, the commute, and the petty decisions of his manager. Now he has it. His pension is adequate. Nothing urgent requires him to leave the house.

At first the absence is pleasant. Later he notices how much of his day had been organised by things he disliked. Colleagues expected him. A task could be finished. Someone noticed if he failed to arrive. He did not love all of this, but it gave time a shape that freedom has not yet replaced.

It would be absurd to conclude that retirement is a mistake or that he needs a worse manager. What has disappeared is a mixture of burdens and supports. Some he is glad to lose. Others he had not recognised as supports until they were gone.

The example concerns a person, not a population of mice. It requires no biological theory of decline. It does raise a question the mouse experiment cannot answer for us: what arrangements make a materially secure life one in which a person can find a place?

\section*{Enough to live on}

Material security matters enormously. A person anxious about food, housing, or heating has reasons for distress that should not be converted into a philosophical lesson about meaning. Poverty does not become ennobling because some wealthy people are miserable.

Yet adequate income does not guarantee friendship, a useful occupation, a tolerable family, or a reason to expect tomorrow with interest. These are different goods. They interact with money without becoming identical to it.

The error is to make either observation refute the other. Someone says that people need more than consumption, and is answered as though she had recommended hunger. Someone says that rent and childcare are too expensive, and is advised to recover a sense of purpose. Each answer avoids the need actually being described.

A serious account of flourishing must be able to say that a person needs resources and that resources need to be usable within a life. A library card does little for someone who cannot reach the library. Free time may feel empty if every friend is elsewhere. Access to a thousand entertainments may not provide anyone who expects a visit.

These failures are not proof that abundance has poisoned the person. They show that provision and participation have been arranged separately. Sometimes the gap can be narrowed through practical changes. Sometimes illness, loss, or circumstance leaves a difficulty no institution can completely remove.

\section*{Being needed, and being used}

People often want to be useful to others. A task can place them within a relationship: someone benefits if it is done and notices if it is not. This can make an ordinary activity more satisfying than an elaborate entertainment designed solely for their consumption.

The wish is easily exploited. An employer describes exhausting work as a calling. A family makes one person indispensable by declining to share the care. A political movement offers belonging in exchange for continuous availability. The person receives a place, but at a price that may leave little room for a life outside it.

We should therefore be exact when praising responsibility. Responsibility can give someone standing and scope for judgement. It can also mean being blamed for a result over which they have little control. The first can develop a person; the second can wear them down while praising their dedication.

Consider a volunteer asked to organise an event. If she has a clear task, resources, authority to make decisions, and others who do their share, the work may be rewarding. If every decision is overruled and every failure assigned to her, the same title conceals a different arrangement. Calling both participation tells us too little.

The retired man may find satisfaction in a useful activity. He should not have to accept exploitation to obtain it. Nor should paid workers be displaced simply because people seeking purpose can be persuaded to work for nothing. Institutions need to distinguish a worthwhile contribution from an inexpensive way to transfer labour to someone eager to belong.

\section*{Entry into adult life}

The difficulty can appear before a working life begins. A young person acquires qualifications but cannot find stable work or independent housing. Family formation seems financially precarious. Opportunities exist in the abstract, yet each requires experience, money, or connections that are difficult to obtain without already having a place.

An older observer may see hesitation and call it immaturity. Sometimes hesitation is avoidable. Sometimes the available route is substantially different from the route the observer remembers. A judgement about character should not substitute for examining the prices, contracts, and expectations now facing the person.

There is also a loss when entry routes become entirely credentialled or informal. A task once learned through supervised practice may require another certificate. An unpaid placement may exclude someone who must earn. Recruitment through personal contacts can make an institution feel open to those already near it and inaccessible to everyone else.

These arrangements can be changed without pretending that every aspiration can be fulfilled. Training can be paid. Qualifications can be required where they protect a real standard rather than where they merely filter applicants. Work can include room to learn. Housing policy can recognise that independence is difficult when every modest dwelling is unaffordable.

None guarantees a meaningful life. It changes whether people have a practicable chance to attempt one. A society that tells the young to take responsibility should provide opportunities in which responsibility involves more than absorbing risks that established people have transferred to them.

\section*{Contact without company}

Digital communication makes forms of contact possible that once required travel or could scarcely occur at all. A person with an uncommon interest can find others. Friends separated by distance can remain in touch. Someone confined at home may participate in conversation, learning, and work.

It would be careless to describe all this as a substitute for real life. It is part of real life for the people involved. The relevant question is what the contact permits and what it leaves absent.

A stream of reactions can make attention abundant while companionship remains scarce. Someone may know that a photograph was liked without knowing whom to call in an emergency. A discussion can be lively but disappear when its participants have no continuing commitment to one another. There are online relationships that become durable and supportive; there are also crowded digital spaces in which a person remains profoundly alone.

The design of the service matters. A platform interested in retaining attention may encourage whatever keeps the user returning, whether that is affection, anxiety, outrage, or the hope of recognition. The user's wish for company and the platform's wish for time are not always aligned.

This observation does not require a claim that technology has caused a universal decline. We can examine a more concrete question: after an hour with this service, what has become easier or harder to do? Has a relationship developed, a task advanced, or a useful exchange occurred? Or has the person mainly been kept in a state of expectation that can be renewed indefinitely?

The answer may differ between uses of the same device. A blanket judgement about screens can obscure the very differences that would help someone make a better choice.

\section*{Solitude is not failure}

We should be equally cautious about judging a life from its visible level of social activity. A person can be alone and content, or surrounded by others and lonely. The WHO's work on social connection distinguishes the structure of social contact from the subjective experience of loneliness.\footnote{WHO Commission on Social Connection, \href{https://www.who.int/groups/commission-on-social-connection/report}{\emph{From Loneliness to Social Connection: Charting a Path to Healthier Societies}}, 2025; see also WHO's \href{https://www.who.int/news-room/questions-and-answers/item/social-connection}{explanation of social connection, isolation, and loneliness}.} Counting encounters does not tell us everything about their quality.

Some solitude is chosen because it allows concentration, rest, reading, walking, or relief from demands. A person who has spent the day attending to others may need an evening without company. There is no reason to interpret this as withdrawal from life.

Other solitude is endured. A person wants contact but cannot find it, cannot travel, fears rejection, or has lost the relationships around which life was organised. The outward appearance can be similar. Advice that assumes one condition may be useless or intrusive in the other.

This matters for the mood of scepticism developed in this book. Freedom from agitation does not require a permanent retreat from other people. Nor does a quieter life have to defend itself by claiming superior wisdom. Some people are less sociable than others. Some periods of life are quieter. The question is whether the person has room to live as they wish, not whether their calendar satisfies a public ideal of engagement.

A campaign against loneliness can itself become oppressive if it treats every solitary person as a problem to be corrected. Institutions can make contact available and remove obstacles without demanding that everyone accept the same quantity of it.

\section*{What disappeared with the old community}

Nostalgia often gathers together goods that did not belong to everyone. The close community remembered for mutual help may also have monitored behaviour, punished departure, and protected powerful members. A person who fitted its expectations can remember security; someone who did not may remember having no place to hide.

The mutual help was not necessarily imaginary. Neither was the coercion. Losing the old arrangement can therefore produce relief and loneliness at once. People gain privacy and lose the neighbour who would have noticed their absence. They escape an intrusive family and discover how much practical care the family supplied.

We cannot restore only the pleasant half by praising community in general. We need to ask how useful forms of interdependence can exist with more freedom to refuse and leave. Associations, public spaces, libraries, clubs, shared projects, and accessible services can help, but their quality depends on how they are run.

A meeting place is not a community because a grant application calls it one. People need reasons to return, something they can do together, and enough influence to avoid feeling like clients of someone else's enthusiasm. The person organising the project must also be able to stop without the whole undertaking collapsing upon them.

There is a modest virtue in institutions that do not promise intimacy. A reliable library, swimming pool, allotment, or local class can provide repeated contact and shared activity without requiring confessions of belonging. People may become friends, or they may simply recognise one another. Both can make a place more inhabitable.

\section*{The right to be ordinary}

Modern life can place an exhausting demand on people to make themselves exceptional. Work should be a passion. Leisure should develop the self. Travel should transform it. Even rest can be judged by whether it improves later performance.

The demand is difficult partly because distinction is comparative. Everyone cannot occupy the most admired position. If a tolerable life depends upon exceptional recognition, many people are being invited to understand an ordinary outcome as a personal failure.

An adequate wage, decent relationships, competence in a task, and some time one can use freely are not contemptible ambitions. Nor is a life diminished because it never becomes an interesting public story. The person who keeps an ordinary service functioning may contribute more than someone who speaks fluently about changing the world.

This does not mean ambition should be discouraged. It means that a society needs room for people whose ambitions are modest, changing, or disappointed. If losing a competition produces ruin or humiliation, the competition has been made to decide too much.

There is a connection with inheritance and merit. When the conditions of a decent life are insecure, parents become more desperate to secure advantage and individuals more dependent on ranking. Providing a tolerable floor does not remove competition, but it can reduce how much of a person's life hangs upon winning it.

The fear that security will make everyone inactive assumes that only danger can move people. We have reasons to work, care, create, and learn that are not reducible to fear. The problem is arranging opportunities for these motives without pretending that everyone will respond in the same way.

\section*{No official can supply the whole meaning}

A state can support conditions in which people form relationships and undertake activities. It should be wary of announcing a shared purpose comprehensive enough to settle the meaning of their lives. Such purposes often become demands for allegiance, with dissent described as selfishness or decay.

An unhappy society may welcome a project that tells everyone where they belong and what their sacrifice means. The promise is powerful precisely where ordinary institutions have failed to offer a credible place. A movement can turn private frustration into collective certainty and make disagreement feel like an attack on the recovered self.

This is one reason material and social failures are politically consequential. People who cannot see a workable future may become receptive to a promise of total renewal. The promise can provide excitement without providing the housing, care, work, or institutional competence that would make daily life better.

We should not infer that dissatisfied people are incapable of sound judgement. Their complaints may be accurate. The danger lies in attaching those complaints to an authority that asks to be released from ordinary restraints because only it can restore meaning.

The next part of the book examines that authority. Before reaching it, we can state a limit: public power should not require people to find fulfilment in obeying it. A person may benefit from good government while finding meaning elsewhere, or while remaining uncertain about the word \emph{meaning} itself.

\section*{What remains after the diagnosis}

The retired man may eventually find a satisfying routine. He may volunteer, care for someone, learn something, or simply discover pleasures that the old working week left little time to notice. He may also remain unhappy. No social theory entitles us to promise a favourable ending to his particular life.

The same restraint is needed in discussing a generation. We can identify obstacles, improve institutions, and make participation less punishing. We cannot guarantee that everyone will want the roles on offer or that every loss will find a replacement. Some people will miss a person who has died rather than a service that can be provided. Some will have possibilities and still struggle to use them.

Melancholy begins, in part, with accepting that these facts do not always contain a hidden lesson. A lost relationship need not be an opportunity. An empty afternoon need not become productive. A society can take suffering seriously without insisting that every suffering person must become an example of resilience.

What we can reject is the laziness of the total diagnosis. The young are not mice spoiled by comfort. The old are not merely dependants. The solitary are not all failed participants. People live within arrangements that can help or obstruct them, and their responses deserve attention at a scale smaller than a verdict on civilisation.

Comfort leaves many human difficulties unsolved. That is a reason to understand those difficulties, not to become nostalgic for cold rooms and hunger. We may have removed one burden and discovered another. The discovery asks for more thought, and perhaps for less confidence in anyone who claims that a single deprivation, technology, or shared mission will make life whole.

\part{Living with Other People's Power}

\chapter{Why We Give Anyone Authority}

An official letter tells a homeowner that part of her land is needed for a new railway. She objects. The line could run elsewhere, she says, or the project could be abandoned. The authority disagrees. After the prescribed process, it may be able to take the land with compensation even though she has never consented.

The railway may be useful. The process may be fair. The compensation may be substantial. The decision still involves compulsion. A public purpose does not turn the owner's refusal into agreement.

This is an imagined case, but it identifies something present throughout government. Taxes are not subscriptions that can be cancelled when the subscriber dislikes the programme. A court's order is not advice. A prohibition can be enforced against a person who regards it as foolish. Public power can compel people to do what they would not otherwise do.

Any defence of government should begin by admitting this. The language of service is incomplete where the service is backed by force. The question is when that force is justified, who may exercise it, and how the person subject to it can challenge an error.

\section*{When private agreement is insufficient}

There are reasons to authorise binding decisions. A transport network can be obstructed if every affected owner has an unlimited veto. Pollution can cross property boundaries. A person defrauded or assaulted needs more than the offender's willingness to discuss the matter. Defence, public health, and shared infrastructure can require coordinated action that voluntary agreement does not reliably supply.

People without an effective public authority are not necessarily free of coercion. They may face the coercion of whoever controls weapons, employment, land, or local access. If a criminal group rules a neighbourhood, the absence of effective policing is not experienced as liberty by those it threatens.

This gives government a difficult double task. It must restrain private power where that power violates rights or undermines a workable public order. It must also restrain the officials entrusted with doing so. The police officer needs authority to stop an assault and limits that prevent the authority from becoming another means of assault.

Thomas Hobbes made the problem of insecurity central to his account of political order. Agreements are fragile where people lack a common power capable of enforcing them and cannot trust others to keep them. His proposed sovereign was far less constrained than the government defended here, but the problem he identified cannot be answered by wishing coercion away.\footnote{Thomas Hobbes, \href{https://www.gutenberg.org/files/3207/3207-h/3207-h.htm}{\emph{Leviathan}}, 1651, especially chapters XIII and XVII.}

John Locke pressed a different part of the difficulty: people need known rules and an authority to which disputes can be referred, yet an absolute ruler leaves the subject without a satisfactory appeal against the ruler's own actions.\footnote{John Locke, \href{https://www.gutenberg.org/files/7370/7370-h/7370-h.htm}{\emph{Second Treatise of Government}}, published in 1689, especially chapters VII and IX.} Establishing an umpire is not enough if the umpire can decide every case in his own favour.

We need not accept either philosopher's entire account to see the conflict. The institution that protects a person from private force acquires force of its own. Its usefulness is a reason to construct it carefully, not a reason to stop asking what it may do.

\section*{The power to decide is not the power to be right}

Government must sometimes decide before agreement is complete. This is one reason decisiveness has appeal. A public exhausted by delay may welcome someone who promises to settle the matter, issue the order, and make the work begin.

Delay can indeed be costly. A bridge may remain unsafe while agencies argue. A planning system may make ordinary housing needlessly difficult to build. A state unable to implement lawful decisions invites people to seek influence outside its procedures. Defending every delay as deliberation gives liberal institutions a reputation for incompetence they need not accept.

But speed is a property of a decision, not evidence of its quality. A ruler can build the wrong project quickly. An administration can make a mistake more efficiently by removing everyone able to object. The cost becomes visible later, when the institutions needed to stop the project have already been weakened.

We should distinguish a capable state from an unchecked state. Capability includes trained staff, reliable records, clear responsibilities, adequate resources, and the ability to carry out a decision. Unchecked power concerns the absence of restraints on what may be decided. The two can coexist, but neither entails the other.

An arbitrary regime can be administratively incompetent. A constitutional government can act quickly within clear authority. Treating the choice as weakness or dictatorship allows bad administration to discredit constraints that were not its cause.

The practical response to a slow permit process may be to simplify rules, provide staff, and identify a responsible decision-maker. It need not be to give a minister permanent discretion to approve friends and obstruct opponents. The first repairs a function; the second may merely change who benefits from its failure.

\section*{The successor enters the room}

Suppose, for the sake of argument, that today's leader is unusually competent and public-spirited. He understands the problem, hears good advice, and intends to use additional powers well. What happens when he leaves?

This is not an objection based on cynicism about every ruler. It is a question about an institution designed to outlast one person. The powers remain available to a successor whose judgement and motives may differ. They may also change the current ruler, who learns that resistance can be removed more easily than answered.

A useful test for a proposed power is to imagine it in the hands of the political opponent one trusts least. Would the authority to close a newspaper, appoint judges at will, direct prosecutions, or classify criticism as a threat still appear necessary and benign? If the defence depends on the present holder's identity, it is a defence of a person rather than a safe arrangement.

Constitutional restraints sacrifice some convenience now to protect people later. That can be frustrating when a sympathetic government encounters them. It is also their purpose. A restraint that disappears whenever those in office regard their cause as urgent is not much of a restraint.

The same issue applies to appointments and informal precedents. An institution can be weakened without a dramatic change in law if its officers learn to serve the governing party instead of the office. A practice justified once as an exceptional necessity becomes available as an ordinary method to the next administration.

Political memory is often remarkably short at this point. People condemn an opponent's use of powers they helped create. The inconsistency is embarrassing, but the deeper loss is practical: the institution no longer provides the protection they now need.

\section*{Which part of life belongs to government?}

Even a competent, accountable government should not direct every activity. People need room to speak, associate, worship or refuse worship, pursue work, form intimate relationships, and use their time without having to show that the activity advances an official purpose.

Isaiah Berlin's distinction between negative and positive liberty helps identify two different questions. One concerns the area left free from interference by others. The other concerns who governs a person and the wish to direct one's own life. Berlin warned that an authority claiming to represent a person's true or rational self could come to dismiss the person's actual refusal as ignorance.\footnote{Isaiah Berlin, \href{https://berlin.wolf.ox.ac.uk/published_works/tcl/tcl-e.pdf}{\emph{Two Concepts of Liberty}}, 1958. The distinction does not imply that all positive accounts of liberty lead to coercive paternalism; the danger discussed here is the authority's claim to override actual persons in the name of their allegedly truer wishes.}

The danger is familiar. An official says that a restriction does not really diminish freedom because it helps citizens become what they ought to want to be. The dissenter's disagreement then becomes evidence that the programme is needed more urgently. The claim is difficult to contest because the authority has appointed itself interpreter of the citizen's proper wishes.

Some public interventions expand practical choices. Education can enable a child to do things previously impossible. Protection from violence can make formal rights usable. A person without resources may have little ability to exercise a freedom that exists on paper. These are reasons to examine what enables agency, not reasons to give the state ownership of the agent's eventual choices.

The educated person may use the education to reject the government's values. The protected person may lead a life officials regard as wasteful. Freedom includes outcomes its defender would not have designed. A state that permits only approved forms of independence has retained the important decision for itself.

\section*{The boundary is contested}

There is no simple line separating every private matter from every public one. A factory owner's activity affects workers and neighbours. Parents' choices affect children. Speech can involve threats as well as opinions. Property can support independence and also give one person power over another.

These conflicts require reasons. Which harm is alleged? How direct is it? What intervention is proportionate? What would be lost through intervention, and who decides? Invoking privacy should not conceal coercion within a home. Invoking a public interest should not make every disliked private choice available for regulation.

The railway case illustrates the difficulty. A route can produce broad benefits while seriously harming particular owners. An absolute veto might prevent a valuable project; unlimited official discretion could permit arbitrary dispossession. Procedures, evidence, compensation, and appeal attempt to make the decision answerable to those affected. They do not make the interests identical.

The losing owner may remain angry after a fair process. The state should not require her to agree that the outcome was good. It owes her the protections the process promises and the freedom to continue criticising the decision. Political legitimacy cannot depend on unanimous emotional reconciliation after each exercise of power.

This is an advantage of treating disagreement as enduring. Institutions can be designed for people who continue to hold different views, rather than for a public expected to discover that the authorities were right once the explanation has been repeated often enough.

\section*{Why elections matter, and what they leave open}

Elections provide a way to authorise and replace rulers without a struggle for physical control each time. They can make office conditional and give opposition a reason to remain within the system. These are substantial achievements even when the candidates and campaigns are disappointing.

An election does not make every act of the winner legitimate. A majority can support censorship, discrimination, or the destruction of the procedures through which a later majority might replace it. The fact that a power was obtained electorally does not settle whether its use preserves democratic government.

Nor does an election prove that voters share one coherent programme. They may support the same party for different reasons, tolerate some policies to obtain others, or choose it chiefly to remove an incumbent. The victorious politician's description of a comprehensive mandate often exceeds what the votes can establish.

The next chapter examines a more precise limit: even when each voter has a clear and consistent ranking, majority comparisons can produce a cycle. We should understand that difficulty before speaking too confidently of what the people have decided as a single mind.

For now, the point is that democratic authority needs institutions around the vote. Opposition must remain lawful. Information must circulate. Rules must permit a future contest. Courts and other checks must be able to address abuses within their proper remit. An election that can be won once and then made impossible to lose is a route to permanent rule, not a continuing system of authorisation.

\section*{Experts cannot relieve us of judgement}

One response to electoral disappointment is to propose that experts govern. They know more about many relevant matters, and some decisions require specialised competence. A bridge should not be designed by a popular vote on load-bearing calculations.

Yet the bridge's location, cost, alternatives, and distribution of burdens involve choices beyond the engineer's technical authority. Expertise can clarify consequences. It does not establish which competing interests should prevail in every case.

Experts also disagree, specialise, acquire institutional loyalties, and operate under incentives. A government by experts must still decide which experts, appointed by whom, responsible to whom, and removable by what process. Avoiding elections does not avoid politics. It relocates political choices to the selection and supervision of the experts.

There are good reasons to give some institutions operational independence. The scope of that independence should be defined, and their work should remain open to appropriate scrutiny. Otherwise technical vocabulary can become a means of excluding the people who bear the consequences.

The person whose land is taken may need an independent technical assessment. She also needs a way to challenge a value judgement disguised as a calculation. Respect for expertise includes being exact about where it ends.

\section*{Consent that never becomes unanimous}

No large political society secures explicit consent to every rule from every person governed. Birth, residence, dependence, and the cost of leaving complicate any simple story that remaining in a country proves agreement with its government.

This should restrain the government's language. The dissenter has not necessarily authorised the particular policy merely because she pays taxes or lacks a practical opportunity to emigrate. She is subject to a common order whose legitimacy must be defended through more than the claim that everyone is free to go.

At the same time, allowing each person to disregard every rule they reject would make many common arrangements impossible. The disagreement cannot be resolved by pretending compulsion is consent or by pretending common rules can operate without compulsion. It must be managed through limits, representation, rights, appeal, and opportunities to change policy.

These devices are imperfect. Access to them can be unequal, and institutions can fail. Their value lies partly in making failures identifiable and contestable without requiring a revolution. A person may lose a case and still retain standing to pursue another.

That continuing standing is more important than the flattering claim that government always expresses the public's deepest wishes. It permits citizens to remain citizens when they disagree with those temporarily in office.

\section*{A government one can live beside}

The best defence of limited constitutional government is comparative. We do not need to believe that its voters are wise, its officials selfless, or its outcomes consistently just. We need to examine whether it offers better ways to protect people, expose error, distribute power, and remove rulers than feasible alternatives.

This is a less stirring defence than a promise of national renewal under a leader who understands everything. It is also less dependent on a remarkable person continuing to be remarkable. Ordinary officials can do useful work within a well-designed institution. Ordinary citizens can correct some errors without becoming full-time political heroes.

Government should make that ordinariness possible. People need to work, rest, raise families, pursue interests, and sometimes pay little attention to politics. A system that requires permanent mobilisation to prevent arbitrary rule is imposing a heavy burden merely to maintain basic freedom.

The railway may still be built, and the homeowner may still regard the decision as wrong. A tolerable government gives reasons, follows rules, provides the remedies it promises, and allows her to speak against it afterwards. It does not demand gratitude for having heard her before taking the land.

Public power will remain necessary and dangerous. We should be suspicious of anyone who finds either half of that sentence sufficient. The task is to give it enough authority to do work we cannot safely leave undone, while preserving a life beyond its direction and a means of saying when it has gone too far.

\chapter{When the Majority Disagrees with Itself}

Three people must choose one project. Anna, Ben, and Clara each understand the alternatives and rank them consistently. There is no misinformation, no corruption, and no confusion about the ballot. We will call the projects A, B, and C so that their merits do not distract us from the decision rule.

Their preferences are as follows. First means most preferred.

\begin{center}
\begin{tabular}{llll}
\toprule
Voter & First & Second & Third \\
\midrule
Anna & A & B & C \\
Ben & B & C & A \\
Clara & C & A & B \\
\bottomrule
\end{tabular}
\end{center}

Ask them to choose between A and B. Anna prefers A; Clara does too. A wins by two votes to one.

Now compare B with C. Anna prefers B; Ben does too. B wins by two votes to one.

Finally compare C with A. Ben prefers C; Clara does too. C wins by two votes to one.

The majority prefers A to B, B to C, and C to A. Every individual has a consistent ranking. The majority comparisons go round in a circle.

This is the Condorcet cycle. Nothing has gone wrong with the voters' arithmetic or sincerity. The difficulty arises when their separate rankings are combined. There is no project in this example that defeats each of the others in a direct majority vote.\footnote{Marquis de Condorcet, \emph{Essai sur l'application de l'analyse à la probabilité des décisions rendues à la pluralité des voix}, 1785. The three-voter example and agenda results are worked out explicitly here; all pairwise outcomes assume sincere voting within each pair.}

\section*{The order changes the winner}

Suppose the rule is to hold two successive votes. The winner of the first contest faces the remaining project. Assume for the moment that everyone votes in each contest for the project they prefer in that pair.

If A faces B first, A wins. It then faces C, which defeats it. The final winner is C.

If B faces C first, B wins. It then faces A, which defeats it. The final winner is A.

If C faces A first, C wins. It then faces B, which defeats it. The final winner is B.

\begin{center}
\begin{tabular}{lll}
\toprule
First contest & First winner & Final winner \\
\midrule
A against B & A & C \\
B against C & B & A \\
C against A & C & B \\
\bottomrule
\end{tabular}
\end{center}

The order has selected among three possible outcomes without anyone changing their underlying preferences. Whoever chooses the agenda has power even without controlling another vote.

That power need not be exercised secretly. A rule adopted openly in advance may specify the sequence. The result can be legitimate under the agreed procedure. What it cannot truthfully be called is the discovery of one unambiguous group ranking that existed independently of the procedure.

This distinction matters in committees and legislatures. The order of amendments, the choice of which proposals reach a vote, and the alternative against which a proposal is compared can affect the outcome. Procedural decisions deserve attention because they are part of the substantive decision.

If voters anticipate the final round, they may also vote strategically in the first. The simple table above assumes sincere pairwise voting precisely to isolate the effect of order. Adding strategy changes the analysis; it does not restore a magical public mind beneath it.

\section*{No one has contradicted herself}

The cycle can initially feel like an error because we are accustomed to describing a group as though it were one person. ``The committee prefers A'' sounds similar to ``Anna prefers A.'' But the majority supporting one comparison is not the majority supporting another.

Anna and Clara defeat B with A. Anna and Ben defeat C with B. Ben and Clara defeat A with C. The coalition changes. Combining the results as though one continuing individual had made all three judgements creates the impression of a single contradictory mind.

The problem is therefore not cured by demanding that voters become more rational. We gave them consistent preferences at the start. Nor does better information necessarily remove it. They may fully understand the consequences and value them differently.

Deliberation might change their preferences or produce a new compromise. That can be useful. But it is another possible process, not a guarantee that the cycle must disappear. People can talk honestly and continue to disagree.

The cycle is possible, not inevitable. Many sets of preferences produce a clear majority winner. Some restrictions on how preferences are arranged can prevent cycles. We should not turn a mathematical possibility into a claim that every actual election is incoherent.

Its importance lies in the guarantee it defeats. Majority rule does not always recover a complete, consistent ranking from individually consistent rankings. A political system must therefore do more than promise to reveal what the people already want as one thing.

\section*{Arrow's harder question}

Perhaps a different rule could avoid the difficulty. Instead of successive majority votes, collect every ranking and use a better method to produce a social ranking. Kenneth Arrow examined whether such a method could satisfy a set of appealing conditions at once.\footnote{Kenneth J. Arrow, \href{https://www.journals.uchicago.edu/doi/abs/10.1086/256963}{``A Difficulty in the Concept of Social Welfare''}, \emph{Journal of Political Economy} 58(4), 1950, pp. 328--346; \emph{Social Choice and Individual Values}, second edition, 1963. The chapter gives the standard ordinal-ranking formulation rather than a proof.}

First, it should work for every possible collection of individual rankings. We should not have to reject the voters' preferences because they happen to be inconvenient for the rule.

Second, the result should be a complete and consistent social ordering. The procedure should be able to compare the alternatives, and its comparisons should not run round in circles.

Third, if everyone prefers A to B, the social ranking should put A above B. This is a unanimity requirement, often called the Pareto condition in this setting.

Fourth, the social comparison between A and B should depend on how individuals compare A and B. If those comparisons remain unchanged, changing their views about C should not reverse the social comparison of A with B. Arrow called this independence of irrelevant alternatives.

Finally, there should be no single individual whose strict preferences always determine the corresponding social comparisons regardless of everyone else. That is the non-dictatorship condition.

With at least three alternatives and the usual finite set of voters, no aggregation rule producing such an ordering can satisfy all these requirements over unrestricted preferences. Something has to give.

The result does not say that no election can be fair, that democracy is impossible, or that a dictator would govern well. It identifies an incompatible set of requirements for a particular kind of aggregation. Its force depends on preserving the requirements, not inflating the conclusion.

\section*{Which requirement would we relax?}

Different procedures accept different compromises. A system may use a rule that does not satisfy Arrow's independence condition. It may restrict the kinds of preferences or alternatives considered. It may allow an incomplete ordering rather than demand a complete one. It may use information beyond rankings, such as scores or comparisons of welfare, and then face questions about what that information means and how it is compared across people.

None of these moves is automatically illegitimate. The point is to say which move has been made and why it suits the decision. A theorem tells us that a promised combination is unavailable. It does not choose the best feasible arrangement for a town, parliament, or association.

Consider the independence condition. It has an intuitive appeal: why should a changed view of C alter the result between A and B when no one has changed their comparison of those two? Yet some ranking methods use the positions of all alternatives to assign points. Their treatment of A and B can therefore depend on where C is placed.

One may judge that acceptable in a particular setting because the method uses more of the ranking information. Another may regard it as an unacceptable vulnerability. That is a real trade-off to discuss, not evidence that the first person has failed to understand democracy.

Likewise, restricting preferences may be reasonable for a narrow decision where alternatives lie along a single dimension and each voter ranks options lower as they move away from a preferred point on either side. These are called single-peaked preferences. It is far less plausible as a general description of politics, where tax, religion, security, environment, and public services interact.

The practical value of the theorem is to make grand claims about a perfect rule harder to sustain. A reform can improve representation or reduce a known defect without becoming a device that satisfies every appealing principle.

\section*{The dictator's consistency}

In the formal setting, using one person's consistent ranking supplies a consistent social ranking. This sometimes provokes the claim that mathematics favours dictatorship. It does nothing of the kind.

Consistency is only one property. A ruler can consistently prefer his own enrichment, punish opponents, and ignore other people's losses. The ranking can be orderly while the government is cruel or incompetent. The theorem does not establish that the chosen individual has good information, sound values, or a willingness to leave office.

The formal term \emph{dictator} describes a decisive individual's role in aggregation. Actual dictatorship involves institutions, coercion, information, and succession. Treating the mathematical term as a complete political recommendation confuses levels of argument.

We accept many difficulties in collective choice because giving one person final authority creates difficulties we judge worse. This is a comparative judgement about how people can live under power. No logical theorem relieves us of making it.

The same applies to the imagined benevolent expert. The expert's ranking may be consistent. That tells us neither whether the ranking should prevail nor how to protect those harmed by it. A coherent decision is not necessarily a legitimate one.

\section*{Voting for the second choice}

Strategy creates another gap between ballots and preferences. Suppose a voter prefers A, believes A has little chance, and fears C. She may vote for B because B is more likely to defeat C. The ballot records support for B without showing that B was her favourite.

Calling this dishonest can be unfair. The voter is responding to the rule and the choices available. If sincere expression makes her least preferred outcome more likely, she has a reason to compromise. The system has asked one mark to perform both expression and influence.

Allan Gibbard and Mark Satterthwaite established a general result about deterministic rules selecting a single winner. With at least three possible winners, unrestricted strict preferences, and every alternative able to win, a non-dictatorial rule cannot guarantee that truthful reporting is always a best strategy for every voter. There will be some circumstances in which misreporting can improve an outcome for someone.\footnote{Allan Gibbard, \href{https://www.eecs.harvard.edu/cs286r/courses/fall11/papers/Gibbard73.pdf}{``Manipulation of Voting Schemes: A General Result''}, \emph{Econometrica} 41(4), 1973, pp. 587--601; Mark A. Satterthwaite, ``Strategy-Proofness and Arrow's Conditions,'' \emph{Journal of Economic Theory} 10(2), 1975, pp. 187--217.}

The assumptions matter. The result is not a claim that every voter can manipulate every election, or that every tactical attempt will succeed. A voter may lack reliable information about others. A manipulation may be difficult to identify or coordinate. The theorem concerns the absence of a universal guarantee against it within the specified class of rules.

Ranked ballots can reduce some familiar incentives associated with vote splitting. They do not abolish every strategic possibility. Approval ballots offer another form of expression, but deciding which candidates to approve can depend on expectations about the contest. Every rule requires a clear account of what the ballot means.

This should make reformers precise rather than discouraged. A system can be less vulnerable to a particularly damaging problem without being invulnerable to all strategy. The claim worth defending is the specific improvement.

\section*{A majority for each promise}

Even without a formal cycle, public opinion can support separately attractive proposals whose combination is difficult. Voters may favour better services, lower taxes, less debt, and generous exemptions. Different majorities may support each. A government then claims a mandate to fulfil the whole set, as if the resources had been supplied along with the votes.

The difficulty can be reduced by presenting packages and consequences clearly. If one option includes a tax increase, that should be visible when support is measured. But a package creates another problem: a vote for it does not show agreement with every component. People may accept one policy reluctantly to obtain another.

Election winners often interpret that reluctant acceptance as enthusiastic endorsement. They have incentives to describe a narrow or mixed victory as an unambiguous instruction. The losers may make the opposite mistake, treating every voter for the winning party as committed to its worst proposal.

The ballots cannot carry all the psychological information assigned to them afterwards. A count establishes a result under a rule. Additional claims about motives require additional evidence.

This restraint matters for liberty. A government should not claim permission to enter every part of life because it won a contest that offered voters only a few bundles of policies and personalities. Constitutional limits protect people partly against this expansion of a limited electoral result into comprehensive authority.

\section*{Procedure has to be chosen too}

There is an apparent regress. If a voting rule affects the outcome, how should the rule itself be chosen? Another vote uses another rule. No procedure can stand entirely outside the political order it helps create.

This does not make all rules equivalent. We can assess transparency, stability, inclusiveness, resistance to manipulation, representation, accountability, and the treatment of minorities. We can examine how a rule performs in practice and whether participants can understand it. We can adopt changes through procedures that seek broad consent and avoid immediate partisan advantage.

Choosing rules before knowing which side benefits in a particular contest can help. So can requiring wider agreement for changes to the rules than for ordinary policy. These arrangements are imperfect, but they reduce the temptation for a temporary majority to redesign the competition solely to remain in office.

The fact that a rule requires judgement is no reason to abandon judgement. It is a reason to make the choice visible. Pretending a procedure is merely technical can conceal the values it privileges and the interests it serves.

There is also value in stability. Changing rules after every disappointing result can make losers suspect that no agreed contest will ever be final. Stability should not protect a seriously defective system forever, but reform needs a case stronger than the fact that one's preferred side lost.

\section*{What a vote can reasonably promise}

A vote can settle which proposal will be acted upon under an accepted procedure. It can distribute offices, reveal support, and permit peaceful alternation. It can leave a public record against which later claims are assessed. These functions are enough to make voting valuable.

It cannot promise that the result is wise, that every voter wanted the same thing, or that another defensible procedure would necessarily have produced the same outcome. It cannot make the losing side's interests vanish. It cannot convert a majority into a person whose single intention governs every future question.

Understanding these limits need not weaken commitment to democracy. It can make the commitment less sentimental. We defend the practice because of what it enables among people who disagree, not because arithmetic will turn disagreement into a flawless common will.

Anna, Ben, and Clara still need to choose a project. They can agree on a procedure, accept its result for the present purpose, and retain the right to discuss what happens afterwards. One will obtain a favourite outcome; others will not. Nothing in the process requires them to pretend that the winning project was secretly everyone's first choice.

Large societies face a more complicated version of this ordinary difficulty. Their achievement is not that everyone becomes satisfied. It is that dissatisfaction can remain lawful, intelligible, and capable of influencing the next decision.

\chapter{The Voter Has Other Things to Do}

A voter comes home from work, makes a meal, deals with a child's school problem, and opens an account that needs paying. Somewhere in the remaining evening she is expected to understand fiscal policy, electricity markets, migration law, pension finance, military strategy, and the records of the people competing to govern her.

She cannot do all of this well. Neither can most of the people who condemn her for failing. Political commentators specialise, rely on sources, and make mistakes. The ordinary voter is asked to reach a conclusion across many fields while having much less time to examine them.

This is a problem for democracy, but it is not an astonishing discovery of stupidity. It is a consequence of scale, complexity, and the division of attention. A serious political arrangement must work with citizens who have lives beyond politics.

The alternative fantasy is a population of permanently attentive public reasoners. They read every bill, inspect each number, remember every promise, and correct themselves without vanity. Such citizens would be useful. Institutions should not depend upon their existence.

\section*{Why knowing more can be costly}

Anthony Downs drew attention to the incentives facing an individual voter. In a large election, one vote is unlikely to determine the winner. Learning more takes time. The private return from mastering a complicated issue can therefore seem small relative to the effort, even when the collective consequences of widespread ignorance are serious.\footnote{Anthony Downs, \emph{An Economic Theory of Democracy}, 1957; see also \href{https://doi.org/10.1086/257897}{``An Economic Theory of Political Action in a Democracy''}, \emph{Journal of Political Economy} 65(2), 1957, pp. 135--150.}

This is often called rational ignorance. The phrase does not mean that an uninformed belief is true, or that failing to learn is admirable. It describes how limited attention can make sense for an individual while producing an unsatisfactory collective result.

The same voter may investigate a used car carefully because her own decision directly determines what she buys. She may know much less about a national transport programme because her individual ballot has little influence on it. The difference need not reflect a deeper love of cars than of the country. It reflects the relationship between effort and effect.

People nevertheless vote and learn about politics for many reasons. They care about others, feel obligations, enjoy discussion, identify with a party, or value participation itself. A calculation based solely on the chance of casting the decisive vote leaves out much of what people think they are doing.

Adding these motives does not remove the information problem. A person can vote dutifully and know little about the policy. Another can know a great deal and use the knowledge chiefly to defend a prior allegiance. Attention, information, and sound judgement are related but distinct.

The institutional question is how to make a limited amount of attention more useful, and how to prevent one poorly informed decision from transferring power that cannot later be corrected.

\section*{The shortcuts we cannot avoid}

Voters use party labels, endorsements, trusted journalists, associations, and personal experience. These shortcuts can be sensible. A party's record may provide more information than a hurried attempt to understand every line of its programme. A professional association may identify a technical concern a voter would otherwise miss.

The usefulness depends on whether the source is reliable for the question. An organisation that represents a trade may know a great deal about practical conditions and still have an interest in restricting competition. A newspaper may investigate corruption rigorously while framing economic issues through a narrow set of assumptions. A trusted friend may have excellent judgement in one area and repeat rumours in another.

The voter does not need to reject all mediation. She needs some ability to compare sources, notice incentives, and recognise when a source has crossed beyond its competence. That is demanding, but less impossible than independently reproducing every relevant inquiry.

Institutions can make this work easier or harder. Clear public records, intelligible explanations, comparable statistics, and accessible audits reduce the effort needed to check a claim. Documents released in unusable form or important changes buried in technical material raise it. Transparency is not achieved merely by making information technically obtainable.

The person with a free evening should be able to learn something consequential without becoming an unpaid investigator. If only specialists with extensive resources can understand what the government has done, the information may be public while accountability remains narrow.

\section*{A party is a useful compromise}

Parties gather positions, recruit candidates, organise legislative work, and provide a continuing identity against which a record can be judged. Without them, voters might face an even less manageable collection of individuals and promises.

The package also creates constraints. A voter may support a party's economic policy and dislike its position on civil liberties. Another may reverse those preferences. Both cast the same ballot for different reasons. Internal party debate can help reconcile such differences, but it can also be suppressed in the name of unity.

A party becomes dangerous when loyalty to the organisation displaces judgement about its actions. A corruption allegation is dismissed because opponents raised it. A policy once condemned becomes acceptable when proposed by one's own leaders. The voter learns not what happened, but which response membership requires.

This behaviour is not confined to people with little education. Knowledge can supply better excuses. A sophisticated supporter may be able to explain away conduct that a less invested observer recognises immediately. The intelligence is real; it has been given a defensive task.

There are moments when a supporter can test the attachment. Would the same action be acceptable if the other party did it? What evidence would justify criticism of the current leader? Is disagreement possible without exclusion from the social group? These questions are uncomfortable because political membership often includes friendships and a sense of personal decency.

The answer need not be withdrawal from every party. Parties can be useful while their members remain capable of refusal. An organisation that can survive internal criticism is less dependent on the fiction that its errors are always smaller than the errors of its enemies.

\section*{Why more voices need not mean more knowledge}

A large group can combine information that no one person possesses. The classic jury argument associated with Condorcet shows how, under demanding assumptions, independent judgements that are individually more likely than not to be correct can produce increasingly reliable majority decisions.

The assumptions are doing work. There must be a correct alternative of the relevant kind, and the errors must have sufficient independence. If everyone receives the same false report, increasing the number of people repeating it does not create new evidence.

The point can be seen without a theorem. Ten witnesses who independently saw an event offer a different kind of evidence from ten people who read one account of it. Counting ten voices in both cases ignores the source of their agreement.

Political disagreement also includes values, not only questions with one objectively correct answer. People can agree on the expected cost of a policy and disagree about whether the benefit justifies it. A majority may settle which choice is made; it has not thereby discovered a fact in the same sense as identifying the correct answer to an arithmetic problem.

This should restrain both praise and contempt of crowds. Distributed knowledge is valuable where institutions allow it to enter the decision. Conformity, common misinformation, and incentives to conceal disagreement can prevent the expected benefit. The result depends on how people learn and speak, not merely on how many there are.

Independent channels matter for this reason. A government should not be able to turn one official error into the only account available. Opposition, investigative reporting, professional inquiry, and local experience can expose different parts of a failure. Their disagreement may look untidy while preserving information a tidier system would lose.

\section*{Blame is easier than attribution}

Voters often judge governments by visible outcomes: prices, employment, safety, waiting times, and public services. This is reasonable. A government should answer for what it does. The difficulty is determining which outcomes it caused or could have prevented.

A global shock can raise prices under a competent government. A favourable external change can improve conditions under an incompetent one. Infrastructure completed today may have been planned years earlier. A neglected system can function for a time before its maintenance failure becomes visible under a successor.

Politicians exploit these lags. They claim favourable outcomes as achievements and unfavourable ones as inheritances. Sometimes the explanation is correct. The public needs a record capable of distinguishing a genuine inheritance from a convenient excuse.

This is another reason to preserve forecasts, timetables, and original budgets. The question is what the government knew, what it promised, what it did, and how conditions changed. A judgement based only on the final number can reward luck and punish competence. A judgement that accepts every explanation from the incumbent can make accountability impossible.

No voter will perform a complete causal analysis. Public institutions and journalism can divide the work. The quality of that division matters more than exhorting every citizen to become an expert in every cause of inflation or every stage of a railway project.

\section*{The complaint about other voters}

Criticism of democracy often contains a revealing asymmetry. The speaker's own imperfect knowledge is a manageable limitation. Other people's imperfect knowledge is a reason they should have less power.

There are serious questions about political competence. A vote helps authorise coercion over others, and ignorance can have consequences. Yet any proposal to distribute voting power according to knowledge must answer who defines the test, who administers it, and how it will resist capture by those whose knowledge is rewarded.

A test can measure familiarity with institutions or facts without measuring judgement, honesty, or awareness of one's limits. It can privilege a form of education associated with class and then present the result as politically neutral. Specialists can share blind spots that people outside their profession notice.

These objections do not prove that public ignorance is harmless. They show that concentrating political power among those certified as knowledgeable creates another problem that must be compared with it. The proposed gatekeepers are people with interests, not an external intelligence unaffected by politics.

There are less sweeping ways to improve decisions: better civic education, clear information, expert testimony open to questioning, deliberative bodies, and public scrutiny of claims. Each can help without requiring a permanent class whose political standing exceeds that of others.

The burden of proof should be especially high for a reform that removes people from equal participation. Its author must explain not only what is wrong with current voters but why the proposed selectors will perform better under the power they acquire.

\section*{Interests do not disappear with education}

James Madison treated faction as a recurring feature of a free society. People form groups around interests and passions, some of which can threaten others' rights or the common good. Eliminating the causes would require measures more destructive than managing the effects. His proposed arrangements sought to make domination harder by multiplying and checking interests.\footnote{James Madison, \emph{The Federalist}, No. 10, 1787. The argument concerns the control of faction's effects, not the elimination of disagreement or a guarantee that a large republic will always prevent domination.}

We need not treat that design as sufficient for every later circumstance to retain the insight. Politics contains conflicting interests because people have different positions in life. A landlord and tenant can understand the same policy and want different outcomes. A worker in a protected industry and a consumer paying more for its goods may not be reconciled by another lesson in economics.

Education can clarify the conflict. It cannot promise to remove it. Nor should every organised interest be dismissed as corrupt. Workers, patients, businesses, tenants, and local residents need ways to make knowledge and claims available to government. The question is when organisation becomes disproportionate control.

An interest can be legitimate and still overstate its case. A union may protect workers while defending a practice costly to outsiders. A business may provide useful employment while seeking rules that exclude competitors. A residents' group may care about a place while making it impossible for others to live there. The public argument must include those who do not belong to the organised group.

The next chapter examines the institutional consequences of these asymmetries. Here they show why better-informed politics will still be politics. People will continue to bargain over resources, rules, and risks even after factual mistakes are corrected.

\section*{A citizen need not live in permanent alarm}

There is a political style that treats every day as an emergency and every pause in attention as complicity. The citizen must monitor each statement, condemn each offence, and remain available for the next mobilisation. The demand is unsustainable for most people and may leave them less capable of careful judgement.

Some emergencies are real. Some moments require unusual attention and action. If the language is used continuously, it becomes harder to distinguish them. The public can become exhausted, habituated, or dependent on the people who announce which outrage deserves its limited energy.

A healthy political order should allow periods in which ordinary people can attend to ordinary life. This does not require indifference. It requires reliable institutions that continue to function without every citizen watching every official every hour.

Voting, local participation, supporting trustworthy organisations, and paying attention to consequential decisions can be meaningful without becoming a total identity. People have different capacities and obligations. The parent caring for a sick child is not failing citizenship because she has not read the latest argument before midnight.

The more comprehensive the powers at stake, the more dangerous limited attention becomes. That is an argument for limits on power as well as for better information. A single election should not place every channel of correction under the winner's control.

\section*{Making attention count}

The ordinary voter should be neither flattered as the source of automatic wisdom nor despised as an obstacle to government by the deserving. She is one person among many, with limited time and a legitimate stake in the decisions made over her life.

Institutions can make her attention more effective by providing clear choices, credible records, contestable claims, and identifiable responsibility. They can make it less effective by blurring responsibility, concealing costs, and turning every disagreement into a test of belonging.

The task is not to eliminate shortcuts but to make better shortcuts available. A trustworthy audit, a well-supported investigation, or a representative with a record of explaining decisions can save thousands of citizens from duplicating the same work. Trust remains conditional, but it need not begin from zero each morning.

There will still be mistakes. Citizens will reward the wrong people, overlook a danger, or become attached to an explanation that later fails. A defensible democracy gives them another opportunity without requiring them to confess that their earlier lives were worthless. It also prevents the beneficiaries of the mistake from closing the process before that opportunity arrives.

Our voter will eventually put the documents aside. The meal needs clearing away, and the child may need help again. Government exists partly so that a person can live this life with some security. If preserving it requires an impossible standard of political vigilance, the fault is not hers alone.

\chapter{Who Controls the People in Charge}

A representative announces that she has secured funding for a new facility in her district. Local people applaud. There will be jobs during construction, and the completed building will be visible evidence that someone has paid attention to them.

Elsewhere, taxpayers each bear a small part of the cost. They may never hear of the project. The people whose more useful project was not funded may not know why theirs lost. The representative receives concentrated credit; the alternatives and costs remain dispersed.

The facility may still be worth building. This imagined case does not establish corruption. It identifies an incentive that can operate without a bribe or a secret agreement. A politician is rewarded for bringing something home while the wider public has difficulty seeing what was given up to obtain it.

The question of control begins here, in arrangements where ordinary people pursuing ordinary incentives can produce a bad result without anyone needing to consider themselves dishonest.

\section*{Everybody's worthwhile project}

Suppose several representatives each want a local project. Each believes the others' proposals are less important, but agrees to support them in exchange for support for her own. The package passes.

Such bargaining can be useful. Different districts have different needs, and agreement often requires exchange. A legislature in which nobody would support anything outside their own constituency could fail to provide valuable common infrastructure.

The exchange can also assemble a package whose total cost exceeds its benefit. Each representative sees the local gain and a fraction of the national cost. Each can tell constituents that refusing the bargain would merely have left them paying for everyone else's project without receiving one themselves.

This is one reason public waste can persist under representatives who are attentive to their voters. Local responsiveness and sound collective allocation are not identical. The mechanism belongs to the broader problem of concentrated benefits and dispersed costs explored by public-choice theorists.\footnote{Mancur Olson, \emph{The Logic of Collective Action}, 1965; James M. Buchanan and Gordon Tullock, \emph{The Calculus of Consent}, 1962. The opening project and bargaining examples are illustrative rather than reports of a particular allocation.}

The answer is not to assume that central officials can identify every local need better. It is to make proposals and comparisons visible. What problem does the facility solve? How many people will use it? What are the lifetime operating costs? Which alternatives were considered? Who will pay when the opening grant is spent?

A building can be financed once and burden a local budget for decades. The politician who wins the construction money may be praised for generosity while successors struggle to heat, staff, and maintain the result. The photograph at the opening does not show that continuing commitment.

\section*{The rule written by its beneficiary}

Political advantage need not arrive as a payment. It can arrive as a regulation. A professional association proposes a licensing requirement in the name of quality. The requirement may protect customers. It may also make entry harder for competitors and preserve the incomes of existing members.

Both effects can occur together. The association's expertise is relevant, and its interest is relevant. Rejecting its advice solely because it benefits would be careless; accepting it solely because it invokes safety would be equally careless.

George Stigler's work on economic regulation directed attention towards the ability of organised industries to obtain rules serving their interests.\footnote{George J. Stigler, ``The Theory of Economic Regulation,'' \emph{Bell Journal of Economics and Management Science} 2(1), 1971, pp. 3--21, doi:10.2307/3003160.} The broader lesson is that regulation should be examined through its actual beneficiaries and incentives, not only through the public purpose written at the top.

An incumbent firm may find a complex standard easy to satisfy because it helped design the standard around its own process. A new firm may face large costs before serving a single customer. A subsidy may be presented as support for innovation while repeatedly flowing to companies with the best relationships. A delay imposed on competitors can be more valuable than a direct grant.

The public often sees only the stated purpose because understanding the detail is expensive. The beneficiary has a strong reason to master it. This asymmetry can shape policy even when every meeting is lawful and every participant speaks politely about the common good.

It is not enough to ask whether officials took money. We should ask whose problem became the government's priority, whose evidence defined the options, and whose alternatives never received a hearing.

\section*{Access before persuasion}

Lobbying is not inherently corrupt. People affected by rules need ways to explain how those rules operate. A patient group, trade union, business, charity, or local association may possess information the government needs.

The problem is unequal access and influence that cannot be inspected. Some groups can afford permanent staff, detailed submissions, legal advice, and repeated contact. Others appear once, badly prepared, after the important choices have already been narrowed. The formal invitation to comment may be equal while the practical opportunity is not.

Officials can become dependent on a narrow circle for information. Over time, its language defines what counts as realistic. A proposal outside the circle sounds naive partly because the people receiving it have spent years hearing the same alternatives from people they know.

This does not require a conspiracy. Professional familiarity, recruitment, shared education, and future employment prospects can produce a durable alignment. Everyone may regard themselves as responsible while treating outsiders as a nuisance to be managed.

Disclosure helps, but disclosure must be usable. A list of meetings without the subjects discussed or proposals supplied may reveal little. Thousands of pages released without a way to identify the decision can overwhelm rather than inform. The relevant record connects contact, proposal, reasoning, and outcome.

Restrictions on conflicts of interest and movement between public office and regulated industries can also help where the risk is serious. They need to be designed so that public service can still recruit capable people without permitting decisions to become auditions for later employment. No slogan about the revolving door settles the appropriate rule in every role.

\section*{Procurement without a friend}

Public purchasing is a practical test of institutional quality. A firm should be able to compete without a political patron. The authority should be able to compare offers on relevant grounds and explain its choice. Poor performance should have consequences that do not depend on the supplier's connections.

Rigid procedure can create waste too. A buyer may be forced to choose a superficially cheap offer that ignores maintenance or quality. Fear of challenge may reward a paperwork-compliant bid over a better solution. The cure for favouritism is not necessarily a process so inflexible that nobody can exercise judgement.

The better arrangement defines criteria, records reasons, checks conflicts, and provides review. It allows discretion where needed while making that discretion answerable. An official should not have to choose between blind compliance and an unrecorded favour.

Small suppliers may need a realistic route into the process. Requirements designed around the largest firms can turn competition into a ceremony among incumbents. At the same time, a public buyer needs evidence that a supplier can perform. The distinction between a necessary safeguard and an unnecessary barrier should be open to scrutiny.

The test is what happens when an unconnected firm challenges a decision. Can it obtain the record? Is review timely enough to matter? Will criticism jeopardise future bids? An institution can publish admirable principles while its participants know that objections are punished informally.

\section*{Who watches the independent body?}

Independent regulators, auditors, courts, and statistical offices can protect the public from immediate political interests. Independence is valuable where the work requires resisting pressure from the government of the day.

Yet an independent body is not outside human incentives. It can protect its reputation, resist correction, acquire a narrow professional culture, or become excessively close to the sector it oversees. Calling it independent tells us something about its formal relationship with government, not everything about its judgement.

Accountability should therefore match its role. Methods can be published, reasons given, conflicts disclosed, decisions appealed, and performance reviewed without making every result subject to ministerial preference. The point is to preserve independence from improper direction while avoiding immunity from criticism.

Appointments matter. A body may remain legally independent while appointments gradually ensure that its members understand loyalty as their unstated qualification. Budgetary pressure can produce similar effects. An institution need not be formally abolished to lose the capacity that justified its existence.

The public should ask what makes an office capable of refusing the person who appointed it. Security of tenure may help; so may transparent appointment procedures and more than one institution sharing the decision. Each arrangement has costs and vulnerabilities. What matters is that the refusal remain possible in practice.

\section*{Selecting citizens by lot}

Elections select people able and willing to compete for office. This favours some combinations of confidence, time, resources, networks, and organisational support. It need not produce a body resembling the population in its daily experience.

Selection by lot offers a different route. A group of citizens is chosen through a random process, often with sampling designed to reflect broad population characteristics. They receive time, information, and opportunities to question witnesses before deliberating on a specified issue. Modern citizens' assemblies and related processes use this approach in various forms.\footnote{OECD, \href{https://www.oecd.org/en/publications/innovative-citizen-participation-and-new-democratic-institutions_339306da-en.html}{\emph{Innovative Citizen Participation and New Democratic Institutions: Catching the Deliberative Wave}}, 2020, doi:10.1787/339306da-en. The book's discussion evaluates the institutional trade-offs rather than assuming that every process using the label citizens' assembly has the same design or quality.}

The attraction is that participants need not build a political career or satisfy donors. People who would never stand for election can contribute practical knowledge and judgement. Given time and a well-run process, ordinary citizens can engage seriously with complex questions.

Random selection does not make participants wise or disinterested. It changes the path by which they arrive. They still need information, and someone must choose how the question is framed, which witnesses appear, and what materials are supplied. Power can move from campaign organisers to process organisers.

Participation also depends on whether selected people can afford to attend. Payment, childcare, accessibility, and protection from employment consequences matter. A nominally random invitation followed by unequal ability to accept can produce a less representative group than the headline suggests.

The assembly should have a clear task and an honest account of its authority. Is it advising, proposing, reviewing, or deciding? What must elected officials do with its recommendations? If the government promises to listen and then quietly shelves an unwelcome result, the exercise may increase distrust rather than repair it.

\section*{What the assembly cannot replace}

A citizens' assembly can produce an informed recommendation without making the whole population informed in the same way. Participants have spent time learning that others have not. Their changed views may be reasonable, but the wider public needs an account of the evidence and reasoning rather than a demand to defer to a miniature version of itself.

The participants also leave. They may not bear electoral responsibility for later results. This is one reason limited, clearly defined roles can be more defensible than treating sortition as a complete replacement for elected government.

A mixed arrangement can use different strengths. An elected legislature remains responsible for a decision. A deliberative body examines a difficult issue and makes its reasoning public. An independent audit assesses implementation. Courts address legal limits. None is infallible; each can expose something another misses.

The danger is creating so many bodies that nobody can identify responsibility. Every institution consults another, and each points elsewhere when the outcome fails. Adding participation should clarify the decision, not bury it under a succession of processes.

A useful design states who must answer at the end. Advice can be rejected, but the rejection should be explained. Evidence can be disputed, but the dispute should be visible. The public should be able to follow the route from proposal to action without needing personal access to the participants.

\section*{Transparency with consequences}

Exposing a failure does not guarantee its correction. A report can be published, criticised for a few days, and forgotten. The officials responsible may remain in place because no institution has authority or incentive to act.

This is why accountability requires more than information. It requires someone able to demand a response, change a rule, recover money where appropriate, impose a consequence, or remove an officeholder. The process must distinguish error from misconduct, and an unpopular judgement from an abuse of power.

Punishing every mistake can make officials conceal uncertainty and avoid worthwhile experiments. Excusing every failure as complexity can make incompetence permanent. A fair assessment asks what was reasonably knowable, which duties applied, what choices were made, and whether information was honestly reported.

The official who identifies a failing programme early should not be treated worse than the official who hides it until after promotion. Yet institutions can create exactly that incentive when reported problems are counted as personal failures and silent deterioration is invisible.

Protecting people who report wrongdoing or serious failure is therefore part of preserving the institution's ability to learn. Protection should include practical safeguards against retaliation, not only a policy on a website. A person who loses work through informal exclusion has not been protected merely because no formal disciplinary action occurred.

\section*{Power that can be interrupted}

No arrangement will eliminate influence, bargaining, or self-interest. The ambition to remove them completely can itself justify concentrated authority: one leader, party, or commission is entrusted to cleanse the system and becomes difficult to scrutinise because scrutiny is described as defence of corruption.

A more credible ambition is to interrupt durable advantages. Make appointments contestable, decisions inspectable, contracts competitive, and offices replaceable. Allow new organisations to enter public argument. Preserve more than one route through which a failure can become known.

These measures will irritate capable people as well as obstruct dishonest ones. A public official may have to explain a decision to critics who understand it poorly. A firm may lose a familiar arrangement that was convenient for both sides. Some delay and inconvenience are part of making power answerable.

They should not become excuses for paralysis. Procedures can be simplified while reasons remain public. Review can be bounded while rights remain real. The task is to reduce arbitrary power without making effective work impossible.

The representative may still secure her local facility. If its case is strong, it should survive comparison with alternatives. If it is weak, the public should be able to discover that before the building becomes a monument to a successful career. Control means that the person who brings home a benefit also has to answer for the cost left elsewhere.

\chapter{Freedom for People We Dislike}

Imagine that a public library lets local associations hire its meeting room. A group whose opinions you despise makes a booking. Its members are not planning an assault. They intend to give speeches, collect subscriptions, and persuade people that the country would be better if some of your convictions disappeared from public life.

You would prefer them to meet somewhere else. You would prefer that they had no audience. When somebody proposes cancelling the booking, the proposal feels less like censorship than a modest improvement to the evening.

The group then points out that the library has hosted your association. You explain why the cases differ. Your association defends freedom. Theirs abuses it.

Perhaps there is a relevant difference. A meeting can be used to organise threats, identify targets, or obstruct other people's political activity. But the fact that you can describe your own cause generously does not establish that difference. Almost any faction can explain why its opponents are dangerous and its own success is a public service.

Freedom becomes inconvenient at exactly this point. A rule that protects only the people we consider responsible allows us to congratulate ourselves on generosity while retaining the power to decide who deserves it. The people who need protection from us are still waiting outside.

\section*{Having a right does not make a person admirable}

We sometimes speak as though defending someone's liberty requires a certificate of good character. Before opposing a dismissal, a prohibition, or an intrusive investigation, we establish that the person is decent, misunderstood, or vulnerable. This may win sympathy. It also leaves the indecent, correctly understood, and unsympathetic person without much of an argument.

There are people whose beliefs are repellent and whose conduct, within the law, is mean. They may publish foolish books, reject sound evidence, or organise against measures we regard as humane. Their freedom does not depend on our discovering a hidden kindness in them.

Equal protection has force because it interrupts the urge to distribute rights according to affection. A fair hearing remains due to someone whose defence is likely to fail. Privacy can protect an embarrassing life. Freedom of religion includes beliefs that an outsider considers absurd. Freedom to leave a religion includes a departure that believers regard as betrayal.

This does not require anyone to offer friendship, admiration, or a private platform. A newspaper can choose its contributors. A club can have a purpose. A person can decline an invitation. The difficult question is which power is being exercised, under what rules, and with what consequences for the person excluded.

A public library applying a general hiring policy is in a different position from a private host choosing dinner guests. An employer deciding whether someone can do a job is in a different position from a friend ending a friendship. Calling every refusal censorship confuses these cases. Calling every exclusion a private choice can conceal institutions whose decisions govern access to an ordinary livelihood.

We need to examine the relationship rather than settle the dispute by choosing a flattering noun. A small discussion group cannot owe the whole world a microphone. A public authority cannot justify denying a common service by saying that it too has preferences.

\section*{Offence is an unreliable boundary}

Offence tells us something about the person experiencing it. It does not, by itself, tell us which restriction would be justified.

A religious believer may be offended by blasphemy. A former believer may be offended by a sermon describing unbelief as depravity. A republican may find hereditary office humiliating. A monarchist may find a satirical portrait disgusting. If sincere distress were sufficient to prohibit the cause, the most restrictive sensibility could set the terms for everyone.

Nor does the intensity of a reaction settle the issue. People can learn that visible outrage brings concessions. Institutions that reward the threat of disruption may end by allowing the most disruptive opponents to decide which meetings take place. The quiet person receives less protection precisely because he can be relied upon not to make trouble.

There are more specific questions to ask. Has someone threatened violence? Is a named person being persistently targeted? Is confidential information being exposed to make intimidation possible? Is an employee being required to endure conduct that prevents her from doing her work? These questions identify acts and relationships that the general word \emph{offence} obscures.

Context matters. Mill's familiar example distinguishes an accusation against corn dealers printed for public discussion from the same accusation addressed to an excited crowd outside a dealer's house.\footnote{John Stuart Mill, \href{https://www.gutenberg.org/files/34901/34901-h/34901-h.htm}{\emph{On Liberty}}, 1859, chapter III. The corn-dealer example concerns the circumstances in which an expression of opinion can become an instigation to harmful action.} Words can form part of an action. Recognising that does not make every unwelcome opinion an assault.

An institution should have to explain the connection. Who is being threatened, by what conduct, and why would a narrower response fail? The answer may justify intervention. It may also reveal that a claim about safety amounts to a wish not to hear a rival case.

This demand for explanation will irritate people who regard their opponents' badness as obvious. Yet obviousness is what every confident censor brings to work. A restriction needs grounds that do more than record the confidence of the person imposing it.

\section*{The useful fragment of Popper}

The paradox of tolerance is attractive to people who have grown tired of tolerating. It appears to offer a philosophical permission slip: intolerance threatens tolerance; therefore the intolerant may be silenced. All that remains is to apply the label.

Popper's actual discussion is less convenient. He warns that unlimited tolerance can allow coercive movements to destroy the conditions of tolerance. He also says that suppressing intolerant philosophies would be unwise while argument and public opinion can keep them in check. He considers movements that refuse argument and teach their followers to answer with force or persecution.\footnote{Karl R. Popper, \emph{The Open Society and Its Enemies}, 1945, volume I, note 4 to chapter 7. The discussion above paraphrases the sequence of Popper's argument; it does not treat the familiar isolated sentence about intolerance as a complete rule for deciding particular cases.}

The warning does not relieve us of the work of judging a case. An opponent who will not be persuaded by us is not thereby an organiser of violence. People may be stubborn, irrational, or devoted to an illiberal ideal without possessing either the intention or the means to coerce their neighbours. A state empowered to suppress everyone it finds impervious to reason will have no shortage of candidates.

There is also a difference between defending the freedom to argue for a political change and agreeing that every change can legitimately be enacted. Someone may advocate hereditary government, religious establishment, or the abolition of elections. The opportunity to express the view does not grant a right to implement it by intimidation or to imprison those who resist it.

Movements can, of course, conceal their intentions. They can use lawful organisations to prepare unlawful coercion. Waiting for the first murder may be a grotesquely late response to a credible threat. Investigation and preventive action can be justified by evidence of preparation, organisation, and capacity. Their grounds should be specific enough to be challenged by someone outside the governing faction.

The phrase \emph{enemy of democracy} is too useful to leave at the unrestricted disposal of a government. An administration may face real enemies of political freedom and still exploit their existence to harass ordinary critics. We have to retain the ability to recognise both activities in the same administration.

Returning to the library, the decision should turn on the relevant rules and evidence. If members have threatened other users, the authority should address that conduct. If the meeting merely promises an evening of objectionable opinions, dislike has not supplied the missing justification. Keeping the room open may be unpleasant. That is part of what the rule costs its supporters.

\section*{A door that can actually be opened}

The freedom to refuse membership matters as much as the freedom to organise. An association can become oppressive when it is easy to enter and ruinous to leave.

Consider an adult who wishes to leave a tightly organised community. The law may permit departure. Her housing, work, friendships, and family relations may nevertheless be concentrated inside it. She can walk away, but the first night outside is not an abstraction. She needs somewhere to sleep.

This does not mean that other members must continue to agree with her, employ her indefinitely, or preserve every friendship. Some losses follow from a genuine change in shared commitments. A government cannot guarantee an unaltered emotional life after a serious disagreement. It can, however, protect her from violence, enforce her property and contractual rights, and provide access to basic services without requiring the community's permission.

Such protections make a legal liberty usable. So can independent education, access to information, and a means of earning money. Where an institution controls every practical route out, its members' continued presence becomes weak evidence of consent.

The same reasoning applies beyond religion. An employer that also controls housing may acquire leverage far beyond the employment contract. A political organisation that distributes jobs can make dissent materially dangerous. A household in which one partner controls all resources can leave the other with a nominal choice and no tolerable next step.

Exit is not a universal remedy. A tenant should not have to move cities to obtain repairs that are already owed. A worker should not have to abandon a career to escape unlawful treatment. Those who cannot leave still need protection, and the departures of the most mobile can reduce pressure to improve conditions for everyone else.

Yet the possibility of departure limits what any one institution can demand. Multiple employers, independent associations, alternative sources of information, and enforceable rights create options. Their value is easily underestimated by someone who has never needed one urgently.

An orderly institution may resent these options. It may describe departure as disloyalty and outside help as interference. Its account deserves no automatic privilege over the account of the person trying to get out.

\section*{A private life without a public defence}

There is another liberty that receives less applause: the freedom to live without making one's choices exemplary.

Someone wants a quiet job and has no ambition to lead. Someone else spends an inheritance travelling. A couple do not want children. A person attends services in a nearly empty church. Another would rather repair an old motorbike than join the local campaign. We may judge these choices foolish, wasteful, or disappointing. That does not establish our authority to reorganise them.

Public argument has a tendency to turn every private decision into a contribution to some aggregate problem. The birth rate is too low, consumption too high, civic participation too weak, the economy insufficiently dynamic. Once a choice is placed in the right column, somebody can explain why a better citizen would choose differently.

Some private conduct does impose costs on others. Pollution, violence, fraud, and neglected legal obligations do not become harmless because the actor calls them personal. But a measurable effect on a social statistic cannot itself create an unlimited claim on a person's life. Otherwise almost no choice would remain outside public direction.

The distinction requires argument about the kind and seriousness of the effect, the rights involved, and the powers a remedy would create. It cannot be settled by observing that society has an interest. Society has an interest in almost everything. That is why limits are needed.

Privacy helps sustain those limits. A person should not have to prove that every concealed fact is innocent before being allowed a concealed life. Private space permits unfinished thoughts, changes of mind, unremarkable eccentricities, and relationships that would suffer under continuous inspection. It also makes wrongdoing harder to discover. We weigh those costs when authorising investigation; we do not abolish the distinction because secrecy can be abused.

A population required to be permanently legible to its rulers will learn to perform legibility. Its members may display correct opinions while withholding anything that could be misunderstood. This is a poor substitute for either honesty or freedom.

\section*{What equality should protect}

Material inequality complicates this account. The wealthy person can often purchase distance from a hostile institution. The poor person may have to remain agreeable. Legal permission to speak is a thin protection if one local employer can end the speaker's livelihood and every realistic alternative.

This is a reason to take bargaining power and access seriously. It does not establish that all differences in possession or outcome must be removed. A society can reduce dependence through social provision, competition, enforceable employment rights, and limits on corruption without requiring everyone to live the same life.

Differences will reappear wherever people can save, give, associate, work different hours, or accept different risks. Some differences will be unjust; some will reflect choices we have reason to permit; many will contain both luck and effort. An administration committed to eliminating every discrepancy would need an increasingly intimate authority to decide which choices count and which must be undone.

The earlier discussion of inheritance showed how quickly a demand for a perfectly level beginning enters the home. A parent reads to a child, spends time teaching, introduces a friend, or provides a quiet room. These advantages matter. So do the relationships through which they arise. Recognising the unfairness of unequal childhoods does not produce a painless method of equalising them.

A workable concern for equality therefore needs a more definite object than the disappearance of all advantage. Equal standing before public institutions, protection from arbitrary power, and a sufficient material basis for refusing abuse give us things to examine and improve. We can ask whether a person can obtain justice, leave a dangerous relationship, or contest a decision without first becoming dependent on a patron.

These aims still involve difficult choices and substantial costs. They are not an excuse to bless whatever distribution happens to exist. They do, however, leave room for people to want different things and to disappoint one another without turning every disappointment into an administrative case.

Freedom will produce lives we would not recommend. If that result is unacceptable, the commitment to freedom has already become conditional on obedience.

\section*{The evening continues}

The library meeting ends. Its audience may have heard nonsense. Some participants may leave more convinced of it than when they arrived. Freedom has not reliably improved their judgement in the course of two hours.

This is a disappointing outcome only if we expected a right to justify itself by producing correct opinions at once. Its more modest achievement is that no faction acquired the routine authority to cancel a rival's meeting simply by announcing that it knew better.

You remain free to answer, organise, publish, or stay home. Staying home need not mean approval. There are evenings when attending to the people at your own table is a reasonable use of a finite life.

The irritating people have not disappeared. They still live in the town, use its rooms, and vote. So do you. A free society leaves much of this dissatisfaction unresolved. The alternative is to give someone the power to resolve it for us, including the dissatisfaction other people feel when they see our name on a booking.

\chapter{Repairs without a Promise of Perfection}

Suppose that a flood cuts a town off from its usual supply route. A bridge needs inspection, residents need temporary accommodation, and equipment must be hired before another night of rain. The ordinary procurement process would take weeks. The mayor authorises emergency expenditure.

This may be exactly what the town needs. Insisting on the usual timetable could protect the purchasing procedure while leaving people in waterlogged houses. Someone must be able to act quickly.

Six months later, the roads are open. The emergency purchasing arrangement continues. A contractor argues that changing it would disrupt recovery. Officials find the reduced paperwork convenient. The mayor says that criticism is an insult to the people who worked through the flood.

Nothing in this imagined sequence requires the original emergency to have been invented. A justified exception can outlive its justification. People acquire contracts, habits, and reputations inside it. Ending the exception then becomes a new political decision with opponents of its own.

The question that should have accompanied the first authorisation was therefore simple: what happens when the water goes down?

\section*{Making the end part of the beginning}

Emergency powers are easier to create than to retire. A measure begins with a vivid threat and ends with an argument about whether the threat has sufficiently receded. The official who ends it may be blamed for anything that goes wrong afterwards. The official who preserves it can continue to invoke caution.

This gives continuation an advantage even when circumstances have changed. It also encourages a peculiar use of uncertainty. At the beginning, uncertainty may support rapid action because waiting is dangerous. Later, uncertainty may support indefinite continuation because nobody can prove that every danger has passed. A power that requires perfect safety before it expires may never expire.

Time limits can reverse part of that burden. When an authorisation ends on a specified date, continuation requires a new decision. The authority should explain which threat remains, what the measure has achieved, what costs it imposes, and why ordinary powers are insufficient. The decision can still be made quickly. It need not be made silently.

A date alone is weak protection if renewal is automatic or if a new declaration restarts the same arrangement without scrutiny. The institution reviewing continuation must have information, time, and some capacity to refuse. Otherwise the sunset is merely a recurring ceremony at which the sun is instructed to remain where it is.

Some consequences also survive the formal expiry. Information collected for an emergency may remain in databases. A temporary restriction may have closed businesses. Equipment may have been purchased under contracts extending for years. A satisfactory ending must address these residues rather than announce that the exceptional legal power has disappeared.

In the flooded town, this would mean reviewing contracts, publishing expenditure, resolving claims, and returning purchasing decisions to ordinary rules. It might also mean retaining a carefully defined emergency procedure for the next flood. Learning from an exception need not mean either preserving it wholesale or pretending it was never needed.

These are principles of institutional design, not a claim that every emergency has the same timetable. War, fire, epidemic disease, and financial disruption present different threats. What they share is the temptation to let urgency answer questions that concern the period after urgency has passed.

\section*{When winning changes the contest}

A more serious difficulty arises when an elected government changes the conditions under which it can be replaced.

It may begin with proposals that can each be defended separately. The broadcasting authority needs reform. The courts are slow. Election administration is inefficient. Public organisations contain political appointees. Some of these complaints may be true. No institution becomes sound merely because a dangerous politician criticises it.

The test is what the remedy does. Does it replace a flawed appointment process with a more open one, or give the governing party control? Does it improve election administration under rules that apply to everyone, or make opposition activity harder? Does it address judicial misconduct through a defensible procedure, or remove judges whose decisions inconvenience the executive?

The cumulative pattern matters. A government can leave the ballot paper intact while weakening the press, punishing donors to rival parties, using public employment to reward loyalty, and applying regulation selectively. Opposition remains legal but grows steadily more difficult to practise. Counting votes fairly does not repair all the unfairness created before polling day.

We should not call every unpopular reform an attack on democracy. The charge becomes useless if it means only that a government is implementing policies we oppose. A tax rise can be bad without making future tax cuts impossible. A planning decision can be objectionable without closing the right to challenge the next one.

Changing the conditions of replacement is different. The public is being deprived of a remedy against the present government's errors. The injury extends beyond the policy disputed today to the future possibility of contesting policy at all.

This is why an election cannot grant unlimited permission. A majority at one moment does not own all later majorities. People change their minds, new voters arrive, and the consequences of a programme become clearer. The winner's authority has to leave room for those developments to matter.

The language of popular sovereignty can conceal this theft of the future. A leader identifies himself with the people and describes every independent obstacle as an attack on them. Once that identification is accepted, his removal begins to look like the removal of the people from their own government. The claim is grandiose. It should be treated as such, including when the leader has won a large and perfectly genuine majority.

\section*{The burden on the leader's supporters}

Opponents will usually object to a concentration of power that disadvantages them. The harder test falls on supporters who expect to benefit from it.

Their government is finally doing things that previous administrations would not do. Its enemies seem smug, obstructive, and undeserving. A procedural shortcut appears a small price for progress. There will be time to restore normal practice after the urgent work is complete.

This is how people with reasonable complaints can help construct an unreasonable power. They do not need to admire arbitrary rule in principle. They need only believe that their own side requires one more exception.

The useful objection comes before the instrument changes hands. If a government gains the power to close an inconvenient newspaper through selective enforcement, the instrument will remain available to its successor. If appointments are converted into rewards for loyalty, the next faction will inherit the method even if it replaces every appointee. A bad precedent does not become harmless because its first use pleases us.

Party colleagues, donors, professional organisations, and sympathetic newspapers therefore have responsibilities that cannot be discharged by pointing to the opposition's faults. They may know those faults intimately. The knowledge does not answer what their own government is doing now.

Breaking with a successful leader can be expensive. A critic may lose office, access, or friends. We should admit that cost rather than present constitutional loyalty as an effortless mark of good character. Institutions need people willing to pay something for their limits. They also need arrangements that reduce the price of doing so.

Secure professional positions, independent sources of income, and protection against retaliation can make resistance less heroic. That is an advantage. A constitution that works only when ordinary officials display exceptional courage has left too much unfinished.

\section*{Begin with a failure somebody can describe}

Most political repair is less dramatic. A service fails, a rule produces an avoidable burden, or a decision takes so long that it becomes a refusal in practice. People who have been waiting for years may hear another promise of institutional reform with complete justification for their boredom.

A useful proposal should begin with a failure that can be stated plainly. Applicants wait nine months for an ordinary decision. A payment intended for eligible households does not reach many of them. A public building consumes money but cannot be used. Each description points towards an investigation. None is improved by announcing a new commitment to excellence.

Take an imagined permit office. It is slow because applications repeatedly arrive incomplete, specialists review them in sequence, and nobody can tell an applicant where a case is stuck. A minister might respond by demanding that staff work faster. Another might buy a new computer system. Either measure could miss the obstruction.

A better examination follows several applications through the actual process. Which information is requested? Which of it is used? Where does a file wait, and why? Are repeated checks catching different risks or reproducing one another? Can an applicant correct a small omission without starting again?

The answers might justify a clearer form, simultaneous review of independent issues, or a named person responsible for the case. They might instead reveal a shortage of qualified staff or a statutory requirement that administrators cannot alter. Identifying the cause prevents a managerial slogan from becoming another task imposed on an already failing office.

The measure of success should follow the purpose. Faster decisions matter, but approving everything would also make an office fast. We need to know whether decisions remain lawful and accurate, whether serious errors are caught, and whether difficult applicants are being pushed aside to improve the average.

Small changes of this kind lack the grandeur of replacing a system. They can nevertheless alter daily life more than a constitutional preamble full of aspirations. An applicant who receives a fair answer in a reasonable time has obtained something real.

\section*{Experiments have conditions}

Trying a reform in a limited setting can reveal mistakes before they spread. It can also create false confidence if the trial receives exceptional staff, extra money, and attention that will disappear when the programme expands.

A pilot project is therefore not merely a small version of the future. It is a particular arrangement whose differences from ordinary operation need to be understood. Volunteers may be unusually motivated. A small group may cooperate through personal contact that will be impossible across an entire country. Other institutions may absorb costs that have not been included in the reported budget.

There should be a plan for deciding what the trial has shown. If every outcome can be described as a valuable learning experience that warrants further expansion, the experiment has no possible adverse result. A reform that cannot fail its evaluation has not submitted to evaluation.

The plan should also say what happens to people who have reorganised their lives around it. A temporary service can create reasonable reliance. Abruptly ending it may cause harm even if the larger programme is abandoned. Reversibility is a property to design, not a word that removes the consequences of reversal.

Some reforms cannot be tried safely on a small scale. Rights should not become privileges allocated experimentally without a compelling justification. Major networks can have costs and benefits that appear only when the whole arrangement operates. In such cases, comparison, modelling, and staged implementation may still help, but uncertainty should remain visible.

The modest attitude is to learn what can be learned and name what cannot. It does not require a timid scale of ambition. Large and urgent improvements may be necessary. What it rejects is the claim that the size of the ambition excuses carelessness about the means.

\section*{The repair can acquire its own defenders}

A reform is also an institution in embryo. It employs people, distributes money, and establishes a vocabulary in which success will be described. After a few years, questioning it may sound much like the obstruction its founders once complained about.

This is not proof that reform was pointless. It is a reason to apply the same questions after the founding enthusiasm has passed. Who benefits from continuation? What has changed? Which costs have become easier to ignore? Does the public still need this particular arrangement, or have its defenders begun to treat their employment as evidence of its necessity?

Review should include the possibility of ending an obsolete function. Governments are often better at adding duties than removing them. New priorities accumulate while older programmes retain supporters and legal obligations. The resulting burden is then attributed to an inexplicably sluggish administration.

Ending a programme is not costless either. Staff deserve fair treatment, useful knowledge should be preserved, and people dependent on a service need a transition. But acknowledging these obligations should help organise an ending, not prohibit one indefinitely.

A mature reformer has to allow that something she built may eventually need to close. Otherwise her demand for change contains an unannounced exception for her own work.

\section*{Disobedience and ordinary disappointment}

No account of repair can assume that lawful channels are always adequate. A law may be gravely unjust. An appeal may exist on paper while retaliation makes its use impossible. People have reasons to protest, refuse, and organise outside procedures designed to contain them.

The legitimacy of resistance cannot be decided by the mere fact that a rule has been broken. Nor can it be established by the protester's sincerity. A sincere conviction can support a just cause or a brutal one. We need to examine the injustice alleged, the means used, the people burdened, and the possibilities for redress.

Nonviolent resistance can make an ignored injury visible, but it too can impose costs on bystanders. A campaign does not acquire ownership of their time or livelihood simply because its purpose is good. The opening supermarket dispute returns here on a larger and sometimes much more serious scale: somebody who has already volunteered for a cause may begin volunteering other people.

This does not make all disruption illegitimate. Effective resistance may disrupt precisely the arrangements that sustain an injustice. It does require an argument about why this disruption is warranted and how avoidable harm will be limited. The answer should concern the actual people affected, rather than congratulate the campaign on the purity of its intentions.

Ordinary political disappointment is a different matter. Losing a vote, receiving a reasoned adverse judgement, or failing to persuade the public does not by itself show that the system has ceased to permit change. If every defeat authorises abandonment of common procedures, those procedures survive only while nobody needs them to settle a disagreement.

Living under shared rules requires a capacity to lose without declaring that loss impossible. It also requires enough honesty to recognise when the rules have been altered so that one side is expected to keep losing. Neither permanent obedience nor permanent rebellion can do this judging for us.

\section*{After the water has gone}

The flooded town will not emerge without damage. Some houses may never be occupied again. A business may close despite competent assistance. The public inquiry may establish that a reasonable decision turned out badly, or that a celebrated decision concealed waste. Recovery will not provide a single verdict that satisfies everyone.

What can improve is more definite. Residents can learn how decisions were made. Claims can be resolved. Unnecessary exceptional powers can end. The next emergency can find better equipment and clearer responsibilities. Officials can leave office without taking the records with them.

There will be no day on which the town is finally beyond mistakes. The people responsible for its repair will retire, quarrel, forget things, and be replaced by others. Some of their improvements will need improving.

That prospect need not make the work futile. A bridge can be repaired without a promise that no bridge will ever fail again. Public life deserves the same permission to accomplish something useful without pretending to finish history.

\chapter{What We Can Leave Unsettled}

The argument is over, but an improved version continues in the mind. This time the reply is immediate. The other person has no answer. Somebody who matters hears the exchange and understands exactly what happened.

Meanwhile, the actual evening passes. A meal goes cold. A person in the room says something that receives only half our attention. The victory we are arranging takes place nowhere, but it uses real time.

There are injuries that deserve years of attention. There are decisions that must be contested until they change. The difficulty is that a mind occupied by a small humiliation can feel just as busy as a mind working on either of those tasks. Intensity does not tell us whether the work remains necessary.

The supermarket quarrel at the beginning of this book did not require a philosophy. There was a question about whose turn it was, enough information to answer it, and a purchase to collect. The rest was something I added. I wanted the exchange to establish more than the order of service. I wanted the people involved to recognise what kind of people they had been.

That second ambition had no obvious finishing point. Even an apology might have been inadequate if I suspected that it was insincere. A few minutes of ordinary discourtesy could supply material for an indefinite inquiry into character.

We can make a whole political life out of this habit. The scale of the subjects changes, and some of the wrongs become very serious. The appetite for a final vindication remains recognisable.

\section*{A pause in the demand for an answer}

Pyrrhonism offers an unfamiliar response to the urge for finality. In Sextus Empiricus's account, the sceptic continues investigating rather than announcing that truth has been found or is impossible to find. Opposing appearances and arguments, when experienced as equally persuasive, lead to \emph{epochē}, suspension of judgement. \emph{Ataraxia}, tranquillity concerning matters of opinion, follows unexpectedly. Ordinary life continues through natural capacities, unavoidable feelings, customs, and acquired skills; hunger and cold do not disappear.\footnote{Sextus Empiricus, \emph{Outlines of Pyrrhonism}, book I, especially sections 1--4, 8--10, and 21--30; available in Mary Mills Patrick's \href{https://www.gutenberg.org/files/17556/17556-h/17556-h.htm}{\emph{Sextus Empiricus and Greek Scepticism}}, 1899, which includes a translation of book I. The ancient account is distinguished here from the book's own practical and political arguments.}

Sextus illustrates the unexpected result with a story about Apelles. Unable to paint the foam at a horse's mouth, the painter throws his sponge at the picture. Its mark achieves the effect he had failed to produce deliberately. The sceptic's tranquillity likewise arrives after an attempted resolution is relinquished.

There is something appealing, and something suspicious, about that story. It releases us from a demand to solve everything. It can also tempt us to turn relinquishment into a technique guaranteed to produce the very satisfaction we have supposedly stopped demanding. We throw the sponge with one eye on the result.

No such guarantee is offered here. A person may stop arguing and remain distressed. A doubt may be honestly acknowledged and continue to hurt. If tranquillity becomes another required achievement, the failure to feel it supplies a fresh accusation against an already tired person.

The attraction of the Pyrrhonian gesture is smaller and more exact. It interrupts the inference from \emph{I cannot settle this} to \emph{I must keep forcing a settlement}. Sometimes the available reasons have been examined and do not sustain a final answer. We can recognise that condition without promoting it into a theory that no answer is possible anywhere.

\section*{This book has made judgements}

A difficulty should now be obvious. These chapters have not suspended judgement on everything. They have argued that some evidence is stronger than other evidence, that some uses of power are unjust, and that some promises cannot be fulfilled as stated. They have defended practices and rejected others. A book that has spent so long doing this should not end by pretending that it has remained neutral.

Nor is modern willingness to revise a conclusion identical to the ancient suspension just described. Saying that a scientific explanation is well supported but corrigible still involves an assent. The political preferences expressed here also go beyond a report that an argument has left the writer unable to choose.

The debt to Pyrrhonism is therefore a disposition towards inquiry and the loosening of a particular demand, rather than an attempt to reconstruct an ancient way of life without alteration. I do not need to claim that the ancient sceptic secretly endorsed constitutional government, modern science, or any proposal in these pages. Those proposals have to stand on the reasons given for them.

It is possible to judge a defined question while leaving a larger one open. We can establish that a statistic was misdescribed without deciding that every person who used it was dishonest. We can reject a policy because its means are inadequate without deciding what a perfect society would be. We can refuse an unfair demand without completing an account of the moral worth of the person who made it.

The restraint lies partly in keeping the conclusion within the reach of the argument. A convincing objection to one claim does not authorise a conquest of everything nearby. The mind enjoys those conquests because they reduce future work. Once a person, institution, or political camp has been classified as wholly fraudulent, each new dispute arrives with its answer attached.

But an answer attached in advance is precisely what scepticism should disturb. If our criticism has made further examination unnecessary, it may have hardened into the kind of assurance it began by questioning.

\section*{Doubt can also become an evasion}

There is a counterfeit scepticism that asks for impossible certainty only when a conclusion is unwelcome.

An institution confronted with evidence of failure demands one more study. A critic accepts a rumour about an opponent while rejecting careful measurement by an expert as merely another opinion. Someone who benefits from an existing arrangement discovers the complexity of the world whenever a reform is proposed. Doubt is being applied selectively to preserve a desired result.

We can detect some of this by asking what would change the person's mind. There may be no decisive experiment, but the person should be able to describe how the preferred conclusion could lose support. If every contrary observation confirms a conspiracy, or every failure proves that the policy needs more time and power, inquiry has become a performance.

The same test belongs on our own desk. It is easy to appreciate uncertainty when it slows an opponent. It is less pleasant when it removes a favourite example or requires us to admit that a person we dislike has made a good point. The willingness to incur that loss is more revealing than a general declaration of open-mindedness.

There are also decisions for which waiting is a decision. A damaged structure remains in use while a committee requests further reassurance. A person remains exposed to a danger while an authority avoids committing itself. Insisting that judgement be postponed can distribute serious risks without acknowledging that they have been distributed.

A responsible choice may therefore have to be made under uncertainty. The reasons should identify what is known, what remains doubtful, and which consequences make delay or action preferable. We can act on the stronger case while arranging to detect error. What we cannot honestly do is claim that uncertainty has relieved us of participation in the outcome.

This is one limit of borrowing tranquillity for public purposes. An official's peace of mind is not a sufficient measure of his administration. The people waiting for an answer may have good reason to disturb it.

\section*{The title's accusation}

The sentimental anti-realism examined in this book begins when a desired description takes precedence over the conditions that would make it true.

A person wants to be generous, so the cost imposed on someone else disappears from the account. An institution wants an inclusive service, so the capacity needed to deliver it becomes an awkward technicality. A movement wants a painless transition, so a difficult remainder is assigned to future innovation. A critic wants a doomed civilisation, so evidence of adaptation appears shallow beside the preferred decline.

The sentiment need not be tender. Anger, nostalgia, disgust, and the pleasure of exposure can organise the same refusal. An account is sentimental in this sense when protecting the feeling becomes more important than finding out what happened or what can be done.

The subtitle gives that refusal a deliberately absurd physical form. Lifting all four limbs and expecting to float is the hope that support will continue after we have withdrawn whatever supplies it. In public argument, the missing support is often less visible than in the image. It may be somebody's unpaid work, an energy source, an institutional restraint, a willingness to bear an unpopular cost, or the confidence that a rule will still apply tomorrow.

The image should not be pushed into a general case against change. Some inherited supports can be replaced. Some were oppressive. A person can mistake a familiar burden for an indispensable foundation because he has never seen another arrangement work. The question is whether the proposed replacement has been thought through, not whether the old structure has earned immunity by being old.

Nor does the title establish a metaphysical position about the existence of an external world. Its accusation is practical: wishing does not supply the missing means, and an agreeable account does not remove a consequence. The word \emph{realism} gives nobody privileged access to every fact. It names a demand we can make of an argument, including an argument that calls itself realistic.

\section*{The critic is still in the room}

After enough exposure to convenient virtue, distrust begins to feel like intelligence. We anticipate the hidden interest, the selective sympathy, the calculation behind the declaration. We can often find one. There is no need to invent selfish motives when the world supplies them so freely.

The temptation is to regard this as the complete explanation of human conduct. Generosity becomes disguised vanity, loyalty becomes fear, conviction becomes an income stream. Nothing can count against the account, because any apparent exception proves that the disguise is effective.

This is a gloomy form of credulity. It believes too readily in the explanation it prefers. A person can want praise and still help. An official can enjoy authority and still do necessary work well. A neighbour can be difficult in one encounter and reliable in another. Mixed motives are ordinary; they do not establish that the useful action was unreal.

Relentless criticism should be relentless about its own conveniences too. The critic may obtain status from seeing through people, exemption from cooperation by declaring every effort compromised, and a reason to avoid disappointment by expecting nothing. These rewards are not less seductive because they are collected in a melancholy voice.

There is a practical loss as well. If we explain away every decent act, we become less able to recognise the arrangements that permit decency to last. We may expose hypocrisy accurately while neglecting the person who performs an unglamorous task without much hypocrisy to expose. Our account of society then becomes distorted by the very acuity we admire in it.

The remedy is not to force an optimistic balance. A corrupt institution may deserve a thoroughly adverse account. A particular person may repeatedly abuse trust. We should state that plainly when the evidence supports it. What should remain open is the possibility that another case will differ, and that our pleasure in condemnation is not itself evidence.

\section*{Some losses do not instruct us}

Melancholy need not be a forecast of universal decline. It can arise from a simpler recognition: even a life that goes reasonably well loses things it cannot replace.

A familiar voice is no longer heard. A parent who once answered every practical question now needs help with one. A house is sold, a friendship thins, a landscape changes. The new arrangement may have advantages. Those advantages do not retrospectively make the lost thing unimportant.

Public language is often impatient with this residue. It asks what can be learned, how resilience will be built, which opportunity has appeared. Such questions can be useful when a decision has to be made. They become intrusive when they require every loss to justify itself by producing an improvement.

There may be nothing to learn that the bereaved person did not already know. The relationship mattered, and it has ended. A service may ease the practical burden. Company may help. Neither has to pretend to replace the person in order to be worthwhile.

The same distinction matters at a social scale. A town may need to change an industry that can no longer support it. Some residents may find better work elsewhere. Others may lose a web of habits and recognition that a higher income cannot reproduce. Taking that loss seriously does not prove that the industry can or should be preserved indefinitely.

Acknowledgement is not a veto. Equally, a decision to proceed is not a finding that nothing valuable is being lost. We can make the decision and allow grief to remain outside its justification. The people affected need not endorse the entire account merely because no better feasible option was available.

This leaves less satisfaction than a story in which every necessary change is eventually vindicated. It may also leave less bitterness. People are sometimes driven to defend an impossible restoration because the only alternative offered is to agree that their former life was worthless. A more truthful account does not demand that agreement.

\section*{What can be enough for today}

Ataraxia is easy to turn into another product: a calmer mind, better habits, a method for protecting oneself from the news. There may be useful practices among such offerings. But tranquillity cannot do all the work of deciding what deserves attention.

A person can be undisturbed because he is well protected from consequences that others endure. Another can be disturbed because she is caring for someone who needs her. Calm does not establish wisdom, and distress does not establish a philosophical mistake.

What may be loosened is the additional demand that every encounter confirm an intelligible and acceptable world. The neighbour can remain irritating without becoming a specimen of civilisational collapse. A political defeat can remain a defeat without proving that everyone is corrupt. An unanswered question can remain on the table without invalidating the meal beside it.

This is not an instruction to reduce every concern to private comfort. Some things on the table need doing: a letter, a visit, a decision that affects another person. The question is whether our thought is helping us do them, or continually replacing them with a larger trial at which the whole world must appear.

There is no universal rule for recognising the difference. Sometimes a recurring worry points to an obligation we have avoided. Sometimes it is the exhausted repetition of an argument whose practical work is complete. We may need another person's judgement to tell them apart. The wish to become entirely self-sufficient in this matter would merely begin another impossible project.

Ordinary activities can continue without resolving the background. We can make something, learn a skill, visit a friend, or prepare food while remaining uncertain about progress, justice, or the final value of a life. Such activities are not made trivial by their failure to answer the larger questions. They are among the things about which those questions were supposed to matter.

\section*{Leaving the argument}

This book has followed arguments from small acts of generosity to institutions with the power to reorganise millions of lives. The scales differ. So do the standards of evidence and the seriousness of the harms. A queue is not a state, a mouse enclosure is not a civilisation, and a household budget is not a model of the climate. Keeping those differences visible has been part of the work.

What recurs is a temptation to ask a pleasing explanation to do more than it can. It may conceal a burden, expand a finding into a prophecy, or turn a preference into an obligation for someone else. Exposing the excess can restore a possibility of judgement. It cannot guarantee that the remaining choice will be pleasant.

We will still sometimes choose badly. We may be unable to repair the consequence. We may understand another person's reasons and continue to oppose what he did. Some disagreements will outlast the people who began them.

There need be no final reconciliation in which every thread is gathered and every loss assigned its proper place. We can leave something unsettled without ceasing to care about it. We can also cease caring about a small quarrel without proving that the quarrel never mattered at all.

At the counter, I had a reasonable objection. I still think so. What I no longer need is a further verdict on the two strangers, their lives, or mine.

I could have collected the fish and gone home. There was still dinner to make.
\end{document}
